
% Revised source: notation audit and structural clarification, 10 September 2026.
\documentclass[oneside]{article}

% Compile with latexmk -pdf, or rerun pdfLaTeX until references stabilize.
\usepackage[utf8]{inputenc}
\usepackage{lmodern}
\usepackage{amsmath,amsthm,amsfonts,amssymb,mathtools}
\usepackage[T1]{fontenc}
\usepackage{graphicx,array,booktabs,longtable}
\usepackage{microtype}
\usepackage[a4paper,total={6in,8in}]{geometry}
\usepackage{enumitem}
\usepackage{tikz-cd}
\usetikzlibrary{decorations.pathmorphing,arrows.meta}
\usepackage{xcolor}
\usepackage[colorlinks=true,linkcolor=black,citecolor=black,urlcolor=blue]{hyperref}

% Notation used throughout the manuscript.
\newcommand{\Boxes}{\mathsf{Boxes}}
\newcommand{\Conf}{\mathsf{Conf}}
\newcommand{\Sym}{\mathsf{Sym}}
\newcommand{\WH}{\mathsf{WH}}
\newcommand{\QG}{\mathsf{QG}}
\newcommand{\Wit}{\mathsf{Wit}}
\newcommand{\Disc}{\mathsf{Disc}}
\newcommand{\Sp}{\mathbf{Sp}}
\newcommand{\Span}{\mathbf{Span}}
\newcommand{\Kar}{\operatorname{Kar}}
\newcommand{\cc}{\mathsf{cc}}
\newcommand{\id}{\operatorname{id}}
\newcommand{\op}{\mathrm{op}}
\newcommand{\trcorner}{\mathbin{\urcorner}}
\newcommand{\leftedge}[2]{\mathsf{L}(#1,#2)}
\newcommand{\bottomedge}[2]{\mathsf{B}(#1,#2)}
\newcommand{\xnrightarrow}[1]{\mathrel{\overset{#1}{\nrightarrow}}}
\newcommand{\act}{\mathsf{act}}
\newcommand{\extprod}{\mathbin{\otimes_{\mathrm{ext}}}}
\newcommand{\boxnw}[1]{\operatorname{nw}(#1)}
\newcommand{\boxne}[1]{\operatorname{ne}(#1)}
\newcommand{\boxsw}[1]{\operatorname{sw}(#1)}
\newcommand{\boxse}[1]{\operatorname{se}(#1)}
\setlength{\emergencystretch}{1.5em}

\theoremstyle{plain}
\newtheorem{theorem}{Theorem}[section]
\newtheorem{corollary}[theorem]{Corollary}
\newtheorem{lemma}[theorem]{Lemma}
\newtheorem{proposition}[theorem]{Proposition}
\theoremstyle{definition}
\newtheorem{definition}[theorem]{Definition}
\newtheorem{example}[theorem]{Example}
\newtheorem{judgement}[theorem]{Judgement}
\newtheorem{construction}[theorem]{Construction}
\newtheorem{problem}[theorem]{Problem}
\theoremstyle{remark}
\newtheorem{remark}[theorem]{Remark}
\newtheorem{notation}[theorem]{Notation}
\newtheorem{convention}[theorem]{Convention}
\newtheorem{observation}[theorem]{Observation}

\title{Quantum Groupoids from Thin Concurrent Games}
\author{Conor Sheridan}
\date{\today}

\begin{document}
\maketitle

\begin{abstract}
Thin concurrent games provide a canonical source of quantum groupoids.
Unique negative--positive factorization of configuration symmetries gives
a vacant double groupoid whose boxes carry an involutive weak Hopf monoid
in spans. We identify its separable Frobenius base with the configuration
set and construct the associated Day--Street quantum groupoid, including
its boundary coactions and inversion data. Tensor products are preserved,
and game negation corresponds to compact transposition. Every distributor
between configuration-symmetry groupoids yields an explicit quantum
module; positive-witness distributors therefore realize thin strategies
as quantum modules. We characterize the positive symmetry transport
retained by these modules and show why full distributor data are needed
for full-coend composition. Copycat agrees with the regular quantum
identity exactly when negative symmetry is discrete. Retaining the
generating distributors gives a symmetric monoidal double category and
a canonical symmetric monoidal oplax realization of thin concurrent games.
\end{abstract}

\tableofcontents

\section{Introduction}
\label{sec:cgqg-introduction}

Day and Street~\cite{DayStreet2004} observed that the role of
star-autonomous categories in linear logic suggests further interaction
between computer science and quantum group theory. Their categorical
formulation of quantum groupoids provides a setting for this connection.
We construct a quantum groupoid from the polarized configuration
symmetries of a thin concurrent game and extend the construction to
strategies through their positive-witness distributors.

Thin concurrent games describe computation using causal dependencies,
concurrent events, and selected symmetries between finite configurations
\cite{CastellanClairambaultWinskel2019}. Their negative and positive
symmetry subgroupoids determine a unique factorization of every selected
symmetry~\cite[Lemma~5.3]{ClairambaultOlimpieriPaquet2025}. This
factorization is the input to our construction. It gives a unique
commuting square with negative horizontal edges and positive vertical
edges for each compatible corner. The resulting vacant double groupoid
is denoted by \(T(A)\).

The algebraic mechanism has two established antecedents.
Andruskiewitsch and Natale~\cite{AndruskiewitschNatale2005} construct weak
Hopf algebras from vacant double groupoids, using vertical composition
and horizontal factorization of boxes. Pastro and Street
\cite{PastroStreet2009} construct quantum groupoids from weak Hopf
monoids with invertible antipode. We identify the required box structure
in thin games and realize its operations in the category of spans.
Spans retain the individual factorization witnesses, so the construction
also applies when a game has infinitely many configurations or selected
symmetries. Its structural equations are proved by bijections of witness
sets.

For every thin concurrent game \(A\), we construct an involutive weak
Hopf monoid \(\WH_{\mathrm{sp}}(A)\) and a Day--Street quantum groupoid
\(\QG_{\mathrm{sp}}(A)\). The configuration set \(\Conf(A)\), with
its canonical separable commutative Frobenius structure, is the
object-of-objects. We exhibit both counital splittings, identify the
cotensor of arrows as a boundary pullback, and give the boundary
coactions, composition, unit, and inversion maps. The inversion map on
the composition object is also described directly by factoring and
pasting boxes. Tensor products act componentwise on this data, and game
negation corresponds to compact transposition. A finite example with
nontrivial symmetry of both polarities makes the incidence operations
explicit.

The strategy construction starts with an ordinary distributor
\(F:\Sym(A)^{\op}\times\Sym(B)\to\mathbf{Set}\). Its endpoint-indexed
carrier has a diagonal comonad, and positive endpoint transport gives a
quantum-module action. Applying this construction to the positive-witness
distributor \(\lVert\sigma\rVert\) assigns a quantum module to each
thin strategy \(\sigma\). We prove the comonad and action equations,
functoriality on natural transformations, and compatibility with external
products.

Two comparison results describe the information retained by this
realization. First, the quantum module depends exactly, at the level of
its carrier and fixed-endpoint module maps, on the restriction of the
distributor to positive symmetry
(Proposition~\ref{prop:cgqg-positive-restriction}). Distributors with the
same such restriction can have different full-coend composites
(Proposition~\ref{prop:cgqg-full-transport-example}). Second, the quantum
module of copycat is isomorphic to the regular quantum module exactly
when the game's negative symmetry groupoid is discrete
(Proposition~\ref{prop:cgqg-copycat-unit-iff}). These results explain the
role of the generating distributor in composition and identity.

Accordingly, our horizontal arrows are pairs \((F,Q_F^{\mathrm{sp}})\),
with composition induced by the full distributor coend. Functors between
the symmetry groupoids supply representable adjunctions, and distributor
squares supply the corresponding quantum-module maps. Adjoining these
prescribed quantum components gives a symmetric monoidal double category
canonically equivalent to the game-labelled distributor double category.
The positive-witness collapse of
\cite[Theorem~5.15]{ClairambaultOlimpieriPaquet2025} then yields a
symmetric monoidal oplax double functor
\[
 \mathfrak Q:\mathbf{TCG}\longrightarrow
 \mathbf{QMod}^{\mathrm{rep}}_{\mathrm{sp}}.
\]
Its identity and tensor comparisons are invertible. A two-move
interaction exhibits a proper composition comparison and identifies the
causal compatibility that the full distributor coend does not impose.

Section~\ref{sec:cgqg-strict-background} recalls the game-semantic input, and
Section~\ref{sec:quantum-categories-and-quantum-modules} fixes the quantum
and span conventions. Section~\ref{sec:cgqg-game-qg} constructs the game
quantum groupoid. Section~\ref{sec:cgqg-strategy-quantum-modules} treats
distributor-induced modules, and
Section~\ref{sec:cgqg-composition-scope} assembles the double-categorical
realization and its comparison results.

\subsection*{Note on use of AI}
Generative AI systems were used during the research and development of this
manuscript as tools for mathematical investigation, discussion, and
editorial refinement.  Responsibility for the mathematical content,
arguments, and conclusions rests with the author.

\section{Thin concurrent games and polarized symmetry}
\label{sec:cgqg-strict-background}

Concurrent games combine causal structure with symmetries between finite executions. In a thin concurrent game these symmetries are polarized: negative symmetries encode Opponent reindexing and positive symmetries encode Player reindexing. Their interaction is especially rigid, since every selected symmetry admits a unique negative--positive decomposition. Full definitions, including the strategy axioms, are given in~\cite{CastellanClairambaultWinskel2019,ClairambaultOlimpieriPaquet2025}; we use the polarized decomposition in the orientation of~\cite[Lemma~5.3]{ClairambaultOlimpieriPaquet2025}. We use the general consistency-family presentation of event structures from~\cite{CastellanClairambaultWinskel2019}. The constructions used below are polarized decomposition, symmetry lifting, and the positive-witness collapse of~\cite{ClairambaultOlimpieriPaquet2025}.

\subsection{Event structures and configurations}

\begin{definition}[Event structure and configuration
{\cite[Def.~2.1]{CastellanClairambaultWinskel2019}}]
\label{def:cgqg-event-structure-background}
    An \textbf{event structure} is a triple
    \[
        E=(|E|,\le_E,\mathsf{Con}_E)
    \]
    where \(|E|\) is a set of events, \(\le_E\) is a partial order with finite principal ideals
    \[
        [e]_E:=\{e'\mid e'\le_E e\},
    \]
    and \(\mathsf{Con}_E\) is a nonempty family of finite subsets of \(|E|\) satisfying:
    \begin{enumerate}
        \item \textbf{singleton consistency}: for every event \(e\in |E|\),
        \[
            \{e\}\in\mathsf{Con}_E;
        \]

        \item \textbf{heredity}: if \(X\in\mathsf{Con}_E\) and \(Y\subseteq X\), then
        \[
            Y\in\mathsf{Con}_E;
        \]

        \item \textbf{causal closure}: if \(X\in\mathsf{Con}_E\), \(e\in X\), and \(e'\le_E e\), then
        \[
            X\cup\{e'\}\in\mathsf{Con}_E.
        \]
    \end{enumerate}
    A finite subset \(X\subseteq |E|\) is \textbf{consistent} when \(X\in\mathsf{Con}_E\). A finite \textbf{configuration} is a finite subset \(x\subseteq |E|\) which is consistent and down-closed:
    \[
        e\in x,\ e'\le_E e
        \quad\Longrightarrow\quad
        e'\in x.
    \]
    We write \(\Conf(E)\) for the set of finite configurations.
\end{definition}

The causal order records which events must already have occurred before a given event can occur, while consistency records which finite collections of events may coexist in one execution. A configuration is therefore a finite causal state of execution. Concurrent events are represented by incomparable, jointly consistent events, and their joint configuration records their coexistence directly.

For example, if \(E\) consists of two concurrent events \(e_1,e_2\), then
\[
    \Conf(E)=
    \{\varnothing,\{e_1\},\{e_2\},\{e_1,e_2\}\}.
\]

The configuration order is the diamond
\begin{equation*}
    \begin{tikzcd}
        & \varnothing \arrow[dr,hook] \arrow[dl,hook'] & \\
        \{e_1\} \arrow[dr,hook] && \{e_2\} \arrow[dl,hook'] \\
        &  \{e_1,e_2\}  &
    \end{tikzcd}
\end{equation*}
whose two singleton configurations are incomparable one-event extensions of the empty configuration. 

A \textbf{polarized event structure} carries a map
\[
    \operatorname{pol}_E:|E|\longrightarrow\{-,+\},
\]
where \(-\) denotes an Opponent/Environment move and \(+\) denotes a Player/System move~\cite[Def.~2.3]{CastellanClairambaultWinskel2019}.

\subsection{Configuration symmetries}

\begin{definition}[Isomorphism family and event structure with symmetry
{\cite[Def.~3.1]{CastellanClairambaultWinskel2019}}]
\label{def:cgqg-ess-background}
    An \textbf{isomorphism family} on an event structure \(E\) is a groupoid \(\Sym(E)\) whose object set is \(\Conf(E)\) and whose arrows are selected bijections
    \[
        \theta:x\xrightarrow{\sim}y,
        \qquad x,y\in\Conf(E),
    \]
    satisfying:
    \begin{enumerate}
        \item \textbf{restriction}: if \(\theta:x\xrightarrow{\sim}y\) is selected and
        \(x'\in\Conf(E)\) with \(x'\subseteq x\), then the restricted bijection
        \[
            \theta|_{x'}:x'\xrightarrow{\sim}\theta(x')
        \]
        is selected;
 
        \item \textbf{extension}: if \(\theta:x\xrightarrow{\sim}y\) is selected and \(x\subseteq x'\in\Conf(E)\), then \(\theta\) extends to some selected \(\theta':x'\xrightarrow{\sim}y'\).
    \end{enumerate}
    An \textbf{event structure with symmetry} is an event structure equipped with an isomorphism family.
\end{definition}

A \textbf{map of event structures} \(f:E\to F\) is a function on events that sends configurations to configurations and is injective on each configuration. It is a \textbf{map of event structures with symmetry} when the pointwise image \(f\theta:f(x)\xrightarrow{\sim}f(y)\) of every selected symmetry \(\theta:x\xrightarrow{\sim}y\) is selected~\cite[Defs.~2.6 and 3.6]{CastellanClairambaultWinskel2019}.

\subsection{Thin games and polarized decomposition}

A thin concurrent game refines the configuration-symmetry groupoid by resolving selected reindexings according to polarity.

\begin{definition}[Thin concurrent game {\cite[Def.~3.16]{CastellanClairambaultWinskel2019}; \cite[Def.~5.2]{ClairambaultOlimpieriPaquet2025}}]
\label{def:cgqg-thin-game-background}
    A \textbf{thin concurrent game} \(A\) is a polarized event structure equipped with three isomorphism families
    \[
        \Sym(A),\qquad \Sym^-(A),\qquad \Sym^+(A),
    \]
    such that
    \[
        \Sym^-(A),\Sym^+(A)\subseteq\Sym(A),
    \]
    every selected symmetry preserves polarity, and the following polarized thinness conditions hold:
    \begin{enumerate}[label=\textup{(\roman*)}]
        \item\label{ax:cgqg-orthogonality} \textbf{orthogonality}: if \(\theta\in\Sym^-(A)\cap\Sym^+(A)\), then \(\theta=\id_x\) for some \(x\in\Conf(A)\);

        \item \textbf{\(-\)-receptivity}: if \(\theta\in\Sym^-(A)\) and \(\theta\subseteq^-\theta'\in\Sym(A)\), then \(\theta'\in\Sym^-(A)\);

        \item \textbf{\(+\)-receptivity}: if \(\theta\in\Sym^+(A)\) and \(\theta\subseteq^+\theta'\in\Sym(A)\), then \(\theta'\in\Sym^+(A)\).
    \end{enumerate}
    Here \(\theta\subseteq^p\theta'\) means that \(\theta'\) extends \(\theta\) and all newly added pairs of events have polarity \(p\).
\end{definition}

We use \emph{thin game} as shorthand for \emph{thin concurrent game} throughout.

Since the three symmetry families are groupoids on the same configuration set and
\[
    \Sym^-(A),\Sym^+(A)\subseteq\Sym(A),
\]
the groupoids \(\Sym^-(A)\) and \(\Sym^+(A)\) are wide subgroupoids of \(\Sym(A)\), with common object set \(\Conf(A)\). When the underlying arrow sets are needed explicitly, we write \(\Sym_1(A)\), \(\Sym^-_1(A)\), and \(\Sym^+_1(A)\). Writing \(\Disc(X)\) for the discrete groupoid on a set \(X\), Condition~\ref{ax:cgqg-orthogonality} gives
\begin{equation}
    \Sym^-(A)\cap \Sym^+(A)=\Disc(\Conf(A)).
\label{eq:cgqg-polarized-trivial-intersection}
\end{equation}

The polarization records which side of the interaction carries a reindexing: positive symmetries are carried by Player moves, while negative symmetries are carried by Opponent moves~\cite{CastellanClairambaultWinskel2019}.

The polarized subgroupoids jointly determine a unique normal form for every selected symmetry.

\begin{theorem}[Polarized decomposition
{\cite[Lemma~5.3]{ClairambaultOlimpieriPaquet2025}}]
\label{thm:cgqg-polarized-decomposition}
    For every thin concurrent game \(A\) and every selected symmetry \(\theta:x\xrightarrow{\sim}z\) in \(\Sym(A)\), there are unique \(y\in \Conf(A)\), \(h:x\xrightarrow{\sim}y\) in \(\Sym^-(A)\), and \(v:y\xrightarrow{\sim}z\) in \(\Sym^+(A)\) such that
    \begin{equation}
        \theta=v\circ h.
    \label{eq:cgqg-polarized-decomposition}
    \end{equation}

    Equivalently, every selected symmetry fits into a unique factorization triangle
    \begin{equation*}
        \begin{tikzcd}[column sep=huge,row sep=large]
            x \arrow[rr,"\theta"] \arrow[dr,"h"'] && z \\
            & y \arrow[ur,"v"'] &
        \end{tikzcd}
    \end{equation*}
    with the first leg negative and the second positive.
\end{theorem}
\begin{proof}
    \leavevmode
    \begin{enumerate}[beginpenalty=10000]
        \item \textbf{Apply polarized decomposition.} The existence and uniqueness are \cite[Lemma~5.3]{ClairambaultOlimpieriPaquet2025}, with \(h=\theta^-\) and \(v=\theta^+\). The factorization equation gives the displayed triangle. \qedhere
    \end{enumerate}
\end{proof}

Polarized decomposition expresses each selected symmetry as an Opponent reindexing followed by a Player reindexing, with the intermediate configuration uniquely determined. 

\subsection{Tensor product and duality}

The dual \(A^\perp\) has the same underlying event structure as \(A\) with reversed event polarity, while \(A\otimes B\) is the tagged disjoint union of \(A\) and \(B\), with causality and consistency inherited componentwise~\cite[Def.~2.8 and the following paragraph]{CastellanClairambaultWinskel2019}. The configuration and polarized symmetry data are compatible with these operations as follows.

\begin{proposition}[Tensor product and duality {\cite[Defs.~2.8, 3.5 and \S\S3.2.2--3.3.3]{CastellanClairambaultWinskel2019}}]
\label{prop:cgqg-background-tensor-duality}
    For thin concurrent games \(A,B\), there are canonical identifications
    \begin{equation}
        \Conf(A\otimes B)
        \cong
        \Conf(A)\times\Conf(B),
    \label{eq:cgqg-conf-tensor}
    \end{equation}
    and
    \begin{equation}
 \begin{aligned}
 \Sym(A^\perp)&=\Sym(A),\\
 \Sym^-(A^\perp)&=\Sym^+(A),\\
 \Sym^+(A^\perp)&=\Sym^-(A),
 \end{aligned}
 \label{eq:cgqg-dual-symmetry}
\end{equation}
\begin{equation}
 \begin{aligned}
 \Sym(A\otimes B)&\cong\Sym(A)\times\Sym(B),\\
 \Sym^-(A\otimes B)&\cong\Sym^-(A)\times\Sym^-(B),\\
 \Sym^+(A\otimes B)&\cong\Sym^+(A)\times\Sym^+(B).
 \end{aligned}
 \label{eq:cgqg-tensor-symmetry}
\end{equation}
\end{proposition}
\begin{proof}
    \leavevmode
    \begin{enumerate}[beginpenalty=10000]
        \item \textbf{Decompose configurations and symmetries.} Configurations and selected symmetries of \(A\otimes B\) decompose componentwise, giving the Cartesian-product identifications above \cite[Defs.~2.8 and 3.5]{CastellanClairambaultWinskel2019}. Tensor product therefore acts componentwise on the polarized symmetry families.

        \item \textbf{Reverse polarity.} Duality reverses event polarity and exchanges the positive and negative symmetry families~\cite[\S\S3.2.2--3.3.3]{CastellanClairambaultWinskel2019}. This gives~\eqref{eq:cgqg-dual-symmetry}. \qedhere
    \end{enumerate}
\end{proof}

\subsection{Strategies and positive-witness distributors}
\label{subsec:cgqg-background-strategies}

A strategy adds a causal realization of behaviour to the game structure introduced above. Its events record concrete occurrences of moves, together with the causal dependencies between them, while a display map records which move of the game each occurrence realizes. Different configurations of the strategy may therefore realize the same configuration of the game.

For a configuration \(x\in\Conf(E)\) and an event \(e\notin x\), write
\[
    x\vdash_E e
\]
when \(x\cup\{e\}\in\Conf(E)\). We also write \(e_1\prec_E e_2\) when \(e_1<_{E}e_2\) is an immediate causal dependency, meaning that no event lies strictly between them~\cite[\S4.1.1]{ClairambaultOlimpieriPaquet2025}.

\begin{definition}[Thin strategy {\cite[Defs.~4.2 and 5.4]{ClairambaultOlimpieriPaquet2025}}]
\label{def:cgqg-background-thin-strategy}
    Let \(A,B\) be thin concurrent games. A \textbf{thin strategy}
    \[
        \sigma:A\longrightarrow B
    \]
    consists of an event structure with symmetry \(\sigma\) together with a display map
    \[
        \partial_\sigma:
        \sigma\longrightarrow A^\perp\otimes B
    \]
    which is a map of event structures with symmetry. The polarity of an event \(s\in|\sigma|\) is inherited from its image:
    \[
        \operatorname{pol}_\sigma(s)
        :=
        \operatorname{pol}_{A^\perp\otimes B}(\partial_\sigma(s)).
    \]
    For \(\theta:x\xrightarrow{\sim}y\) in \(\Sym(\sigma)\), write
    \[
        \partial_\sigma\theta:
        \partial_\sigma x\xrightarrow{\sim}\partial_\sigma y
    \]
    for its pointwise image under the display map. The following conditions are required.
    \begin{enumerate}
        \item \textbf{receptivity}: if \(x\in\Conf(\sigma)\) and a negative event \(a\) satisfies
        \[
            \partial_\sigma x\vdash_{A^\perp\otimes B} a,
        \]
        then there is a unique event \(s\) such that
        \[
            x\vdash_\sigma s,
            \qquad
            \partial_\sigma(s)=a;
        \]

        \item \textbf{courtesy}: if \(s_1\prec_\sigma s_2\) and either \(\partial_\sigma(s_1)\) is positive or \(\partial_\sigma(s_2)\) is negative, then
        \[
            \partial_\sigma(s_1)
            \prec_{A^\perp\otimes B}
            \partial_\sigma(s_2);
        \]

        \item \textbf{symmetry-receptivity}: let \(\theta:x\xrightarrow{\sim}y\) be in \(\Sym(\sigma)\). If \(a,b\) are negative events,
        \[
            (a,b)\notin\partial_\sigma\theta,
        \]
        and
        \[
            \partial_\sigma\theta\cup\{(a,b)\}
        \]
        is a selected symmetry of \(A^\perp\otimes B\), then there are unique events \(s,t\) with
        \[
            \partial_\sigma(s)=a,
            \qquad
            \partial_\sigma(t)=b,
        \]
        such that
        \[
            \theta\cup\{(s,t)\}:
            x\cup\{s\}\xrightarrow{\sim}y\cup\{t\}
        \]
        belongs to \(\Sym(\sigma)\);

        \item \textbf{thinness}: if \(\theta:x\xrightarrow{\sim}y\) belongs to \(\Sym(\sigma)\) and
        \[
            \partial_\sigma\theta
            \in\Sym^+(A^\perp\otimes B),
        \]
        then
        \[
            x=y,
            \qquad
            \theta=\id_x.
        \]
    \end{enumerate}
\end{definition}

Receptivity and courtesy control the causal behaviour of the strategy relative to the game. Symmetry-receptivity performs the corresponding lifting for Opponent-carried reindexings, while thinness makes the remaining Player-carried symmetry rigid inside the strategy. Together these conditions give a canonical lifting property for arbitrary selected symmetries.

\begin{lemma}[Symmetry lifting up to positive symmetry
{\cite[Lemma~5.5]{ClairambaultOlimpieriPaquet2025}}]
\label{lem:cgqg-background-strategy-symmetry-lifting}
    Let \(\sigma:A\to B\) be a thin strategy, \(x^\sigma\in\Conf(\sigma)\), and
    \[
        \psi:
        \partial_\sigma x^\sigma
        \xrightarrow{\sim}
        y
    \]
    a selected symmetry of \(A^\perp\otimes B\). There are unique
    \[
        \varphi:
        x^\sigma\xrightarrow{\sim}z^\sigma
        \quad\text{in }\Sym(\sigma)
    \]
    and
    \[
        \theta^+:
        \partial_\sigma z^\sigma
        \xrightarrow{\sim} y
        \quad\text{in }\Sym^+(A^\perp\otimes B)
    \]
    such that
    \begin{equation}
        \theta^+\circ\partial_\sigma\varphi=\psi.
    \label{eq:cgqg-background-strategy-symmetry-lifting}
    \end{equation}
\end{lemma}
\begin{proof}
    \leavevmode
    \begin{enumerate}[beginpenalty=10000]
        \item \textbf{Lift the selected endpoint symmetry.} Apply \cite[Lemma~5.5]{ClairambaultOlimpieriPaquet2025} to the display map \(\partial_\sigma:\sigma\to A^\perp\otimes B\). It supplies a strategy symmetry \(\varphi:x^\sigma\to z^\sigma\) and a positive residual symmetry \(\theta^+:\partial_\sigma z^\sigma\to y\) satisfying \(\theta^+\circ\partial_\sigma\varphi=\psi\).

        \item \textbf{Verify uniqueness.} For another such pair \(\varphi',\theta^{+\prime}\),
        \[
            \partial_\sigma(\varphi'\circ\varphi^{-1})
              =(\theta^{+\prime})^{-1}\circ\theta^+
              \in\Sym^+(A^\perp\otimes B).
        \]
        Thinness makes \(\varphi'\circ\varphi^{-1}\) an identity, hence
        \(\varphi'=\varphi\) and \(\theta^{+\prime}=\theta^+\). \qedhere
    \end{enumerate}
\end{proof}

For \(x^\sigma\in\Conf(\sigma)\), the tensor-product identification of configurations gives unique
\[
    x_A^\sigma\in\Conf(A),
    \qquad
    x_B^\sigma\in\Conf(B)
\]
such that
\[
    \partial_\sigma x^\sigma
    =
    x_A^\sigma\otimes x_B^\sigma.
\]
If \(\varphi:x^\sigma\xrightarrow{\sim}y^\sigma\) lies in \(\Sym(\sigma)\), write
\[
    \partial_\sigma\varphi
    =
    \varphi_A^\sigma\otimes\varphi_B^\sigma
\]
for its components under the tensor-product identification, so that
\(\varphi_A^\sigma:x_A^\sigma\xrightarrow{\sim}y_A^\sigma\) and
\(\varphi_B^\sigma:x_B^\sigma\xrightarrow{\sim}y_B^\sigma\).

\begin{definition}[\(+\)-covered configuration
{\cite[\S4.1.4]{ClairambaultOlimpieriPaquet2025}}]
\label{def:cgqg-background-plus-covered}
    A configuration \(x^\sigma\in\Conf(\sigma)\) is \textbf{\(+\)-covered} when every maximal event of \(x^\sigma\) is positive. Write
    \[
        \Conf^+(\sigma)\subseteq\Conf(\sigma)
    \]
    for the set of \(+\)-covered configurations.
\end{definition}

Thus the display restricts to an endpoint map
\begin{equation*}
    \begin{tikzcd}[column sep=huge]
        \Conf^+(\sigma)
          \arrow[r,"{x^\sigma\mapsto(x_A^\sigma,x_B^\sigma)}"]
        & \Conf(A)\times\Conf(B).
    \end{tikzcd}
\end{equation*}
By the tensor and duality identifications above,
\[
    \Sym^+(A^\perp\otimes B)
    \cong
    \Sym^-(A)\times\Sym^+(B).
\]
Accordingly, the positive residual symmetry in Lemma~\ref{lem:cgqg-background-strategy-symmetry-lifting} consists of a negative reindexing at the source and a positive reindexing at the target. This gives the endpoint-normalized configurations used below.

\begin{definition}[Positive witness {\cite[\S5.2]{ClairambaultOlimpieriPaquet2025}}]
\label{def:cgqg-background-positive-witness}
    Let \(\sigma:A\to B\) be a thin strategy. A \textbf{positive witness} from \(a\in\Conf(A)\) to \(b\in\Conf(B)\) is a triple
    \begin{equation}
        (\theta_A^-,x^\sigma,\theta_B^+),
        \qquad
        \theta_A^-:
        a\longrightarrow x_A^\sigma
        \text{ in }\Sym^-(A),
        \qquad
        \theta_B^+:
        x_B^\sigma\longrightarrow b
        \text{ in }\Sym^+(B),
    \label{eq:cgqg-background-positive-witness}
    \end{equation}
    with \(x^\sigma\in\Conf^+(\sigma)\). Write
    \[
        \Wit_\sigma^+(a,b)
    \]
    for the set of positive witnesses from \(a\) to \(b\).
\end{definition}

The polarized frames make these witnesses stable under arbitrary selected symmetries of their endpoints. The lifting lemma supplies the corresponding renormalization uniquely.

\begin{proposition}[Endpoint transport {\cite[Proposition~5.6]{ClairambaultOlimpieriPaquet2025}}]
\label{prop:cgqg-background-positive-witness-transport}
    Let
    \[
        w=(\theta_A^-,x^\sigma,\theta_B^+)
        \in\Wit_\sigma^+(a,b),
    \]
    and let
    \[
        \Omega_A:a'\longrightarrow a
        \quad\text{in }\Sym(A),
        \qquad
        \Omega_B:b\longrightarrow b'
        \quad\text{in }\Sym(B).
    \]
    There are unique
    \[
        \varphi:x^\sigma\xrightarrow{\sim}y^\sigma
        \quad\text{in }\Sym(\sigma),
    \]
    \[
        \vartheta_A^-:
        a'\longrightarrow y_A^\sigma
        \quad\text{in }\Sym^-(A),
        \qquad
        \vartheta_B^+:
        y_B^\sigma\longrightarrow b'
        \quad\text{in }\Sym^+(B),
    \]
    for which the diagrams
    \begin{equation*}
        \begin{tikzcd}[column sep=huge, row sep=large]
            a
              \arrow[r,"\theta_A^-"]
            &
            x_A^\sigma
              \arrow[d,"\varphi_A^\sigma"]
            \\
            a'
              \arrow[u,"\Omega_A"]
              \arrow[r,"\vartheta_A^- "']
            &
            y_A^\sigma
        \end{tikzcd}
        \qquad
        \begin{tikzcd}[column sep=huge, row sep=large]
            x_B^\sigma
              \arrow[r,"\theta_B^+"]
              \arrow[d,"\varphi_B^\sigma"']
            &
            b
              \arrow[d,"\Omega_B"]
            \\
            y_B^\sigma
              \arrow[r,"\vartheta_B^+"']
            &
            b'
        \end{tikzcd}
    \end{equation*}
    commute. Hence endpoint transport is defined by
    \begin{equation}
        \Omega_{B*}\Omega_A^*w
        :=
        (\vartheta_A^-,y^\sigma,\vartheta_B^+)
        \in\Wit_\sigma^+(a',b').
    \label{eq:cgqg-background-positive-witness-transport}
    \end{equation}
\end{proposition}
\begin{proof}
    \leavevmode
    \begin{enumerate}[beginpenalty=10000]
        \item \textbf{Lift the endpoint transport.} Apply Lemma~\ref{lem:cgqg-background-strategy-symmetry-lifting} to the selected endpoint symmetry determined by \(\Omega_A,\Omega_B,\theta_A^-\), and \(\theta_B^+\).

        \item \textbf{Identify the polarized residuals.} The tensor and duality formulas identify the positive residual symmetry with the displayed pair \(\vartheta_A^-,\vartheta_B^+\), giving the two commuting diagrams. Uniqueness follows from the lifting lemma. This is the endpoint transport of \cite[Proposition~5.6]{ClairambaultOlimpieriPaquet2025}. \qedhere
    \end{enumerate}
\end{proof}

\begin{proposition}[Positive-witness distributor {\cite[Proposition~5.7]{ClairambaultOlimpieriPaquet2025}}]
\label{prop:cgqg-background-positive-witness-distributor}
    For every thin strategy \(\sigma:A\to B\), endpoint transport makes the positive-witness sets into a distributor
    \begin{equation}
        \lVert\sigma\rVert:
        \Sym(A)^\op\times\Sym(B)
        \longrightarrow\mathbf{Set},
        \qquad
        (a,b)\longmapsto\Wit_\sigma^+(a,b).
    \label{eq:cgqg-positive-witness-distributor}
    \end{equation}
\end{proposition}
\begin{proof}
    \leavevmode
    \begin{enumerate}[beginpenalty=10000]
        \item \textbf{Verify identity transport.} Transport along identity symmetries returns the original witness by uniqueness in Proposition~\ref{prop:cgqg-background-positive-witness-transport}.

        \item \textbf{Verify composition.} Successive endpoint transports and transport along the corresponding composite symmetries satisfy the same commuting diagrams. Uniqueness therefore makes them equal. These are the identity and composition laws for
        \[
            \Sym(A)^\op\times\Sym(B)\longrightarrow\mathbf{Set}.\qedhere
        \]
    \end{enumerate}
\end{proof}

\subsection{The double category of thin concurrent games}
\label{subsec:cgqg-background-double-category}

Throughout, double categories may be weak in horizontal composition:
their horizontal arrows form bicategories.

Thin games, strategies, and their morphisms assemble into a two-dimensional structure. Besides strategies relating games horizontally, maps of games act vertically, while positive morphisms compare strategies compatibly with those maps.

A \textbf{map of thin concurrent games}
\[
    h:A\longrightarrow C
\]
is a map of event structures with symmetry that is rigid, meaning that it preserves causal dependency, and that preserves event polarities and the polarized symmetry families, in the sense that~\cite[\S5.3.1]{ClairambaultOlimpieriPaquet2025}
\[
    h(\Sym^-(A))\subseteq\Sym^-(C),
    \qquad
    h(\Sym^+(A))\subseteq\Sym^+(C).
\]
For such a map, \(h^\perp:A^\perp\to C^\perp\) denotes the same underlying map on events, regarded between the dual games.

Given thin strategies
\[
    \sigma:A\longrightarrow B,
    \qquad
    \tau:C\longrightarrow D
\]
and maps of thin games
\[
    h:A\longrightarrow C,
    \qquad
    k:B\longrightarrow D,
\]
a \textbf{positive morphism of strategies}
\[
    \begin{tikzcd}[column sep=huge,row sep=large]
        A \arrow[r,Rightarrow,"\sigma"] \arrow[d,"h"']
        & B \arrow[d,"k"] \\
        C \arrow[r,Rightarrow,"\tau"']
        & D
        \arrow[ul,Rightarrow,from=1-2,to=2-1,phantom,"\scriptstyle f"]
    \end{tikzcd}
\]
is a rigid map of event structures with symmetry
\[
    f:\sigma\longrightarrow\tau
\]
for which the display square commutes up to positive symmetry:
\[
    \partial_\tau\circ f
    \simeq^+
    (h^\perp\otimes k)\circ\partial_\sigma.
\]
Here \(\simeq^+\) means that for every \(x\in\Conf(\sigma)\), the canonical pointwise bijection
\[
    \bigl\{
        ((h^\perp\otimes k)(\partial_\sigma(s)),\partial_\tau(f(s)))
        \mid s\in x
    \bigr\}
\]
belongs to \(\Sym^+(C^\perp\otimes D)\)~\cite[Def.~5.8]{ClairambaultOlimpieriPaquet2025}.

\begin{proposition}[The double category of thin concurrent games {\cite[Prop.~5.11]{ClairambaultOlimpieriPaquet2025}}]
\label{prop:cgqg-background-tcg-double-category}
    There is a symmetric monoidal double category \(\mathbf{TCG}\) whose
    \begin{enumerate}[label=\textup{(\roman*)}]
        \item objects are thin concurrent games;
    
        \item vertical arrows are maps of thin concurrent games;
    
        \item horizontal arrows are thin strategies;
    
        \item squares are positive morphisms of strategies.
    \end{enumerate}
    Horizontal composition is interaction followed by hiding, horizontal identities are copycat strategies, and the symmetric monoidal product is induced by tensor product.
\end{proposition}
\begin{proof}
    \leavevmode
    \begin{enumerate}[beginpenalty=10000]
        \item \textbf{Describe interaction and hiding.}
        For composable strategies
        \[
            \sigma:A\longrightarrow B,
            \qquad
            \tau:B\longrightarrow C,
        \]
        interaction synchronizes configurations \(x^\sigma\in\Conf^+(\sigma)\) and \(x^\tau\in\Conf^+(\tau)\) whose displayed \(B\)-configurations agree and whose synchronization is causally compatible, meaning that the relation on the synchronization graph induced by the two causal orders is acyclic~\cite[\S4.2.2]{ClairambaultOlimpieriPaquet2025}. Hiding the synchronized \(B\)-events then produces a configuration of the composite
        \[
            \tau\odot\sigma:A\longrightarrow C.
        \]
        More precisely, there is an order-isomorphism
        \[
            \left\{
                (x^\sigma,x^\tau)
                \in\Conf^+(\sigma)\times\Conf^+(\tau)
                \;\middle|\;
                \begin{array}{c}
                x_B^\sigma=x_B^\tau\\
                \text{and the pair is causally compatible}
                \end{array}
            \right\}
            \xrightarrow{\cong}
            \Conf^+(\tau\odot\sigma),
        \]
        and the resulting composite configuration has exterior boundary
        \[
            (x_A^\sigma,x_C^\tau).
        \]
        This is the interaction-and-hiding description of strategy composition~\cite[Prop.~4.7]{ClairambaultOlimpieriPaquet2025}.

        \item \textbf{Identify the horizontal identity.}
        The horizontal identity on \(A\) is the copycat strategy
        \[
            \cc_A:A\longrightarrow A.
        \]
        Its two copies of \(A\) are linked so that corresponding occurrences are causally related from the negative occurrence to the positive one. Its \(+\)-covered configurations are canonically indexed by configurations of \(A\):
        \[
            \Conf(A)\xrightarrow{\cong}\Conf^+(\cc_A),
            \qquad
            x\longmapsto\cc_A x,
        \]
        with
        \[
            \partial_{\cc_A}(\cc_A x)=x\otimes x.
        \]
        See \cite[Prop.~4.6]{ClairambaultOlimpieriPaquet2025}.

        \item \textbf{Assemble the double-category structure.}
        For thin strategies, synchronization also determines the selected symmetries of the composite, while synchronization up to symmetry gives horizontal composition of positive morphisms~\cite[Props.~5.9--5.10]{ClairambaultOlimpieriPaquet2025}. Together with componentwise tensor product, these constructions satisfy the symmetric monoidal double-category axioms~\cite[Prop.~5.11]{ClairambaultOlimpieriPaquet2025}. \qedhere
    \end{enumerate}
\end{proof}

Here \(\mathbf{Dist}\) denotes the double category of groupoids, functors, distributors, and natural transformations~\cite[Example~3.3]{ClairambaultOlimpieriPaquet2025}.

\begin{theorem}[Positive-witness collapse
{\cite[Thm.~5.15]{ClairambaultOlimpieriPaquet2025}}]
\label{thm:cgqg-background-positive-witness-double-functor}
    The assignments
    \[
        A\longmapsto\Sym(A),
        \qquad
        \sigma\longmapsto\lVert\sigma\rVert
    \]
    extend to a symmetric monoidal oplax double functor
    \[
        \lVert-\rVert:
        \mathbf{TCG}\longrightarrow\mathbf{Dist}.
    \]
    On vertical arrows it sends
    \[
        h:A\longrightarrow B
    \]
    to the induced functor
    \[
        \Sym(h):\Sym(A)\longrightarrow\Sym(B),
    \]
    and on positive morphisms of strategies it acts by the corresponding natural transformations of positive-witness distributors.

    The horizontal identity comparison is the natural isomorphism
    \[
        \lVert\cc_A\rVert
        \xrightarrow{\cong}
        \Sym(A)(-,-),
    \]
    while for composable thin strategies
    \[
        \sigma:A\longrightarrow B,
        \qquad
        \tau:B\longrightarrow C,
    \]
    the compositor is the natural transformation
    \begin{equation}
        \gamma_{\sigma,\tau}:
        \lVert\tau\odot\sigma\rVert
        \longrightarrow
        \lVert\tau\rVert\circ\lVert\sigma\rVert.
    \label{eq:cgqg-background-positive-witness-compositor}
    \end{equation}
\end{theorem}
\begin{proof}
    \leavevmode
    \begin{enumerate}[beginpenalty=10000]
        \item \textbf{Construct the assignments.} The action on objects and horizontal arrows is the positive-witness distributor construction above. Maps of thin games induce functors between their configuration-symmetry groupoids, and positive morphisms induce natural transformations by endpoint transport~\cite[\S\S5.4.1--5.4.2]{ClairambaultOlimpieriPaquet2025}.

        \item \textbf{Use the collapse comparisons.} The copycat unitor and strategy-composition compositor are constructed in~\cite[Props.~5.12--5.13]{ClairambaultOlimpieriPaquet2025}. Their coherence and compatibility with tensor give the symmetric monoidal oplax double functor~\cite[Thm.~5.15]{ClairambaultOlimpieriPaquet2025}. \qedhere
    \end{enumerate}
\end{proof}

The one-dimensional organization of the same game-semantic operations is the compact closed category \(\mathbf{Tcg}\) of~\cite[Thm.~3.37]{CastellanClairambaultWinskel2019}, whose tensor product is \(\otimes\) and whose duality is \(A\mapsto A^\perp\). The double-category presentation additionally retains maps of games, positive morphisms of strategies, and the oplax compositor above.

\section{Quantum categories and quantum modules}
\label{sec:quantum-categories-and-quantum-modules}

Quantum groupoids have their roots in several problems in mathematical physics and
operator algebra in which natural symmetry structures turned out to be more general
than ordinary Hopf algebras. In low-dimensional quantum theory and conformal field
theory, Mack and Schomerus were led to weak quasi-Hopf symmetries in order to
accommodate truncated tensor products of physical sectors
\cite{MackSchomerus1992}. In solvable lattice models, Hayashi introduced face
algebras to describe quantum-group symmetry of interaction-round-a-face models
whose partition functions depend on boundary conditions, with applications to the
Jones index problem \cite{Hayashi1993,Hayashi1998}. Closely related weak Hopf and
finite quantum-groupoid structures arose in the study of finite-index and
finite-depth subfactors, where they describe finite-index inclusions,
intermediate subalgebras, bimodule categories, and principal-graph data
\cite{NikshychVainerman2000,NikshychVainerman2002}. These
developments were organized algebraically in the theory of weak Hopf algebras of
Böhm--Nill--Szlachányi \cite{BohmNillSzlachanyi1999}: much of the familiar Hopf
calculus survives, but the unit and counit laws are weakened so that nontrivial
source and target bases can remain. Day--Street subsequently formulated quantum
categories and quantum groupoids categorically, with a quantum groupoid intended
as a ``Hopf algebra with several objects'' \cite{DayStreet2004}. Pastro--Street
related this formulation directly to weak Hopf monoids in braided monoidal
categories \cite{PastroStreet2009}, and Chikhladze developed the corresponding
theory of quantum modules \cite{Chikhladze2010QuantumModules}. We develop the span and comodule formulations used in the construction.

\subsection{Spans as coefficient-free incidence morphisms}

\begin{definition}[Span category]
\label{def:cgqg-span-category}
    Let \(\mathcal E\) be a locally small category with finite limits, in a fixed universe in which the isomorphism classes of spans between fixed endpoints form sets. Define
    \[
        \Sp(\mathcal E):=\Span_{\cong}(\mathcal E)
    \]
    to be the category whose objects are those of \(\mathcal E\) and whose morphisms \(X\to Y\) are isomorphism classes of spans
    \[
        X\xleftarrow{\ell}S\xrightarrow{r}Y.
    \]
    Composition is pullback. Cartesian product makes \(\Sp(\mathcal E)\) symmetric monoidal. For \(\mathcal E=\mathbf{Set}\), write simply \(\Sp\).

    Thus if
    \begin{equation*}
       \begin{tikzcd}
            & S \ar[dr] \ar[dl]& \\
            X & & Y
        \end{tikzcd}
        \quad 
        \text{and}
        \quad
        \begin{tikzcd}
            & T \ar[dr] \ar[dl] & \\
            Y & & Z
        \end{tikzcd} 
    \end{equation*}    
    are composed, the witness object is displayed by
    \begin{equation*}
        \begin{tikzcd}
            && S\times_YT
                \ar[dr]
                \ar[dl]
                \ar[dd, phantom, "\rotatebox{-45}{$\lrcorner$}" description, very near start]
            && \\
            & S \ar[dr] \ar[dl] && T \ar[dr] \ar[dl] & \\
            X && Y && Z.
        \end{tikzcd}
    \end{equation*}
\end{definition}

\subsection{Weak Hopf monoids and incidence}

\begin{definition}[Weak bimonoid {\cite[Def.~1.1]{PastroStreet2009}}]
\label{def:cgqg-background-weak-bimonoid}
    Let
    \[
        \mathcal V=(\mathcal V,\otimes,I,c)
    \]
    be a braided monoidal category. A \textbf{weak bimonoid} \(H\) consists of
    the following data and compatibility conditions:
    \begin{enumerate}
        \item \textbf{monoid structure}: morphisms
        \[
            m:H\otimes H\longrightarrow H,
            \qquad
            \eta:I\longrightarrow H,
        \]
        making \((H,m,\eta)\) a monoid, so that \(m\) is associative and \(\eta\) is its two-sided unit;

        \item \textbf{comonoid structure}: morphisms
        \[
            \Delta:H\longrightarrow H\otimes H,
            \qquad
            \varepsilon:H\longrightarrow I,
        \]
        making \((H,\Delta,\varepsilon)\) a comonoid, so that \(\Delta\) is coassociative and \(\varepsilon\) is its counit;

        \item \textbf{weak bimonoid compatibility}: multiplication and comultiplication satisfy the braided multiplicativity law, namely the diagram
        \begin{equation}
            \begin{tikzcd}[column sep=small, row sep=huge]
                & H\otimes H\otimes H\otimes H \ar[rr, "1\otimes c_{H,H}\otimes1"] & & H\otimes H\otimes H\otimes H \ar[dr, "m \otimes m"] & \\
                H\otimes H \ar[ur, "\Delta \otimes \Delta"] \ar[drr, swap, "m"] &&&& H\otimes H \\
                && H \ar[urr, swap, "\Delta"] &&
            \end{tikzcd}
            \label{eq:cgqg-background-weak-bimonoid-multiplicativity}
        \end{equation}
        commutes, while the bimonoid compatibility of the unit and counit with the opposite structure is replaced by the weak unit and weak counit equations of
        \cite[Def.~1.1]{PastroStreet2009}.
    \end{enumerate}
    The weak unit and weak counit equations determine canonical source and target counital idempotents
    \[
        \varepsilon_s,\varepsilon_t:H\longrightarrow H.
    \]
\end{definition}

In the symmetric case used below, the counital maps are the composites
\begin{align}
 \varepsilon_t&=(\varepsilon m\otimes1)(1\otimes c)(\Delta\eta\otimes1),
 \label{eq:cgqg-canonical-target-composite}\\
 \varepsilon_s&=(1\otimes\varepsilon m)(c\otimes1)(1\otimes\Delta\eta).
 \label{eq:cgqg-canonical-source-composite}
\end{align}
These formulas fix the counital convention independently of the later
antipode calculation.

The adjective ``weak'' refers specifically to the interaction of the unit and counit with the opposite structure: the weakened unit and counit laws allow the resulting object to retain nontrivial source and target bases.

\begin{definition}[Weak Hopf monoid {\cite[Def.~2.1]{PastroStreet2009}}]
\label{def:cgqg-background-weak-hopf-monoid}
    A \textbf{weak Hopf monoid} is a weak bimonoid \(H\) equipped with a morphism
    \[
        S:H\longrightarrow H,
    \]
    called the \textbf{antipode}, satisfying the following commutative diagrams:
    \begin{equation}
        \begin{tikzcd}[column sep=huge,row sep=large]
            & H
                \arrow[d,"\Delta"]
                \arrow[dl,"\varepsilon_t"']
                \arrow[dr,"\varepsilon_s"]
            & \\
            H
            &
            H\otimes H
                \arrow[l,"m(1\otimes S)"]
                \arrow[r,"m(S\otimes1)"']
            &
            H
        \end{tikzcd}
        \qquad
        \begin{tikzcd}[column sep=huge,row sep=large]
            H
                \arrow[r,"\Delta^{(3)}"]
                \arrow[d,"S"']
            &
            H\otimes H\otimes H
                \arrow[d,"m^{(3)}(S\otimes1\otimes S)"]
            \\
            H
                \arrow[r,"1_H"']
            &
            H .
        \end{tikzcd}
        \label{eq:cgqg-background-antipode}
    \end{equation}
\end{definition}

For an ordinary Hopf monoid the counital idempotents collapse to
\(\eta\varepsilon\). For a weak Hopf monoid, splitting the counital idempotents gives the
source and target base objects, as in the Karoubi-completion construction
below. These bases provide the object structure of the corresponding
quantum category.

\subsection{Separable Frobenius bases and boundary comodules}

\begin{definition}[Frobenius, commutative, and separable Frobenius monoids]
\label{def:cgqg-background-frobenius-monoid}
    Following \cite[\S A.4 and Def.~A.5]{PastroStreet2009} for the Frobenius
    and separability conventions, let \((\mathcal V,\otimes,I)\) be a monoidal category. A \textbf{Frobenius monoid} \(C\) consists of a monoid
    \[
        (C,\mu,\eta)
    \]
    and a comonoid
    \[
        (C,\delta,\epsilon)
    \]
    on the same object, satisfying the \textbf{Frobenius equations}
    \begin{equation}
        (\mu\otimes1_C)(1_C\otimes\delta)
        =
        \delta\mu
        =
        (1_C\otimes\mu)(\delta\otimes1_C).
        \label{eq:cgqg-background-frobenius}
    \end{equation}
    Equivalently, suppressing coherence isomorphisms
    \begin{equation*}
        \begin{tikzcd}[column sep=huge,row sep=large]
            && C\otimes C
            \arrow[dl,"\delta\otimes1_C"']
            \arrow[d,"\mu"]
            \arrow[dr,"1_C\otimes\delta"]
            && \\
            & C\otimes C\otimes C
            \arrow[dr,"1_C\otimes\mu"']
            & C
            \arrow[d,"\delta"]
            & C\otimes C\otimes C
            \arrow[dl,"\mu\otimes1_C"]
            & \\
            && C\otimes C,
            &&
        \end{tikzcd}
    \end{equation*}
    commutes.

    If \(\mathcal V\) is braided, a Frobenius monoid is \textbf{commutative} when its multiplication and comultiplication satisfy
    \begin{equation*}
        \begin{tikzcd}
            C\otimes C
                \arrow[rr,"c_{C,C}"]
                \arrow[dr,"\mu"']
            &&
            C\otimes C
                \arrow[dl,"\mu"]
            \\
            &
            C
        \end{tikzcd}
        \qquad
        \begin{tikzcd}
            &
            C
                \arrow[dl,"\delta"']
                \arrow[dr,"\delta"]
            &
            \\
            C\otimes C
                \arrow[rr,"c_{C,C}"']
            &&
            C\otimes C .
        \end{tikzcd}
    \end{equation*}
    Thus \(\mu c_{C,C}=\mu\) and \(c_{C,C}\delta=\delta\).

    A Frobenius monoid is \textbf{separable} when
    \begin{equation}
        \begin{tikzcd}[column sep=huge, row sep=large]
            & C \otimes C \ar[dr, "\mu"] & \\
            C \ar[rr, "1_C"'] \ar[ur, "\delta"] && C
        \end{tikzcd}
    \label{eq:cgqg-background-frobenius-separable}
    \end{equation}
    commutes, equivalently
    \[
        \mu\delta=1_C.
    \]
\end{definition}

\begin{proposition}[Canonical Frobenius structure in spans
{\cite[Examples~3.7 and 4.8]{ContrerasKellerMehta2022}}]
\label{prop:cgqg-background-frobenius-base}
    For every set \(P\), the spans
    \[
    \begin{aligned}
    \mu_P:\quad
    &\begin{tikzcd}
        & P \arrow[dl, "\Delta_P"'] \arrow[dr, "1_P"] & \\
        P\times P & & P
    \end{tikzcd}
    \qquad
    \eta_P:\quad
    \begin{tikzcd}
        & P \arrow[dl] \arrow[dr, "1_P"] & \\
        1 & & P
    \end{tikzcd}
    \\[1.5em]
    \delta_P:\quad
    &\begin{tikzcd}
        & P \arrow[dl, "1_P"'] \arrow[dr, "\Delta_P"] & \\
        P & & P\times P
    \end{tikzcd}
    \qquad
    \epsilon_P:\quad
    \begin{tikzcd}
        & P \arrow[dl, "1_P"'] \arrow[dr] & \\
        P & & 1
    \end{tikzcd}
    \end{aligned}
    \]
    define a canonical separable commutative Frobenius monoid
    \[
        C_P=(P,\mu_P,\eta_P,\delta_P,\epsilon_P)
    \]
    in \(\Sp\).
\end{proposition}
\begin{proof}
    \leavevmode
    \begin{enumerate}[beginpenalty=10000]
        \item \textbf{Identify the Frobenius spans.} Specializing the constructions of~\cite[Examples~3.7 and 4.8]{ContrerasKellerMehta2022} to the discrete groupoid on \(P\) gives the displayed Frobenius spans.

        \item \textbf{Verify commutativity and separability.} Commutativity follows from the diagonal. The pullback computing \(\mu_P\delta_P\) is canonically \(P\), with both exterior legs the identity, so
        \[
            \mu_P\delta_P=1_P.\qedhere
        \]
    \end{enumerate}
\end{proof}

\begin{remark}[Pastro--Street convention
{\cite[\S2.1, Def.~2.1 and Prop.~2.3(15)]{PastroStreet2009}}]
\label{rem:cgqg-ps-convention}
    We use categorical composition order. Relative to the notation of~\cite{PastroStreet2009}, our counital maps are
    \begin{equation}
        r^{\mathrm{PS}}=\varepsilon_t,
        \qquad
        t^{\mathrm{PS}}=\varepsilon_s.
        \label{eq:cgqg-ps-an-convention-bridge}
    \end{equation}
    Indeed, \cite[Def~2.1]{PastroStreet2009} gives
    \[
        m(S\otimes1)\Delta=t^{\mathrm{PS}},
        \qquad
        m(1\otimes S)\Delta=r^{\mathrm{PS}}.
    \]
    Moreover \(r^{\mathrm{PS}}=S s^{\mathrm{PS}}\) by \cite[Prop.~2.3(15)]{PastroStreet2009}; hence, when \(S\) is invertible,
    \[
        s^{\mathrm{PS}}=S^{-1}\varepsilon_t.
    \]
\end{remark}

\subsection{Quantum categories and quantum groupoids}

The bicategorical description of an ordinary category provides a useful model for the quantum notion. A small category with object set \(A_0\) and arrow set \(A_1\) is equivalently a monad on \(A_0\) in the bicategory \(\Span(\mathbf{Set})\)~\cite[Example~2.3]{McCrudden2002}. Here \(\Span(\mathbf{Set})\) denotes the bicategory of spans and maps of spans; \(\Sp\) denotes the category of isomorphism classes of spans introduced above. The underlying endomorphism is the span
\begin{equation*}
    \begin{tikzcd}
        & A_1
            \arrow[dl,"s"']
            \arrow[dr,"t"]
        & \\
        A_0 && A_0,
    \end{tikzcd}
\end{equation*}
while the multiplication and unit of the monad are maps of spans
\begin{equation*}
    \begin{tikzcd}[column sep=huge,row sep=large]
        &
        A_1\times_{A_0} A_1
            \arrow[dl,"s\pi_1"']
            \arrow[d,"m"]
            \arrow[dr,"t\pi_2"]
        &
        \\
        A_0
        &
        A_1
            \arrow[l,"s"]
            \arrow[r,"t"']
        &
        A_0
    \end{tikzcd}
    \qquad
    \begin{tikzcd}[column sep=huge,row sep=large]
        &
        A_0
            \arrow[dl,"1_{A_0}"']
            \arrow[d,"u"]
            \arrow[dr,"1_{A_0}"]
        &
        \\
        A_0
        &
        A_1
            \arrow[l,"s"]
            \arrow[r,"t"']
        &
        A_0 .
    \end{tikzcd}
\end{equation*}
Thus span composition forms the pullback of composable arrows, and the monad
associativity and unit laws are precisely the ordinary category axioms.

In Chikhladze's general comodule setting, where the required coreflexive
equalizers are available and preserved by tensoring, the corresponding span-like
bicategory is
\[
    \mathbf{Comod}(\mathcal V)
\]
of comonoids, bicomodules, and comodule maps
\cite[\S3]{Chikhladze2011QuantumCategories}. A \(1\)-cell
\[
    M:C\rightsquigarrow{} D
\]
is a \(C\)--\(D\)-bicomodule. For composable bicomodules
\[
    M:C\rightsquigarrow{} D,
    \qquad
    N:D\rightsquigarrow{} E,
\]
their composite is the cotensor product
\[
    M\boxtimes_DN,
\]
defined by the equalizer
\begin{equation}
    \begin{tikzcd}[column sep=huge]
        M\boxtimes_DN
            \arrow[r,hook,"j"]
        &
        M\otimes N
            \arrow[r,shift left=1.2ex,"\rho_M\otimes1_N"]
            \arrow[r,shift right=1.2ex,"1_M\otimes\lambda_N"']
        &
        M\otimes D\otimes N .
    \end{tikzcd}
    \label{eq:cgqg-background-comodule-composition}
\end{equation}
Thus cotensor product plays in \(\mathbf{Comod}(\mathcal V)\) the role played by pullback composition in \(\Span(\mathbf{Set})\). For configuration bases, Theorem~\ref{thm:cgqg-split-cotensor} constructs the cotensors explicitly as split absolute equalizers.

For \(\mathcal V=\mathbf{Set}\), each set has its diagonal comonoid, and a \(C\)--\(D\)-bicomodule is equivalently a span
\[
    C\xleftarrow{}M\xrightarrow{}D.
\]
Under this identification the equalizer above is the ordinary pullback, so that
\[
    \mathbf{Comod}(\mathbf{Set})
    \simeq
    \Span(\mathbf{Set}).
\]
Correspondingly, quantum categories in \(\mathbf{Set}\) are precisely ordinary small categories~\cite[Def.~5.2, p.~13]{Chikhladze2011QuantumCategories}.

For a comonoid
\[
    A_0=(A_0,\delta_0,\epsilon_0)
\]
in the braided monoidal category \(\mathcal V\), define its \textbf{opposite comonoid}
\[
    A_0^\circ
    :=
    \bigl(A_0,c_{A_0,A_0}\delta_0,\epsilon_0\bigr).
\]
Thus \(A_0^\circ\) has the same underlying object and counit as \(A_0\), but its comultiplication is reversed by the braiding. In \(\mathbf{Comod}(\mathcal V)\), the comonoid \(A_0^\circ\) is a right bidual of \(A_0\) \cite[\S3]{Chikhladze2011QuantumCategories}. The corresponding evaluation and coevaluation are \(1\)-cells
\[
    e:
    A_0\otimes A_0^\circ
    \rightsquigarrow{} I,
    \qquad
    n:
    I
    \rightsquigarrow{} A_0^\circ\otimes A_0,
\]
satisfying the biduality triangle identities.

The resulting \textbf{enveloping pseudomonoid}
\[
    A_0^e:=A_0^\circ\otimes A_0
\]
has multiplication and unit
\begin{equation*}
    \begin{tikzcd}[column sep=huge]
        A_0^e\otimes A_0^e
            \arrow[r,
                "p_A^e=1_{A_0^\circ}\otimes e\otimes1_{A_0}"]
        &
        A_0^e
        &
        I
            \arrow[l,"j_A^e=n"']
        .
    \end{tikzcd}
\end{equation*}

\begin{definition}[Quantum category {\cite[\S12]{DayStreet2004}; \cite[\S5.2]{PastroStreet2009}; \cite[Def.~5.2]{Chikhladze2011QuantumCategories}}]
\label{def:cgqg-background-quantum-category}
    A \textbf{quantum category} \((A_0,A_1)\) in \(\mathcal V\) consists of a comonoid \(A_0\), called the \textbf{object-of-objects}, together with a monoidal comonad
    \[
        A_1:A_0^e\rightsquigarrow{} A_0^e
    \]
    on the enveloping pseudomonoid
    \[
        A_0^e=A_0^\circ\otimes A_0
    \]
    in \(\mathbf{Comod}(\mathcal V)\). We call \(A_1\) the \textbf{arrow object}.
\end{definition}

Unpacked, the underlying comonad determines an arrow comonoid \(A_1\) with source and target maps
\begin{equation*}
    \begin{tikzcd}
        & A_1
            \arrow[dl,"s"']
            \arrow[dr,"t"]
        & \\
        A_0^\circ && A_0,
    \end{tikzcd}
\end{equation*}
while the monoidal structure on the comonad supplies a composition morphism and unit
\begin{equation*}
    \begin{tikzcd}[column sep=huge]
        A_1\boxtimes_{A_0}A_1
            \arrow[r,"\mu"]
        &
        A_1
        &
        A_0
            \arrow[l,"u"']
        .
    \end{tikzcd}
\end{equation*}
The quantum-category axioms are the corresponding monoidal-comonad compatibility conditions. In this form the analogy with an ordinary category is direct: the object set \(A_0\) is replaced by the object comonoid \(A_0\), the arrow set \(A_1\) by the arrow comonoid \(A_1\), spans by bicomodules, and pullback of composable arrows by cotensor product.

The weak-bimonoid formulation introduced above is related to this bicategorical description by the following correspondence.

\begin{theorem}[Weak-bimonoid/quantum-category correspondence
{\cite[Thm.~6.1]{PastroStreet2009}}]
\label{thm:cgqg-background-pastro-street}
In the setting of \cite[\S5]{PastroStreet2009}, where tensoring preserves
its coreflexive equalizers, weak bimonoids in the Cauchy (Karoubi)
completion \(\Kar(\mathcal V)\) correspond to quantum categories whose
object-of-objects is a separable Frobenius monoid. An invertible antipode
induces the corresponding quantum-groupoid structure.
\end{theorem}

The passage to \(\Kar(\mathcal V)\) allows the source and target counital idempotents of a weak bimonoid to split. Recall that an object of \(\Kar(\mathcal V)\) is a pair \((X,e)\) with
\[
    e:X\longrightarrow X,
    \qquad
    e^2=e,
\]
and that a splitting
\begin{equation*}
    \begin{tikzcd}[column sep=huge]
        Y
            \arrow[r,"i"]
        &
        X
            \arrow[r,"r"]
        &
        Y
    \end{tikzcd}
    \qquad
    ri=1_Y,
    \qquad
    ir=e
\end{equation*}
identifies the idempotent with its image \(Y\). In the forward Pastro--Street construction, splitting the target counital idempotent \(t^{\mathrm{PS}}\) produces the separable Frobenius object-of-objects \(A_0\), while the corresponding idempotent on \(A_1\otimes A_1\) selects the object of composable arrows.

For the inversion data we also need the composition comodules. Write
\(\Delta_1,\varepsilon_1\) for the arrow comonoid structure and put
\[
 \gamma_l=c^{-1}_{A_0,A_1}(1\otimes s)\Delta_1,
 \qquad \gamma_r=(1\otimes t)\Delta_1.
\]
The opposite-comonoid notation is understood in these underlying
morphisms of \(\mathcal V\). Let
\[
 P=A_1\boxtimes_{A_0}A_1,
 \qquad \iota_P:P\longrightarrow A_1\otimes A_1
\]
be the cotensor and its equalizer inclusion. We use ``cotensor'' for this
equalizer; it is the tensor product over the base in the terminology of
\cite{PastroStreet2009}.

The quantum-category structure induces commuting coactions
\[
 \lambda_P:P\longrightarrow A_1\otimes A_1\otimes P,
 \qquad \rho_P:P\longrightarrow P\otimes A_1,
\]
characterized by
\begin{align}
 (1\otimes1\otimes\iota_P)\lambda_P
   &=(1\otimes c_{A_1,A_1}\otimes1)
       (\Delta_1\otimes\Delta_1)\iota_P,
 \label{eq:cgqg-composition-left-arrow-coaction}\\
 (\iota_P\otimes1)\rho_P
   &=(1\otimes1\otimes\mu)\lambda_P.
 \label{eq:cgqg-composition-right-arrow-coaction}
\end{align}
The equalizer universal property gives the first map; the
quantum-category compatibility equations give the factorization in the
second. Their combined coaction is
\begin{equation}
 \chi_P=(\lambda_P\otimes1)\rho_P
       =(1\otimes1\otimes\rho_P)\lambda_P:
 P\longrightarrow A_1\otimes A_1\otimes P\otimes A_1.
 \label{eq:cgqg-composition-combined-coaction}
\end{equation}
These are the composition comodules of
\cite[\S\S5.1--5.3]{PastroStreet2009}.

\begin{definition}[Quantum groupoid {\cite[\S5.3]{PastroStreet2009}}]
\label{def:cgqg-background-quantum-groupoid}
    A \textbf{quantum groupoid} is a quantum category \(A=(A_0,A_1)\)
    equipped with comonoid isomorphisms
    \[
        \upsilon:A_0^{\circ\circ}\xrightarrow{\cong}A_0,
        \qquad \nu:A_1^\circ\xrightarrow{\cong}A_1,
    \]
    and an isomorphism \(\theta:P_l\xrightarrow{\cong}P_r\) of left
    \(A_1^{\otimes3}\)-comodules. Here \(P_l\) and \(P_r\) have the
    common carrier \(P\) and the respective coactions
    \begin{align}
      \lambda_{P_l}
       &=(1\otimes1\otimes c_{P,A_1})
         (1\otimes1\otimes1_P\otimes\nu)\chi_P,
       \label{eq:cgqg-groupoid-left-composition-coaction}\\
      \lambda_{P_r}
       &=c^{-1}_{A_1\otimes A_1\otimes P,A_1}
         (1\otimes1\otimes1_P\otimes\nu^{-1})\chi_P.
       \label{eq:cgqg-groupoid-right-composition-coaction}
    \end{align}
    These data satisfy:
    \begin{enumerate}[label=\textup{(G\arabic*)}]
        \item\label{ax:cgqg-G1} \(s\nu=t\);
        \item\label{ax:cgqg-G2} \(t\nu=\upsilon s\);
        \item\label{ax:cgqg-G3}
        \[
            (1_{A_0}\otimes1_{A_0}\otimes\upsilon)
            c_{A_0,A_0\otimes A_0}\,\varsigma
            =\varsigma\theta,
        \]
        where
        \[
            \varsigma
            =(s\otimes1_{A_0}\otimes t)(\gamma_r\otimes1_{A_1})\iota_P
            =(s\otimes1_{A_0}\otimes t)(1_{A_1}\otimes\gamma_l)\iota_P.
        \]
    \end{enumerate}
\end{definition}

\subsection{Quantum modules and multitensors}

Quantum modules play for quantum categories the role that distributors play for
ordinary categories \cite{Chikhladze2010QuantumModules}. In the module calculus
distinguish the formal Chikhladze orientation
\[
 \widehat E_{A,B}:=B_0^\circ\otimes A_0
\]
from its right-mate presentation
\[
 E_{A,B}:=A_0^\circ\otimes B_0,\qquad E_A:=E_{A,A}.
\]
The formal module comonad is on \(\widehat E_{A,B}\); all explicit boundary,
action, and multitensor calculations below use the equivalent right-mate
presentation on \(E_{A,B}\), made explicit for configuration bases in
Proposition~\ref{prop:cgqg-configuration-qmodule-right-mate}.
The canonical action of the enveloping pseudomonoids is
\begin{equation}
 \act_{A,B}:=
 1_{A_0^\circ}\otimes e_{A_0}\otimes e_{B_0}\otimes1_{B_0}:
 E_A\otimes E_{A,B}\otimes E_B\rightsquigarrow E_{A,B}.
 \label{eq:cgqg-background-canonical-module-action}
\end{equation}
Its associativity and unit constraints are induced by the biduality identities.

\begin{definition}[Quantum module {\cite[Defs.~3 and 7]{Chikhladze2010QuantumModules}}]
\label{def:cgqg-background-quantum-module}
Let \(A=(A_0,A_1)\) and \(B=(B_0,B_1)\) be quantum categories. Formally, a
\textbf{quantum module} \(M:A\xnrightarrow{}B\) has its module comonad on
\(\widehat E_{A,B}\). In the right-mate presentation used below it is represented
by a comonad \(M\) on \(E_{A,B}\) and a \(2\)-cell
\begin{equation}
 \kappa_M:
 \act_{A,B}(A_1\otimes M\otimes B_1)
 \Longrightarrow M \act_{A,B}.
 \label{eq:cgqg-background-quantum-module-direct-action}
\end{equation}
Put \(D=A_1\otimes M\otimes B_1\). The pair
\((\act_{A,B},\kappa_M)\) is required to be a map of comonads, meaning
\begin{align}
 (\epsilon_M \act_{A,B})\kappa_M
   &=\act_{A,B}\epsilon_D,
 \label{eq:cgqg-background-module-counit}\\
 (\delta_M \act_{A,B})\kappa_M
   &=(M\kappa_M)(\kappa_M D)(\act_{A,B}\delta_D).
 \label{eq:cgqg-background-module-comultiplication}
\end{align}
Write \(p_A^{e},j_A^{e}\) and \(p_B^{e},j_B^{e}\) for the multiplications and units
of the two enveloping pseudomonoids, and write
\[
\begin{aligned}
\phi_A &:p_A^{e}(A_1\otimes A_1)\Rightarrow A_1p_A^{e},
&\phi_A^0 &:j_A^{e}\Rightarrow A_1j_A^{e},\\
\phi_B &:p_B^{e}(B_1\otimes B_1)\Rightarrow B_1p_B^{e},
&\phi_B^0 &:j_B^{e}\Rightarrow B_1j_B^{e}
\end{aligned}
\]
for the monoidal structure cells. Put \(\act=\act_{A,B}\) and
\[
\pi=p_A^{e}\otimes1_{E_{A,B}}\otimes p_B^{e},\qquad
\lambda=1_{E_A}\otimes \act\otimes1_{E_B},\qquad
\jmath=j_A^{e}\otimes1_{E_{A,B}}\otimes j_B^{e}.
\]
The canonical action constraints identify \(\act\pi\cong \act\lambda\) and
\(\act\jmath\cong1_{E_{A,B}}\). Under these identifications the action equations are
\begin{align}
(\kappa_M\pi)\act(\phi_A\otimes1_M\otimes\phi_B)
 &=(\kappa_M\lambda)\act(1_{A_1}\otimes\kappa_M\otimes1_{B_1}),
 \label{eq:cgqg-background-module-associativity}\\
(\kappa_M\jmath)\act(\phi_A^0\otimes1_M\otimes\phi_B^0)
 &=1_M.
 \label{eq:cgqg-background-module-unit}
\end{align}
The first equation compares multiplication followed by action with two
successive actions; the second identifies the action of the units.

A map of quantum modules is a comonad transformation intertwining the actions.
\end{definition}

The direct action has a Kan-extension presentation whose inputs are endocomodules with source and target boundary maps.

\begin{definition}[Multitensors {\cite[\S3, Eqs.~(4)--(5)]{Chikhladze2010QuantumModules}}]
\label{def:cgqg-background-multitensors}
Let \(C_0,\ldots,C_n\) be comonoids and let
\[
 X_i:C_{i-1}^\circ\otimes C_i
       \rightsquigarrow C_{i-1}^\circ\otimes C_i
       \qquad(1\leq i\leq n)
\]
be endocomodules. For \(n\geq2\), let \(q_n\) contract the consecutive
middle boundaries by the evaluations \(e_{C_i}\). When the indicated
left Kan extension exists, set
\[
 T_n(X_1,\ldots,X_n)
 :=\operatorname{Lan}_{q_n}
       \bigl(q_n(X_1\otimes\cdots\otimes X_n)\bigr).
\]
Also set \(T_1(X)=X\) and
\(T_0(C)=\operatorname{Lan}_{n_C}n_C\), where
\(n_C:I\rightsquigarrow C^\circ\otimes C\) is coevaluation.
\end{definition}

For comonad inputs, the Kan-extension universal property induces a comonad
on \(T_n\). Under this correspondence, the direct action
\eqref{eq:cgqg-background-quantum-module-direct-action} is equivalently a
map of comonads
\begin{equation}
 \alpha_M:T_3(A_1,M,B_1)\longrightarrow M
 \label{eq:cgqg-background-quantum-module-action}
\end{equation}
satisfying the induced action equations. Section~\ref{sec:cgqg-strategy-quantum-modules} constructs the direct action and verifies these equations by identifying the corresponding spans.

\begin{proposition}[Span multitensors]
\label{prop:cgqg-background-span-multitensors}
Let \(X_i\) be an endocomodule on \(C_{P_{i-1}\times P_i}\), with
source and target boundary maps
\[
 d_i^{\diamond}(x_i)=(a_i^{\diamond}(x_i),b_i^{\diamond}(x_i)),\qquad\diamond\in\{L,R\}.
\]
For \(n\geq2\), the carrier of \(T_n(X_1,\ldots,X_n)\) is canonically
\begin{equation}
 K_n=
 \left\{(x_1,\ldots,x_n)\ \middle|\
 \begin{array}{l}
 b_i^{L}(x_i)=a_{i+1}^{L}(x_{i+1}),\\
 b_i^{R}(x_i)=a_{i+1}^{R}(x_{i+1})
 \end{array}
 \text{ for }1\leq i<n\right\},
 \label{eq:cgqg-background-span-multitensor}
\end{equation}
with boundary maps
\[
 d^{\diamond}(x_1,\ldots,x_n)
   =(a_1^{\diamond}(x_1),b_n^{\diamond}(x_n)).
\]
The carrier of \(T_0(C_P)\) is \(P\times P\), with source
\((p,p)\) and target \((q,q)\) at \((p,q)\). Its comultiplication has
witnesses
\begin{equation}
 (p,q)\longmapsto((p,z),(z,q))\quad(z\in P),
 \label{eq:cgqg-background-nullary-comultiplication}
\end{equation}
and its counit is supported on \(p=q\).
\end{proposition}
\begin{proof}
    \leavevmode
    \begin{enumerate}[beginpenalty=10000]
        \item \textbf{Express coactions by boundary maps.} For a coaction over a canonical diagonal comonoid, the counit equation identifies the coaction apex with its carrier. Thus coactions are graphs of boundary functions, and comodule maps are spans preserving these functions.

        \item \textbf{Construct the Kan extension.} The composite \(q_n(X_1\otimes\cdots\otimes X_n)\) imposes the
        target-boundary equations. A \(2\)-cell from this composite to \(Zq_n\)
        has support where the source labels also satisfy the
        source-boundary equations. On that locus, the source labels uniquely
        determine the contraction witness in \(Zq_n\).
        Forgetting or inserting this witness gives inverse operations
        \[
        \begin{aligned}
        &\{\text{boundary spans }q_n(X_1\otimes\cdots\otimes X_n)\Rightarrow Zq_n\}\\
        &\hspace{3em}\cong\{\text{boundary spans }K_n\Rightarrow Z\},
        \end{aligned}
        \]
        natural in \(Z\). This is the left Kan-extension universal property.

        \item \textbf{Construct the nullary carrier.} A \(2\)-cell \(n_P\Rightarrow Zn_P\) is supported at elements of \(Z\) with source \((p,p)\) and target \((q,q)\). Inserting or forgetting those two labels identifies it with a boundary-preserving span from the displayed \(P^2\)-carrier to \(Z\).

        \item \textbf{Compute the nullary comonad.} The induced comultiplication splits \((p,q)\) at an intermediate label \(z\), giving~\eqref{eq:cgqg-background-nullary-comultiplication}. Its counit selects \(p=q\). The two iterated comultiplications have witnesses \((p,z_1,z_2,q)\), and either counit removes the corresponding repeated endpoint. This is the matrix coalgebra in spans. \qedhere
    \end{enumerate}
\end{proof}

For quantum modules \(M:A\xnrightarrow{}B\) and
\(N:B\xnrightarrow{}C\), Chikhladze's intrinsic composition involves
the balancing pair
\begin{equation}
 \begin{tikzcd}[column sep=large]
 T_3(M,B_1,N)
   \arrow[r,shift left=.7ex,"d_r"]
   \arrow[r,shift right=.7ex,"d_l"']
 &T_2(M,N)\arrow[r]&N\bullet_B M .
 \end{tikzcd}
 \label{eq:cgqg-background-quantum-module-composition}
\end{equation}
Chikhladze's Theorem~11 constructs this composition on classes of modules closed under the indicated coequalizers and compatible multitensors~\cite{Chikhladze2010QuantumModules}. The generated composition of Section~\ref{sec:cgqg-composition-scope} is defined on pairs consisting of a distributor and its associated quantum module.

\begin{proposition}[Classical specialization {\cite[\S4]{Chikhladze2010QuantumModules}}]
\label{prop:cgqg-background-quantum-modules-set}
For \(\mathcal V=\mathbf{Set}\), quantum categories are ordinary
small categories, quantum modules are distributors, and
\eqref{eq:cgqg-background-quantum-module-composition} is the ordinary
coend composition
\[
 (N\circ M)(a,c)=\int^{b\in B}M(a,b)\times N(b,c).
\]
\end{proposition}
\begin{proof}
    \leavevmode
    \begin{enumerate}[beginpenalty=10000]
        \item \textbf{Identify the module data.} In the cartesian setting the comonad and action data reduce to a boundary-indexed set and its commuting endpoint actions. These are the data of a distributor.

        \item \textbf{Identify composition.} The balancing pair identifies the two ways of applying a middle arrow. This is exactly the defining relation of the displayed coend. \qedhere
    \end{enumerate}
\end{proof}

\section{Quantum symmetries of thin concurrent games}
\label{sec:cgqg-game-qg}
We now begin the main construction. Let \(A\) be a thin concurrent game. We use the polarized decomposition of its configuration symmetries to extract the structure from which the quantum groupoid associated to \(A\) will be built.

The construction proceeds as
\[
    A\longmapsto T(A)\longmapsto\WH_{\mathrm{sp}}(A)\longmapsto\QG_{\mathrm{sp}}(A),
\]
from polarized decomposition to vacant boxes, then incidence weak Hopf monoid, and finally the Pastro--Street quantum groupoid.

\subsection{Polarized boxes of a thin game}

A \textbf{polarized box} of \(A\) is a commuting square of selected symmetries
\begin{equation}
    \begin{tikzcd}[column sep=huge,row sep=huge]
        p
            \arrow[r,"h_t"]
            \arrow[d,"v_l"']
        & q
            \arrow[d,"v_r"]\\
        r
            \arrow[r,"h_b"']
        & s,
    \end{tikzcd}
    \qquad
    h_b\circ v_l=v_r\circ h_t,
    \label{diag:cgqg-game-resource-box}
\end{equation}
where
\[
    h_t,h_b\in\Sym^-(A),
    \qquad
    v_l,v_r\in\Sym^+(A).
\]
For a box \(X\) displayed as in~\eqref{diag:cgqg-game-resource-box}, set
\[
 \boxnw{X}=p,\quad\boxne{X}=q,\quad\boxsw{X}=r,\quad\boxse{X}=s.
\]
We keep these corner names separate from the boundary functions
\(\ell_X,r_X\) of an arbitrary bicomodule. Horizontal and vertical
pasting are the evident pastings of commuting squares.

\begin{definition}[Double and vacant double groupoids
{\cite[\S1.4, Def.~2.1 and Prop.~2.2]{AndruskiewitschNatale2005}}]
\label{def:cgqg-vacant-double-groupoid}
    A \textbf{double groupoid} consists of a set \(P\) of objects, horizontal and vertical groupoids
    \[
        H\rightrightarrows P,
        \qquad
        V\rightrightarrows P,
    \]
    and a set \(B\) of boxes
    \[
        \begin{tikzcd}[column sep=huge,row sep=huge]
            p \arrow[r,"h_t"] \arrow[d,"v_l"']
            & q \arrow[d,"v_r"] \\
            r \arrow[r,"h_b"']
            & s,
        \end{tikzcd}
    \]
    such that the boxes form groupoids under both horizontal and vertical composition, with objects \(V\) and \(H\), respectively, and the two compositions satisfy the interchange law.

    A double groupoid is \textbf{vacant} when every compatible pair of adjacent horizontal and vertical edges determines a unique box. Equivalently, every compatible southwest corner
    \begin{equation*}
        \begin{tikzcd}[column sep=huge,row sep=huge]
            p \arrow[d,"v"'] \arrow[r,dashed]
            & q \arrow[d,dashed] \\
            r \arrow[r,"h"']
            & s
        \end{tikzcd}
    \end{equation*}
    with \(v\in V\) and \(h\in H\) has a unique filler.
\end{definition}

\begin{proposition}[Polarized box double groupoid]
\label{prop:cgqg-game-vacant-boxes}
    The polarized boxes of \(A\) form a vacant double groupoid, denoted
    \begin{equation}
        T(A).
        \label{eq:cgqg-game-resource-double-groupoid}
    \end{equation}
    Its horizontal edge groupoid is \(\Sym^-(A)\), its vertical edge groupoid is \(\Sym^+(A)\), and its object set is \(\Conf(A)\).
\end{proposition}
\begin{proof}
    \leavevmode
    \begin{enumerate}[beginpenalty=10000]
        \item \textbf{Construct the two box compositions.}
        Pasting polarized boxes preserves the commuting-square equation.
        The horizontal edges compose in \(\Sym^-(A)\) and the vertical
        edges in \(\Sym^+(A)\), so the pasted squares remain polarized.

        \item \textbf{Identify identities and inverses.} Horizontal and vertical identities are the evident identity squares. Horizontal and vertical inverses are obtained by reversing the corresponding direction, and remain polarized because \(\Sym^-(A)\) and \(\Sym^+(A)\) are groupoids.

        \item \textbf{Verify the double-groupoid laws.} Associativity and the identity and inverse laws in each direction are inherited from composition in \(\Sym(A)\). The interchange law follows from the equality of the two ways of composing the boundary arrows of a pasted \(2\times2\) array of commuting squares. Hence the polarized boxes form a double groupoid.

        \item \textbf{Form the symmetry around a southwest corner.} Let
        \[
            v:p\longrightarrow r
            \quad\text{in }\Sym^+(A),
            \qquad
            h:r\longrightarrow s
            \quad\text{in }\Sym^-(A)
        \]
        be a compatible southwest corner. Their composite
        \[
            h\circ v:p\longrightarrow s
        \]
        is a selected symmetry of \(A\).

        \item \textbf{Apply polarized decomposition.} By Theorem~\ref{thm:cgqg-polarized-decomposition}, there are unique
        \(q\in\Conf(A)\),
        \[
            h':p\longrightarrow q
            \quad\text{in }\Sym^-(A),
            \qquad
            v':q\longrightarrow s
            \quad\text{in }\Sym^+(A),
        \]
        such that
        \[
            h\circ v=v'\circ h'.
        \]

        \item \textbf{Identify the unique box filler.} The preceding factorization is precisely the commuting square
        \[
            \begin{tikzcd}[column sep=huge,row sep=huge]
                p \arrow[r,"h'"] \arrow[d,"v"']
                & q \arrow[d,"v'"] \\
                r \arrow[r,"h"']
                & s.
            \end{tikzcd}
        \]
        Thus every compatible southwest corner has a unique polarized box filler. Therefore the double groupoid \(T(A)\) is vacant. \qedhere
    \end{enumerate}
\end{proof}

\begin{lemma}[Top-right determination]
\label{lem:cgqg-polarized-top-right-determination}
    Let
    \[
        h:p\longrightarrow q
        \quad\text{in }\Sym^-(A),
        \qquad
        v:q\longrightarrow s
        \quad\text{in }\Sym^+(A).
    \]
    There is a unique polarized box whose top horizontal edge is \(h\) and whose right vertical edge is \(v\).
\end{lemma}
\begin{proof}
    \leavevmode
    \begin{enumerate}[beginpenalty=10000]
        \item \textbf{Construct the filler.} Apply polarized decomposition to \((v\circ h)^{-1}\), say
        \[
            (v\circ h)^{-1}=u\circ k
        \]
        with \(k:s\to r\) negative and \(u:r\to p\) positive.
        Inverting gives \(v\circ h=k^{-1}\circ u^{-1}\), hence the required filler.

        \item \textbf{Verify uniqueness.} Any other filler gives, after inversion, a negative--positive decomposition of \((v\circ h)^{-1}\). Uniqueness follows from Theorem~\ref{thm:cgqg-polarized-decomposition}. \qedhere
    \end{enumerate}
\end{proof}

Write \(\Boxes_A\) for the set of polarized boxes. For compatible \(h\in\Sym^-(A)\) and \(v\in\Sym^+(A)\), write
\[
    h\trcorner v
\]
for the unique box of Lemma~\ref{lem:cgqg-polarized-top-right-determination}. Its left and bottom edges are denoted
\[
    \leftedge{h}{v},
    \qquad
    \bottomedge{h}{v},
\]
respectively.

For a polarized box \(X\), write \(X^{-1}\) for its total inverse, obtained by inverting in both double-groupoid directions; interchange makes the two orders equal.

\begin{lemma}[Polarized box calculus]
\label{lem:cgqg-vacant-cell-calculus}
    The following identities and factorization properties hold.
    \begin{enumerate}
        \item \textbf{Horizontal pasting.}
        \begin{equation}
            (h_2\trcorner v)
            \circ_{\mathrm h}
            (h_1\trcorner\leftedge{h_2}{v})
            =
            (h_2\circ h_1)\trcorner v.
            \label{eq:cgqg-vacant-cell-horizontal}
        \end{equation}

        \item \textbf{Vertical pasting.}
        \begin{equation}
            (\bottomedge{h}{v_1}\trcorner v_2)
            \circ_{\mathrm v}
            (h\trcorner v_1)
            =
            h\trcorner(v_2\circ v_1).
            \label{eq:cgqg-vacant-cell-vertical}
        \end{equation}

        \item \textbf{Vertical identities.}
        \begin{equation}
            \id_h^{\mathrm v}
            =
            h\trcorner\id_{t(h)}.
            \label{eq:cgqg-vacant-cell-identities}
        \end{equation}

        \item \textbf{Horizontal identities.}
        \begin{equation}
            \id_v^{\mathrm h}
            =
            \id_{s(v)}\trcorner v.
            \label{eq:cgqg-vacant-cell-horizontal-identity}
        \end{equation}

        \item \textbf{Total inversion.}
        \begin{equation}
            \bigl(h\trcorner v\bigr)^{-1}
            =
            \bigl(\bottomedge{h}{v}\bigr)^{-1}
            \trcorner
            \bigl(\leftedge{h}{v}\bigr)^{-1}.
            \label{eq:cgqg-vacant-cell-inversion}
        \end{equation}

        \item \textbf{Horizontal factorizations.} A horizontal factorization of \(h\trcorner v\) is uniquely determined by a factorization
        \[
            h=h_2\circ h_1
        \]
        in \(\Sym^-(A)\), and its two factors are
        \begin{equation}
            h_1\trcorner\leftedge{h_2}{v},
            \qquad
            h_2\trcorner v.
            \label{eq:cgqg-vacant-cell-horizontal-factorization}
        \end{equation}

        \item \textbf{Vertical factorizations.} A vertical factorization of \(h\trcorner v\) is uniquely determined by a factorization
        \[
            v=v_2\circ v_1
        \]
        in \(\Sym^+(A)\), and its two factors are
        \begin{equation}
            h\trcorner v_1,
            \qquad
            \bottomedge{h}{v_1}\trcorner v_2.
            \label{eq:cgqg-vacant-cell-vertical-factorization}
        \end{equation}

        \item \textbf{Factors of identity boxes.} Every horizontal factor of a vertical identity is a vertical identity, and every vertical factor of a horizontal identity is a horizontal identity.
    \end{enumerate}
\end{lemma}
\begin{proof}
    \leavevmode
    \begin{enumerate}[beginpenalty=10000]
        \item \textbf{Paste horizontally.} The two boxes
        \[
            h_1\trcorner\leftedge{h_2}{v},
            \qquad
            h_2\trcorner v
        \]
        are horizontally composable by definition of
        \(\leftedge{h_2}{v}\). Their paste has top edge
        \(h_2\circ h_1\) and right edge \(v\). By
        Lemma~\ref{lem:cgqg-polarized-top-right-determination}, it is
        therefore
        \[
            (h_2\circ h_1)\trcorner v.
        \]
        This proves~\eqref{eq:cgqg-vacant-cell-horizontal}.

        \item \textbf{Paste vertically.}
        The boxes
        \[
            h\trcorner v_1,
            \qquad
            \bottomedge{h}{v_1}\trcorner v_2
        \]
        are vertically composable by definition of
        \(\bottomedge{h}{v_1}\). Their paste has top edge \(h\) and
        right edge \(v_2\circ v_1\). Top-right determination therefore
        gives~\eqref{eq:cgqg-vacant-cell-vertical}.

        \item \textbf{Identify the two kinds of identity box.}
        The vertical identity on \(h\) has top edge \(h\) and right edge
        \(\id_{t(h)}\), so top-right determination gives
        \[
            \id_h^{\mathrm v}
            =
            h\trcorner\id_{t(h)}.
        \]
        Likewise, the horizontal identity on \(v\) has top edge
        \(\id_{s(v)}\) and right edge \(v\), hence
        \[
            \id_v^{\mathrm h}
            =
            \id_{s(v)}\trcorner v.
        \]

        \item \textbf{Compute total inversion.}
        Total inversion exchanges the top edge with the inverse of the
        bottom edge and the right edge with the inverse of the left edge.
        Thus the top and right edges of
        \((h\trcorner v)^{-1}\) are
        \[
            \bigl(\bottomedge{h}{v}\bigr)^{-1},
            \qquad
            \bigl(\leftedge{h}{v}\bigr)^{-1}.
        \]
        Top-right determination gives
        \eqref{eq:cgqg-vacant-cell-inversion}.

        \item \textbf{Classify horizontal factorizations.}
        Suppose
        \[
            h\trcorner v
            =
            X_2\circ_{\mathrm h}X_1.
        \]
        If \(h_1,h_2\) are the top edges of \(X_1,X_2\), respectively,
        then
        \[
            h=h_2\circ h_1.
        \]
        Since \(X_2\) has top edge \(h_2\) and right edge \(v\),
        top-right determination forces
        \[
            X_2=h_2\trcorner v.
        \]
        Its left edge is therefore
        \(\leftedge{h_2}{v}\), which is the right edge of \(X_1\).
        Hence
        \[
            X_1
            =
            h_1\trcorner\leftedge{h_2}{v}.
        \]
        Conversely, Step~1 shows that these two boxes paste to
        \(h\trcorner v\). This proves the horizontal factorization
        statement and its uniqueness.

        \item \textbf{Classify vertical factorizations.}
        Suppose
        \[
            h\trcorner v
            =
            X_2\circ_{\mathrm v}X_1.
        \]
        If \(v_1,v_2\) are the right edges of \(X_1,X_2\), respectively,
        then
        \[
            v=v_2\circ v_1.
        \]
        The upper factor has top edge \(h\) and right edge \(v_1\), so
        \[
            X_1=h\trcorner v_1.
        \]
        Its bottom edge is \(\bottomedge{h}{v_1}\), which is the top
        edge of \(X_2\). Hence
        \[
            X_2
            =
            \bottomedge{h}{v_1}\trcorner v_2.
        \]
        Conversely, Step~2 gives the required paste. This proves the
        vertical factorization statement and its uniqueness.

        \item \textbf{Classify factors of identity boxes.}
        Let
        \[
            \id_h^{\mathrm v}
            =
            h\trcorner\id_{t(h)}
        \]
        be a vertical identity, and suppose
        \[
            h=h_2\circ h_1.
        \]
        By the horizontal factorization classification, its right
        horizontal factor is
        \[
            h_2\trcorner\id_{t(h)}.
        \]
        Since \(t(h_2)=t(h)\), the identity formula of Step~3 gives
        \[
            h_2\trcorner\id_{t(h)}
            =
            h_2\trcorner\id_{t(h_2)}
            =
            \id_{h_2}^{\mathrm v}.
        \]
        The left edge of this box is therefore
        \[
            \id_{s(h_2)}
            =
            \id_{t(h_1)}.
        \]
        Hence the left horizontal factor is
        \[
            h_1\trcorner\id_{t(h_1)}
            =
            \id_{h_1}^{\mathrm v}.
        \]
        Thus every horizontal factor of a vertical identity is a
        vertical identity.

        Dually, let
        \[
            \id_v^{\mathrm h}
            =
            \id_{s(v)}\trcorner v
        \]
        be a horizontal identity and suppose
        \[
            v=v_2\circ v_1.
        \]
        The vertical factorization classification gives the upper
        factor
        \[
            \id_{s(v)}\trcorner v_1.
        \]
        Since \(s(v_1)=s(v)\), this is
        \[
            \id_{s(v_1)}\trcorner v_1
            =
            \id_{v_1}^{\mathrm h}.
        \]
        Its bottom edge is \(\id_{t(v_1)}=\id_{s(v_2)}\), so the lower
        factor is
        \[
            \id_{s(v_2)}\trcorner v_2
            =
            \id_{v_2}^{\mathrm h}.
        \]
        Thus every vertical factor of a horizontal identity is a
        horizontal identity. \qedhere
    \end{enumerate}
\end{proof}

\begin{proposition}[Tensor product of polarized boxes]
\label{prop:cgqg-resource-tensor}
    For thin concurrent games \(A,B\),
    \[
        T(A\otimes B)\cong T(A)\times T(B).
    \]
\end{proposition}
\begin{proof}
    \leavevmode
    \begin{enumerate}[beginpenalty=10000]
        \item \textbf{Factor the box data.} Equations~\eqref{eq:cgqg-conf-tensor} and~\eqref{eq:cgqg-tensor-symmetry} make polarized decomposition and every box operation componentwise. The resulting bijection preserves both groupoid structures. \qedhere
    \end{enumerate}
\end{proof}

\begin{proposition}[Negation as transposition]
\label{prop:cgqg-resource-negation}
    For every thin concurrent game \(A\),
    \begin{equation}
        T(A^\perp)\cong T(A)^{\mathsf t},
    \label{eq:cgqg-resource-negation}
    \end{equation}
    where transposition exchanges the horizontal and vertical edge groupoids and the two box compositions.
\end{proposition}
\begin{proof}
    \leavevmode
    \begin{enumerate}[beginpenalty=10000]
        \item \textbf{Transpose the box data.} Equation~\eqref{eq:cgqg-dual-symmetry} exchanges \(\Sym^-(A)\) and \(\Sym^+(A)\). Transposition preserves the commuting-square equation and exchanges the two box compositions. \qedhere
    \end{enumerate}
\end{proof}

\subsection{Box incidence and the weak Hopf monoid}

The two box directions now give transverse incidence operations. Put
\[
    M_{\mathrm v}(A)=\{(X,Y)\mid Y\circ_{\mathrm v}X\text{ is defined}\},
    \qquad
    M_{\mathrm h}(A)=\{(L,R)\mid R\circ_{\mathrm h}L\text{ is defined}\}.
\]
Define spans in \(\Sp\)

\begin{equation}
\begin{aligned}
m_A:\quad&
 \Boxes_A^2
 \xleftarrow{(X,Y)\mapsto(X,Y)}M_{\mathrm v}(A)
 \xrightarrow{(X,Y)\mapsto Y\circ_{\mathrm v}X}\Boxes_A,\\[.8em]
\eta_A:\quad&
 1\longleftarrow\Sym^-_1(A)
 \xrightarrow{h\mapsto\id_h^{\mathrm v}}\Boxes_A,\\[.8em]
\Delta_A:\quad&
 \Boxes_A
 \xleftarrow{(L,R)\mapsto R\circ_{\mathrm h}L}M_{\mathrm h}(A)
 \xrightarrow{(L,R)\mapsto(L,R)}\Boxes_A^2,\\[.8em]
\varepsilon_A:\quad&
 \Boxes_A\xleftarrow{v\mapsto\id_v^{\mathrm h}}
 \Sym^+_1(A)\longrightarrow1,\\[.8em]
S_A:\quad&
 \Boxes_A\xrightarrow{X\mapsto X^{-1}}\Boxes_A.
\end{aligned}
\label{eq:cgqg-game-wh-structure-maps}
\end{equation}
where the last map is regarded as its graph span.

For reference, in the symmetric category \(\Sp\) the weak unit and counit equations used below are
\begin{align}
    (\Delta_A\otimes1)\Delta_A\eta_A
    &=(1\otimes m_A\otimes1)(\Delta_A\eta_A\otimes\Delta_A\eta_A)
    \notag\\
    &=(1\otimes m_A\otimes1)(1\otimes c\otimes1)(\Delta_A\eta_A\otimes\Delta_A\eta_A),
    \label{eq:cgqg-game-weak-unit-equations}\\
    \varepsilon_A m_A(m_A\otimes1)
    &=(\varepsilon_A m_A\otimes\varepsilon_A m_A)(1\otimes\Delta_A\otimes1)
    \notag\\
    &=(\varepsilon_A m_A\otimes\varepsilon_A m_A)(1\otimes c\Delta_A\otimes1).
    \label{eq:cgqg-game-weak-counit-equations}
\end{align}

\begin{lemma}[Polarized rectangular decomposition]
\label{lem:cgqg-game-rectangular-decomposition}
    Let \(X,Y\in\Boxes_A\) be vertically composable, and suppose their vertical paste admits a horizontal factorization
    \[
        Y\circ_{\mathrm v}X
        =
        R\circ_{\mathrm h}L.
    \]
    Then there is a unique \(2\times2\) grid of polarized boxes
    \[
        \begin{matrix}
            X_1 & X_2\\
            Y_1 & Y_2
        \end{matrix}
    \]
    such that
    \[
        X=X_2\circ_{\mathrm h}X_1,
        \qquad
        Y=Y_2\circ_{\mathrm h}Y_1,
    \]
    the two columns are vertically composable, and
    \[
        L=Y_1\circ_{\mathrm v}X_1,
        \qquad
        R=Y_2\circ_{\mathrm v}X_2.
    \]
\end{lemma}
\begin{proof}
    \leavevmode
    \begin{enumerate}[beginpenalty=10000]
        \item \textbf{Fix the common paste.}
        Put
        \[
            P:=Y\circ_{\mathrm v}X
            =
            R\circ_{\mathrm h}L.
        \]

        \item \textbf{Factor the right column.}
        Since
        \[
            P=Y\circ_{\mathrm v}X,
        \]
        the right edge of \(P\) factors as the right edge of \(X\)
        followed by the right edge of \(Y\). Since
        \[
            P=R\circ_{\mathrm h}L,
        \]
        the right edge of \(P\) is the right edge of \(R\).
        By the vertical factorization statement of
        Lemma~\ref{lem:cgqg-vacant-cell-calculus}, there are therefore
        unique vertically composable boxes \(X_2,Y_2\) such that
        \[
            R=Y_2\circ_{\mathrm v}X_2,
        \]
        with the right edges of \(X_2\) and \(Y_2\) equal respectively
        to the right edges of \(X\) and \(Y\).

        \item \textbf{Factor the left column.}
        Since \(L\) and \(R\) are horizontally composable, the right
        edge of \(L\) is the left edge of \(R\). From
        \[
            R=Y_2\circ_{\mathrm v}X_2,
        \]
        that left edge factors as the left edge of \(X_2\) followed by
        the left edge of \(Y_2\). Applying the vertical factorization
        statement again gives unique vertically composable boxes
        \(X_1,Y_1\) such that
        \[
            L=Y_1\circ_{\mathrm v}X_1,
        \]
        where the right edge of \(X_1\) is the left edge of \(X_2\),
        and the right edge of \(Y_1\) is the left edge of \(Y_2\).
        Hence the two rows
        \[
            X_1,\ X_2
            \qquad\text{and}\qquad
            Y_1,\ Y_2
        \]
        are horizontally composable.

        \item \textbf{Identify the upper row with \(X\).}
        The horizontal paste
        \[
            X_2\circ_{\mathrm h}X_1
        \]
        has right edge equal to the right edge of \(X\). Its top edge is
        the horizontal composite of the top edges of \(L\) and \(R\),
        hence is the top edge of
        \[
            R\circ_{\mathrm h}L=P.
        \]
        Since \(P=Y\circ_{\mathrm v}X\), this is the top edge of \(X\).
        Thus \(X_2\circ_{\mathrm h}X_1\) and \(X\) have the same top
        and right edges. By
        Lemma~\ref{lem:cgqg-polarized-top-right-determination},
        \[
            X=X_2\circ_{\mathrm h}X_1.
        \]

        \item \textbf{Identify the lower row with \(Y\).}
        By interchange,
        \begin{align*}
            P
            &=
            R\circ_{\mathrm h}L
            \\
            &=
            \bigl(
                Y_2\circ_{\mathrm v}X_2
            \bigr)
            \circ_{\mathrm h}
            \bigl(
                Y_1\circ_{\mathrm v}X_1
            \bigr)
            \\
            &=
            \bigl(
                Y_2\circ_{\mathrm h}Y_1
            \bigr)
            \circ_{\mathrm v}
            \bigl(
                X_2\circ_{\mathrm h}X_1
            \bigr)
            \\
            &=
            \bigl(
                Y_2\circ_{\mathrm h}Y_1
            \bigr)
            \circ_{\mathrm v}X.
        \end{align*}
        But also
        \[
            P=Y\circ_{\mathrm v}X.
        \]
        Since the boxes form a groupoid under vertical composition,
        cancellation of the common upper factor \(X\) gives
        \[
            Y=Y_2\circ_{\mathrm h}Y_1.
        \]
        By construction,
        \[
            L=Y_1\circ_{\mathrm v}X_1,
            \qquad
            R=Y_2\circ_{\mathrm v}X_2.
        \]

        \item \textbf{Prove uniqueness.}
        The right-edge factorization of \(R\) determined by the given
        vertical factorization
        \[
            P=Y\circ_{\mathrm v}X
        \]
        uniquely determines \(X_2,Y_2\) by the vertical factorization
        classification. Their left edges then uniquely determine the
        vertical factorization of \(L\), hence uniquely determine
        \(X_1,Y_1\). Therefore the entire \(2\times2\) grid is unique. \qedhere
    \end{enumerate}
\end{proof}

\begin{lemma}[Incidence weak bimonoid]
\label{lem:cgqg-game-incidence-weak-bimonoid}
    The spans
    \[
        m_A,\eta_A,\Delta_A,\varepsilon_A
    \]
    make \(\Boxes_A\) a weak bimonoid in \(\Sp\).

    Its canonical target and source counital maps are represented respectively by the spans
    \begin{equation}
        \Boxes_A
        \xleftarrow{\;(v,h)\mapsto\id_v^{\mathrm h}\;}
        \bigl\{
            (v,h)\in\Sym^+_1(A)\times\Sym^-_1(A)
            \mid s(v)=s(h)
        \bigr\}
        \xrightarrow{\;(v,h)\mapsto\id_h^{\mathrm v}\;}
        \Boxes_A
        \label{eq:cgqg-game-target-counital-map}
    \end{equation}
    and
    \begin{equation}
        \Boxes_A
        \xleftarrow{\;(v,h)\mapsto\id_v^{\mathrm h}\;}
        \bigl\{
            (v,h)\in\Sym^+_1(A)\times\Sym^-_1(A)
            \mid t(v)=t(h)
        \bigr\}
        \xrightarrow{\;(v,h)\mapsto\id_h^{\mathrm v}\;}
        \Boxes_A.
        \label{eq:cgqg-game-source-counital-map}
    \end{equation}
\end{lemma}
\begin{proof}
    \leavevmode
    \begin{enumerate}[beginpenalty=10000]
        \item \textbf{Verify the monoid laws.}
        The multiplication span records vertical composition of boxes,
        and \(\eta_A\) records vertical identity boxes. The apex of
        either parenthesization of the triple product is the set of
        vertically composable triples
        \[
            (X_1,X_2,X_3),
        \]
        and both composites send such a triple to its total vertical
        paste. Associativity follows from associativity in the vertical
        box groupoid.

        For the unit laws, pulling back against \(\eta_A\) inserts the
        appropriate vertical identity box. The two resulting composites
        are therefore the identity span on \(\Boxes_A\). Thus
        \((m_A,\eta_A)\) is a monoid.

        \item \textbf{Verify the comonoid laws.}
        Dually, \(\Delta_A\) records horizontal factorizations and
        \(\varepsilon_A\) records horizontal identity boxes. Both
        iterated coproducts have as apex the set of horizontally
        composable triples
        \[
            (X_1,X_2,X_3)
        \]
        together with their common total horizontal paste. Associativity
        in the horizontal box groupoid identifies the two apexes
        canonically, with the same exterior legs.

        Pullback against \(\varepsilon_A\) forces the corresponding
        factor to be a horizontal identity, and the horizontal identity
        laws then identify either counit composite with the identity
        span on \(\Boxes_A\). Hence \((\Delta_A,\varepsilon_A)\) is a
        comonoid.

        \item \textbf{Prove multiplicativity.}
        A witness for
        \[
            \Delta_A m_A
        \]
        consists of a vertically composable pair \(X,Y\), together with
        a horizontal factorization
        \[
            Y\circ_{\mathrm v}X
            =
            R\circ_{\mathrm h}L.
        \]
        By Lemma~\ref{lem:cgqg-game-rectangular-decomposition}, this is
        uniquely equivalent to a \(2\times2\) grid
        \[
            \begin{matrix}
                X_1 & X_2\\
                Y_1 & Y_2
            \end{matrix}
        \]
        whose rows and columns are composable. Such a grid is exactly a
        witness for
        \[
            (m_A\otimes m_A)
            (1\otimes c\otimes1)
            (\Delta_A\otimes\Delta_A).
        \]
        Under this correspondence the four exterior boxes are unchanged.
        Hence the two spans are equal.

        \item \textbf{Verify the weak unit equations.}
        Put
        \[
            U^-_n
            :=
            \left\{
                (h_1,\ldots,h_n)
                \in\Sym^-_1(A)^n
                \;\middle|\;
                h_n\circ\cdots\circ h_1
                \text{ is defined}
            \right\}.
        \]
        By the factor-of-identity statement in
        Lemma~\ref{lem:cgqg-vacant-cell-calculus}, a horizontal
        factorization of a vertical identity
        \[
            \id_h^{\mathrm v}
        \]
        consists entirely of vertical identities. More precisely,
        horizontal factorizations of \(\id_h^{\mathrm v}\) are in
        bijection with factorizations
        \[
            h=h_2\circ h_1,
        \]
        and the two factors are
        \[
            \id_{h_1}^{\mathrm v},
            \qquad
            \id_{h_2}^{\mathrm v}.
        \]
        Consequently \(\Delta_A\eta_A\) is represented by the span whose
        witness set is \(U^-_2\), with output
        \[
            (h_1,h_2)
            \longmapsto
            \bigl(
                \id_{h_1}^{\mathrm v},
                \id_{h_2}^{\mathrm v}
            \bigr).
        \]

        Iterating the coproduct shows that either coassociative
        threefold factorization of the unit has witness set \(U^-_3\),
        with output
        \[
            (h_1,h_2,h_3)
            \longmapsto
            \bigl(
                \id_{h_1}^{\mathrm v},
                \id_{h_2}^{\mathrm v},
                \id_{h_3}^{\mathrm v}
            \bigr).
        \]

        Consider now either product-side composite occurring in the two
        weak-unit equations of
        \cite[Def.~1.1]{PastroStreet2009}. A witness consists of two
        horizontal factorizations of vertical identity boxes together
        with a vertical multiplication of the two middle factors.
        Every factor is a vertical identity by the preceding paragraph,
        and two vertical identity boxes are vertically composable
        precisely when their underlying horizontal edges agree. Thus
        the two middle witnesses are forced to carry the same horizontal
        edge.

        Hence in either weak-unit composite the data reduce uniquely to
        a composable triple
        \[
            (h_1,h_2,h_3)\in U^-_3,
        \]
        and the exterior output is again
        \[
            \bigl(
                \id_{h_1}^{\mathrm v},
                \id_{h_2}^{\mathrm v},
                \id_{h_3}^{\mathrm v}
            \bigr).
        \]
        The symmetry appearing in the second weak-unit equation only
        exchanges the two middle witness coordinates before the same
        equality is imposed. Since \(\Sp\) is symmetric, the resulting
        witness set and exterior legs are the same. Thus both weak-unit
        equations hold.

        \item \textbf{Verify the weak counit equations.}
        Put
        \[
            U^+_n
            :=
            \left\{
                (v_1,\ldots,v_n)
                \in\Sym^+_1(A)^n
                \;\middle|\;
                v_n\circ\cdots\circ v_1
                \text{ is defined}
            \right\}.
        \]
        A witness for the counit of a total vertical paste is a
        vertically composable family of boxes whose total paste is a
        horizontal identity. By the factor-of-identity statement in
        Lemma~\ref{lem:cgqg-vacant-cell-calculus}, every vertical factor
        of such a paste is itself a horizontal identity. Thus a
        threefold witness is uniquely of the form
        \[
            \bigl(
                \id_{v_1}^{\mathrm h},
                \id_{v_2}^{\mathrm h},
                \id_{v_3}^{\mathrm h}
            \bigr)
        \]
        for a composable triple
        \[
            (v_1,v_2,v_3)\in U^+_3.
        \]

        On either product-side composite in the weak-counit equations of
        \cite[Def.~1.1]{PastroStreet2009}, the middle box is first
        horizontally factorized and the two resulting vertical pastes
        are then required by the two counits to be horizontal identities.
        Each factor in those vertical pastes is therefore a horizontal
        identity. Since two horizontally composable horizontal identity
        boxes must be identities on the same vertical edge, the inserted
        horizontal factorization is uniquely the identity factorization
        on that common edge.

        Hence both product-side composites again have witness set
        \(U^+_3\), with the same exterior legs as the total-paste side.
        The second weak-counit equation differs only by the symmetry
        interchanging the two middle coproduct coordinates, and gives
        the same witness set in the symmetric category \(\Sp\).
        Therefore both weak-counit equations hold.

        \item \textbf{Compute the target counital map.}
        Expand the canonical target counital composite of the weak
        bimonoid. By the identity calculations above, a witness reduces
        uniquely to a pair
        \[
            (v,h)
            \in
            \Sym^+_1(A)\times\Sym^-_1(A)
        \]
        satisfying
        \[
            s(v)=s(h).
        \]
        Its input box is
        \[
            \id_v^{\mathrm h},
        \]
        and its output box is
        \[
            \id_h^{\mathrm v}.
        \]
        Hence the target counital map is represented by
        \[
            \Boxes_A
            \xleftarrow{\;(v,h)\mapsto\id_v^{\mathrm h}\;}
            \bigl\{
                (v,h)
                \mid
                s(v)=s(h)
            \bigr\}
            \xrightarrow{\;(v,h)\mapsto\id_h^{\mathrm v}\;}
            \Boxes_A,
        \]
        which is~\eqref{eq:cgqg-game-target-counital-map}.

        \item \textbf{Compute the source counital map.}
        Reflecting the preceding identity calculation exchanges sources
        with targets. The surviving witnesses are therefore the pairs
        \[
            (v,h)
            \in
            \Sym^+_1(A)\times\Sym^-_1(A)
        \]
        satisfying
        \[
            t(v)=t(h),
        \]
        again with input \(\id_v^{\mathrm h}\) and output
        \(\id_h^{\mathrm v}\). This is precisely
        \eqref{eq:cgqg-game-source-counital-map}. \qedhere
    \end{enumerate}
\end{proof}

\begin{theorem}[Quantum symmetry weak Hopf monoid]
\label{thm:cgqg-span-weak-hopf}
    For every thin concurrent game \(A\), the polarized box object carries a canonical involutive weak Hopf monoid structure
    \begin{equation}
        \WH_{\mathrm{sp}}(A):=\Boxes_A
        \label{eq:cgqg-game-wh-object}
    \end{equation}
    in \(\Sp\), with structure maps \(m_A,\eta_A,\Delta_A,\varepsilon_A,S_A\) above and
    \[
        S_A^2=1_{\Boxes_A}.
    \]
\end{theorem}
\begin{proof}
    \leavevmode
    \begin{enumerate}[beginpenalty=10000]
        \item \textbf{Use the incidence weak bimonoid.} By Lemma~\ref{lem:cgqg-game-incidence-weak-bimonoid},
        \[
            (\Boxes_A,m_A,\eta_A,\Delta_A,\varepsilon_A)
        \]
        is a weak bimonoid. It remains to verify the three antipode equations for total box inversion.

        \item \textbf{Set up the target antipode convolution.} Write an input box as
        \[
            X=h\trcorner v.
        \]
        By Lemma~\ref{lem:cgqg-vacant-cell-calculus}, a horizontal factorization of \(X\) is uniquely determined by
        \[
            h=h_2\circ h_1
        \]
        and has factors
        \[
            Z
            =
            h_1\trcorner\leftedge{h_2}{v},
            \qquad
            Y
            =
            h_2\trcorner v.
        \]
        A witness for
        \[
            m_A(1\otimes S_A)\Delta_A
        \]
        requires \(Z\) and \(S_A(Y)\) to be vertically composable. Since the top edge of \(S_A(Y)\) is
        \[
            \bigl(\bottomedge{h_2}{v}\bigr)^{-1},
        \]
        this condition is
        \[
            \bottomedge{h_1}{\leftedge{h_2}{v}}
            =
            \bigl(\bottomedge{h_2}{v}\bigr)^{-1}.
        \]

        \item \textbf{Force the input to be a horizontal identity.} The bottom edge of the horizontal paste \(X=Y\circ_{\mathrm h}Z\) is
        \[
            \bottomedge{h_2}{v}
            \circ
            \bottomedge{h_1}{\leftedge{h_2}{v}}.
        \]
        Hence the preceding composability condition is equivalent to
        \[
            \bottomedge{h}{v}=\id.
        \]
        The commuting-square equation for \(X\) then gives
        \[
            \leftedge{h}{v}
            =
            v\circ h.
        \]
        The left-hand side lies in \(\Sym^+(A)\). Comparing the polarized decomposition
        \[
            v\circ h
        \]
        with the purely positive decomposition
        \[
            \leftedge{h}{v}\circ\id
        \]
        and using uniqueness in Theorem~\ref{thm:cgqg-polarized-decomposition} forces
        \[
            h=\id,
            \qquad
            \leftedge{h}{v}=v.
        \]
        Thus the convolution is supported precisely on horizontal identity boxes.

        \item \textbf{Identify the target counital map.} If \(X=\id_v^{\mathrm h}\), then a factorization
        \[
            \id=h_2\circ h_1
        \]
        has \(h_2=h_1^{-1}\), and the corresponding antipode witness exists uniquely. Its vertical paste is \(\id_{h_1}^{\mathrm v}\). Therefore
        \[
            m_A(1\otimes S_A)\Delta_A
        \]
        is exactly the target counital span~\eqref{eq:cgqg-game-target-counital-map}.

        \item \textbf{Identify the source counital map.} Reflecting the preceding calculation gives
        \[
            m_A(S_A\otimes1)\Delta_A
        \]
        equal to the source counital span~\eqref{eq:cgqg-game-source-counital-map}.

        \item \textbf{Set up the third antipode equation.} A threefold horizontal factorization of
        \[
            X=h\trcorner v
        \]
        is uniquely determined by
        \[
            h=h_3\circ h_2\circ h_1
        \]
        and has factors
        \[
            A_1
            =
            h_1
            \trcorner
            \leftedge{h_2}{\leftedge{h_3}{v}},
            \qquad
            A_2
            =
            h_2\trcorner\leftedge{h_3}{v},
            \qquad
            A_3
            =
            h_3\trcorner v.
        \]
        A witness for
        \[
            m_A^{(3)}
            (S_A\otimes1\otimes S_A)
            \Delta_A^{(3)}
        \]
        requires \(S_A(A_1),A_2,S_A(A_3)\) to be vertically composable.

        \item \textbf{Determine the unique threefold witness.} The first composability condition gives
        \[
            h_2=h_1^{-1}.
        \]
        The second gives
        \[
            \bottomedge{h_2}{\leftedge{h_3}{v}}
            =
            \bigl(\bottomedge{h_3}{v}\bigr)^{-1}.
        \]
        Hence the horizontal paste
        \[
            A_3\circ_{\mathrm h}A_2
            =
            (h_3\circ h_2)\trcorner v
        \]
        has identity bottom edge. By Step~3,
        \[
            h_3\circ h_2=\id.
        \]
        Therefore
        \[
            h_2=h_1^{-1},
            \qquad
            h_3=h_1,
            \qquad
            h=h_1.
        \]
        Moreover \(A_3=X\), and the left vertical edge of the horizontal identity \(A_3\circ_{\mathrm h}A_2\) is \(v\); hence the right edge of \(A_1\) is \(v\), so \(A_1=X\) as well. Thus there is exactly one such threefold witness over each input box.

        \item \textbf{Compute the threefold output.} The total vertical paste has top edge
        \[
            \bigl(\bottomedge{h}{v}\bigr)^{-1}
        \]
        and right edge
        \[
            \bigl(\leftedge{h}{v}\bigr)^{-1}
            \circ
            \leftedge{h}{v}
            \circ
            \bigl(\leftedge{h}{v}\bigr)^{-1}
            =
            \bigl(\leftedge{h}{v}\bigr)^{-1}.
        \]
        By top-right determination, this box is precisely
        \[
            \bigl(h\trcorner v\bigr)^{-1}
            =
            S_A(X).
        \]
        Hence
        \[
            m_A^{(3)}
            (S_A\otimes1\otimes S_A)
            \Delta_A^{(3)}
            =
            S_A.
        \]

        \item \textbf{Conclude involutivity.} Total inversion in a double groupoid is involutive, so
        \[
            S_A^2=1_{\Boxes_A}.
        \]
        The three antipode equations therefore make \(\WH_{\mathrm{sp}}(A)\) an involutive weak Hopf monoid. \qedhere
    \end{enumerate}
\end{proof}

When \(A\) is fixed, we suppress the subscript \(A\) on the weak Hopf structure maps \(m,\eta,\Delta,\varepsilon,S\).

\begin{example}[A finite game with symmetry of both polarities]
\label{ex:cgqg-mixed-polarity-boxes}
Let \(A\) have four concurrent events \(n_1,n_2,p_1,p_2\), with
the \(n_i\) negative and the \(p_i\) positive, all finite subsets
consistent, and all polarity-preserving bijections selected. Negative
symmetries fix the positive events pointwise, and positive symmetries
fix the negative events pointwise. At the full configuration
\(x=\{n_1,n_2,p_1,p_2\}\), let \(h\) exchange the two negative
events and let \(v\) exchange the two positive events. Write
\(X_{ij}=h^i\trcorner v^j\), with indices in \(\mathbb Z/2\mathbb Z\).
The four boxes at this configuration satisfy
\[
 m_A(X_{ij},X_{k\ell})=X_{i,j+\ell}\quad\text{exactly when }i=k,
\]
and the horizontal factorization witnesses of \(X_{ij}\) are
\[
 (X_{aj},X_{bj})\quad\text{with }a+b=i.
\]
In particular, \(m_A(X_{11},X_{11})=X_{10}\), while
\(\Delta_A(X_{11})\) has the two witnesses
\((X_{01},X_{11})\) and \((X_{11},X_{01})\).
\end{example}
\begin{proof}
    \leavevmode
    \begin{enumerate}[beginpenalty=10000]
        \item \textbf{Check the polarized game structure.}
        Restriction and extension hold by extending the bijections
        separately on each polarity. The polarized families intersect
        in the identities and satisfy their respective receptivity
        conditions, so \(A\) is a thin concurrent game.

        \item \textbf{Display the nontrivial polarized box.}
        The permutations \(h\) and \(v\) commute because they act on
        disjoint sets of events. Thus
        \[
        \begin{tikzcd}[column sep=huge,row sep=large]
          x\arrow[r,"h"]\arrow[d,"v"']&x\arrow[d,"v"]\\
          x\arrow[r,"h"']&x
        \end{tikzcd}
        \qquad X_{11}=h\trcorner v.
        \]
        More generally, \(X_{ij}\) has top and bottom edges \(h^i\)
        and left and right edges \(v^j\). Every symmetry out of the
        full configuration has that same full configuration as target,
        so its factorization witnesses remain in this component.

        \item \textbf{Compute multiplication and comultiplication.}
        Vertical pasting requires equality of the intermediate
        horizontal edges, hence \(i=k\), and composes the positive
        edges to \(v^{j+\ell}\). Horizontal pasting requires equal
        positive edges and multiplies the negative edges; consequently
        the factorizations are exactly the displayed pairs with
        \(a+b=i\). This proves the formulas and the two particular
        calculations.

        \item \textbf{Identify the remaining operations.}
        In this component the unit witnesses are \(X_{00},X_{10}\),
        the counit is supported on \(X_{00},X_{01}\), and total
        inversion fixes every \(X_{ij}\). The unit and counit therefore
        select different box directions even in this finite example.
        \qedhere
    \end{enumerate}
\end{proof}

\subsection{The counital base}

For reference, the relation between the two counital conventions and
the base-valued maps is
\[
\begin{array}{c|c}
\text{Pastro--Street notation}&\text{configuration-span presentation}\\\hline
r_A^{\mathrm{PS}}&\varepsilon_t=i_t r_t\\
t_A^{\mathrm{PS}}&\varepsilon_s=i_s r_s\\
s_A^{\mathrm{PS}}&i_s r_t\\
\bar s_A,\bar t_A&r_t,r_s
\end{array}
\]
The spans below construct these presentations and prove the displayed
identities.

Let \(\Theta_a\) denote the double identity at \(a\in\Conf(A)\). Define the spans
\begin{equation}
    \begin{aligned}
        i_t:\Conf(A)&\longrightarrow\WH_{\mathrm{sp}}(A),
        \qquad
        &
        \begin{tikzcd}
            & \Sym^-_1(A)
                \arrow[dl,"s"']
                \arrow[dr,"{h\mapsto\id_h^{\mathrm v}}"]
            & \\
            \Conf(A)
            &&
            \Boxes_A
        \end{tikzcd}
        \\[1.5em]
        i_s:\Conf(A)&\longrightarrow\WH_{\mathrm{sp}}(A),
        &
        \begin{tikzcd}
            & \Sym^-_1(A)
                \arrow[dl,"t"']
                \arrow[dr,"{h\mapsto\id_h^{\mathrm v}}"]
            & \\
            \Conf(A)
            &&
            \Boxes_A
        \end{tikzcd}
        \\[1.5em]
        r_t:\WH_{\mathrm{sp}}(A)&\longrightarrow\Conf(A),
        &
        \begin{tikzcd}
            & \Sym^+_1(A)
                \arrow[dl,"{v\mapsto\id_v^{\mathrm h}}"']
                \arrow[dr,"s"]
            & \\
            \Boxes_A
            &&
            \Conf(A)
        \end{tikzcd}
        \\[1.5em]
        r_s:\WH_{\mathrm{sp}}(A)&\longrightarrow\Conf(A),
        &
        \begin{tikzcd}
            & \Sym^+_1(A)
                \arrow[dl,"{v\mapsto\id_v^{\mathrm h}}"']
                \arrow[dr,"t"]
            & \\
            \Boxes_A
            &&
            \Conf(A).
        \end{tikzcd}
    \end{aligned}
\label{eq:cgqg-span-counital-retract-maps}
\end{equation}

\begin{lemma}[Counital retract spans]
\label{lem:cgqg-game-counital-retract-spans}
    The composites of the configuration retract spans are
    \begin{equation}
        i_tr_t
        =
        \left[
        \Boxes_A
        \xleftarrow{\;(v,h)\mapsto\id_v^{\mathrm h}\;}
        \bigl\{
            (v,h)\in\Sym^+_1(A)\times\Sym^-_1(A)
            \mid s(v)=s(h)
        \bigr\}
        \xrightarrow{\;(v,h)\mapsto\id_h^{\mathrm v}\;}
        \Boxes_A
        \right],
        \label{eq:cgqg-game-itrt-span}
    \end{equation}
    and
    \begin{equation}
        i_sr_s
        =
        \left[
        \Boxes_A
        \xleftarrow{\;(v,h)\mapsto\id_v^{\mathrm h}\;}
        \bigl\{
            (v,h)\in\Sym^+_1(A)\times\Sym^-_1(A)
            \mid t(v)=t(h)
        \bigr\}
        \xrightarrow{\;(v,h)\mapsto\id_h^{\mathrm v}\;}
        \Boxes_A
        \right].
        \label{eq:cgqg-game-isrs-span}
    \end{equation}
    Hence \(i_tr_t\) and \(i_sr_s\) are respectively the target and source counital maps of the incidence weak bimonoid.
\end{lemma}
\begin{proof}
    \leavevmode
    \begin{enumerate}[beginpenalty=10000]
        \item \textbf{Compose \(r_t\) with \(i_t\).} The composite
        \[
            \Boxes_A
            \xrightarrow{r_t}
            \Conf(A)
            \xrightarrow{i_t}
            \Boxes_A
        \]
        is formed by pulling back the two apex maps
        \[
            s:\Sym^+_1(A)\longrightarrow\Conf(A),
            \qquad
            s:\Sym^-_1(A)\longrightarrow\Conf(A).
        \]
        Its apex is therefore
        \[
            \bigl\{
                (v,h)\in\Sym^+_1(A)\times\Sym^-_1(A)
                \mid s(v)=s(h)
            \bigr\},
        \]
        with exterior legs
        \[
            (v,h)\longmapsto\id_v^{\mathrm h},
            \qquad
            (v,h)\longmapsto\id_h^{\mathrm v}.
        \]
        This gives \eqref{eq:cgqg-game-itrt-span}.

        \item \textbf{Compose \(r_s\) with \(i_s\).} Likewise, the pullback defining
        \[
            \Boxes_A
            \xrightarrow{r_s}
            \Conf(A)
            \xrightarrow{i_s}
            \Boxes_A
        \]
        imposes
        \[
            t(v)=t(h),
        \]
        giving \eqref{eq:cgqg-game-isrs-span}.

        \item \textbf{Identify the counital maps.} These are exactly the target and source counital spans computed in Lemma~\ref{lem:cgqg-game-incidence-weak-bimonoid}. \qedhere
    \end{enumerate}
\end{proof}

\begin{proposition}[Explicit counital maps]
\label{prop:cgqg-incidence-counital-maps}
    For \(\WH_{\mathrm{sp}}(A)\),
    \begin{equation}
        m(1\otimes S)\Delta=i_tr_t,
        \qquad
        m(S\otimes1)\Delta=i_sr_s.
        \label{eq:cgqg-explicit-counital-incidence}
    \end{equation}
    Moreover
    \begin{equation}
        r_ti_t=r_si_s=r_si_t=r_ti_s=1_{\Conf(A)},
        \label{eq:cgqg-base-cross-splittings}
    \end{equation}
    \begin{align}
        Si_t&=i_s,
        &
        Si_s&=i_t,
        \label{eq:cgqg-antipode-exchanges-base-inclusions}\\
        r_tS&=r_s,
        &
        r_sS&=r_t.
        \label{eq:cgqg-antipode-exchanges-base-retractions}
    \end{align}
\end{proposition}

\begin{proof}
    \leavevmode
    \begin{enumerate}[beginpenalty=10000]
        \item \textbf{Identify the two counital convolutions.} By Theorem~\ref{thm:cgqg-span-weak-hopf}, the antipode satisfies
        \[
            m(1\otimes S)\Delta=\varepsilon_t,
            \qquad
            m(S\otimes1)\Delta=\varepsilon_s.
        \]
        Lemma~\ref{lem:cgqg-game-counital-retract-spans} identifies these counital maps with
        \[
            \varepsilon_t=i_tr_t,
            \qquad
            \varepsilon_s=i_sr_s.
        \]
        This proves~\eqref{eq:cgqg-explicit-counital-incidence}.

        \item \textbf{Compute the four configuration composites.} Consider, for example, \(r_ti_t\). Its pullback apex consists of pairs
        \[
            (h,v)\in\Sym^-_1(A)\times\Sym^+_1(A)
        \]
        satisfying
        \[
            \id_h^{\mathrm v}
            =
            \id_v^{\mathrm h}.
        \]
        The same equation occurs for each of the other three composites, with the exterior configuration labels read using the appropriate source or target maps.

        \item \textbf{Identify the surviving boxes.} A box which is simultaneously a vertical identity and a horizontal identity is necessarily a double identity
        \[
            \Theta_a
        \]
        at a unique \(a\in\Conf(A)\). Conversely every \(\Theta_a\) gives such a witness. In each of the four composites the input and output configuration are therefore both \(a\), with a unique witness. Hence
        \[
            r_ti_t
            =
            r_si_s
            =
            r_si_t
            =
            r_ti_s
            =
            1_{\Conf(A)}.
        \]

        \item \textbf{Track inversion on the inclusions.} Total box inversion sends a vertical identity to
        \[
            S(\id_h^{\mathrm v})
            =
            \id_{h^{-1}}^{\mathrm v}.
        \]
        Since
        \[
            s(h^{-1})=t(h),
            \qquad
            t(h^{-1})=s(h),
        \]
        inversion exchanges the source- and target-labelled inclusions:
        \[
            Si_t=i_s,
            \qquad
            Si_s=i_t.
        \]

        \item \textbf{Track inversion on the retractions.} Likewise,
        \[
            S(\id_v^{\mathrm h})
            =
            \id_{v^{-1}}^{\mathrm h}.
        \]
        On the common apex \(\Sym^+_1(A)\),
        \[
            r_tS(\id_v^{\mathrm h})
            =
            s(v^{-1})
            =
            t(v)
            =
            r_s(\id_v^{\mathrm h}),
        \]
        so
        \[
            r_tS=r_s.
        \]
        Applying \(S^2=1_{\Boxes_A}\) gives
        \[
            r_sS=r_t.\qedhere
        \]
    \end{enumerate}
\end{proof}

By Proposition~\ref{prop:cgqg-background-frobenius-base}, the configuration set \(\Conf(A)\) carries the canonical separable commutative Frobenius monoid
\[
    C_{\Conf(A)}
    :=
    \bigl(
        \Conf(A),
        \mu_{\Conf(A)},
        \eta_{\Conf(A)},
        \delta_{\Conf(A)},
        \epsilon_{\Conf(A)}
    \bigr)
\]
in \(\Sp\).

\begin{proposition}[Configuration Frobenius retract]
\label{prop:cgqg-configuration-frobenius-retract}
    Transport the Pastro--Street target counital retract through
    \[
        \Conf(A)
        \xrightarrow{i_s}
        \bigl(
            \WH_{\mathrm{sp}}(A),
            t_A^{\mathrm{PS}}
        \bigr)
        \xrightarrow{r_s}
        \Conf(A).
    \]
    The transported Frobenius structure is
    \begin{equation*}
\begin{aligned}
\mu_A^C:\quad&
 \Conf(A)^2\xleftarrow{\Delta_{\Conf(A)}}\Conf(A)
 \xrightarrow{1_{\Conf(A)}}\Conf(A),\\[.8em]
\eta_A^C:\quad&
 1\longleftarrow\Conf(A)\xrightarrow{1_{\Conf(A)}}\Conf(A),\\[.8em]
\delta_A^C:\quad&
 \Conf(A)\xleftarrow{1_{\Conf(A)}}\Conf(A)
 \xrightarrow{\Delta_{\Conf(A)}}\Conf(A)^2,\\[.8em]
\epsilon_A^C:\quad&
 \Conf(A)\xleftarrow{1_{\Conf(A)}}\Conf(A)\longrightarrow1.
\end{aligned}
\end{equation*}
    Hence the transported separable commutative Frobenius monoid is
    precisely \(C_{\Conf(A)}\).
\end{proposition}
\begin{proof}
    \leavevmode
    \begin{enumerate}[beginpenalty=10000]
        \item \textbf{Set up the transported Frobenius structure.}
        By \cite[Proposition~1.4]{PastroStreet2009}, transporting the target Frobenius structure along \(i_s,r_s\) gives the four composites
        \[
            r_sm(i_s\otimes i_s),
            \qquad
            r_s\eta,
            \qquad
            (r_s\otimes r_s)\Delta i_s,
            \qquad
            \varepsilon i_s.
        \]
        We calculate them as spans.

        \item \textbf{Calculate the transported multiplication.} The multiplication is
        \[
            r_sm(i_s\otimes i_s).
        \]
        A multiplication witness begins with vertical identities \(\id_{h_1}^{\mathrm v}\) and \(\id_{h_2}^{\mathrm v}\). Vertical composability forces
        \[
            h_1=h_2=:h,
        \]
        and their paste is \(\id_h^{\mathrm v}\). The final application of \(r_s\) requires this paste also to be a horizontal identity. Thus
        \[
            \id_h^{\mathrm v}
            =
            \id_v^{\mathrm h}
        \]
        for some \(v\in\Sym^+(A)\), which forces the common box to be the double identity \(\Theta_a\) for a unique \(a\in\Conf(A)\).

        Hence the two input labels are both \(a\), the output label is \(a\), and the surviving witness is unique. Therefore
        \[
            \mu_A^C
            =
            \left[
                \Conf(A)\times\Conf(A)
                \xleftarrow{\Delta_{\Conf(A)}}
                \Conf(A)
                \xrightarrow{1_{\Conf(A)}}
                \Conf(A)
            \right].
        \]

        \item \textbf{Calculate the transported unit.} The unit is
        \[
            r_s\eta.
        \]
        Its pullback apex consists of pairs
        \[
            (h,v)
            \in
            \Sym^-_1(A)\times\Sym^+_1(A)
        \]
        satisfying
        \[
            \id_h^{\mathrm v}
            =
            \id_v^{\mathrm h}.
        \]
        Again the common box is uniquely a double identity \(\Theta_a\). Therefore
        \[
            \eta_A^C
            =
            \left[
                1
                \xleftarrow{}
                \Conf(A)
                \xrightarrow{1_{\Conf(A)}}
                \Conf(A)
            \right].
        \]

        \item \textbf{Calculate the transported comultiplication.} The comultiplication is
        \[
            (r_s\otimes r_s)\Delta i_s.
        \]
        An \(i_s\)-witness over \(a\) is a vertical identity \(\id_h^{\mathrm v}\) with \(t(h)=a\). By Lemma~\ref{lem:cgqg-vacant-cell-calculus}, a horizontal
        factorization of \(\id_h^{\mathrm v}\) is uniquely determined by a factorization
        \[
            h=h_2\circ h_1,
        \]
        with factors
        \[
            \id_{h_1}^{\mathrm v},
            \qquad
            \id_{h_2}^{\mathrm v}.
        \]
        For both factors to survive the two applications of \(r_s\), each must also be a horizontal identity. Hence \(h_1,h_2\), and therefore \(h\), are identity arrows.

        The unique surviving factorization is the double identity at \(a\), and its two output labels are \((a,a)\). Thus
        \[
            \delta_A^C
            =
            \left[
                \Conf(A)
                \xleftarrow{1_{\Conf(A)}}
                \Conf(A)
                \xrightarrow{\Delta_{\Conf(A)}}
                \Conf(A)\times\Conf(A)
            \right].
        \]

        \item \textbf{Calculate the transported counit.} Finally, \(\varepsilon i_s\) has the same double-identity pullback as the unit calculation, now with exterior legs \(1_{\Conf(A)}\) and the unique map to \(1\). Hence
        \[
            \epsilon_A^C
            =
            \left[
                \Conf(A)
                \xleftarrow{1_{\Conf(A)}}
                \Conf(A)
                \xrightarrow{}
                1
            \right].
        \]

        \item \textbf{Identify the canonical Frobenius structure.} The four transported spans are exactly the canonical Frobenius structure on \(\Conf(A)\) from Proposition~\ref{prop:cgqg-background-frobenius-base}. Hence the transported Frobenius monoid is \(C_{\Conf(A)}\). \qedhere
    \end{enumerate}
\end{proof}

Write
\[
    C_A:=C_{\Conf(A)}.
\]

With the Pastro--Street convention
of~\eqref{eq:cgqg-ps-an-convention-bridge},
Proposition~\ref{prop:cgqg-incidence-counital-maps} gives
\[
    r_A^{\mathrm{PS}}
    =
    \varepsilon_t
    =
    i_tr_t,
    \qquad
    t_A^{\mathrm{PS}}
    =
    \varepsilon_s
    =
    i_sr_s.
\]
Moreover,
\[
    s_A^{\mathrm{PS}}
    =
    S^{-1}\varepsilon_t.
\]
Since \(S^2=1\) and \(Si_t=i_s\),
\[
    s_A^{\mathrm{PS}}
    =
    S i_tr_t
    =
    i_sr_t.
\]
Thus
\begin{align}
    r_A^{\mathrm{PS}}
    &=
    i_tr_t,
    \\
    t_A^{\mathrm{PS}}
    &=
    i_sr_s,
    \\
    s_A^{\mathrm{PS}}
    &=
    i_sr_t.
\end{align}
After splitting the target retract by \(\Conf(A)\), the base-valued
source and target maps are therefore
\begin{equation}
    \bar s_A=r_t,
    \qquad
    \bar t_A=r_s.
    \label{eq:cgqg-base-valued-source-target}
\end{equation}

\subsection{Boundary coactions and composition}

Define the Pastro--Street coactions
\begin{align}
    \gamma_l
    &:=c_{\WH_{\mathrm{sp}}(A),C_A}(1\otimes\bar s_A)\Delta:
    \WH_{\mathrm{sp}}(A)\longrightarrow C_A\otimes\WH_{\mathrm{sp}}(A),
    \label{eq:cgqg-ps-left-coaction}\\
    \gamma_r
    &:=(1\otimes\bar t_A)\Delta:
    \WH_{\mathrm{sp}}(A)\longrightarrow\WH_{\mathrm{sp}}(A)\otimes C_A.
\label{eq:cgqg-ps-right-coaction}
\end{align}

For a box
\[
    \begin{tikzcd}[column sep=huge,row sep=huge]
         p\arrow[r,"h_t"]\arrow[d,"v_l"']
         &q\arrow[d,"v_r"]\\
         r\arrow[r,"h_b"']&s,
    \end{tikzcd}
\]
put
\begin{equation}
    \ell_A(X):=q,
    \qquad
    r_A(X):=s.
\label{eq:cgqg-box-boundary-functions}
\end{equation}

\begin{proposition}[Boundary graph coactions]
\label{prop:cgqg-ps-coactions-boundary}
    The left coaction is the graph of
    \[
      X\longmapsto(\ell_A(X),X),
    \]
    and the right coaction is the graph of
    \[
      X\longmapsto(X,r_A(X)).
    \]

    Equivalently, the two graph coactions are
    \begin{equation*}
        \begin{tikzcd}[column sep=huge]
            \WH_{\mathrm{sp}}(A) \arrow[r,"{X\mapsto(\ell_A(X),X)}"] & \Conf(A)\times\WH_{\mathrm{sp}}(A)
        \end{tikzcd}
        \qquad
        \begin{tikzcd}[column sep=huge]
            \WH_{\mathrm{sp}}(A) \arrow[r,"{X\mapsto(X,r_A(X))}"] & \WH_{\mathrm{sp}}(A)\times \Conf(A).
        \end{tikzcd}
    \end{equation*}
\end{proposition}
\begin{proof}
    \leavevmode
    \begin{enumerate}[beginpenalty=10000]
        \item \textbf{Compute the left boundary coaction.} In \((1\otimes\bar s_A)\Delta(X)\), a horizontal factorization survives exactly when its right factor lies in the domain of \(r_t\), hence is a horizontal identity. The unique such factorization has right factor \(\id_{v_r}^{\mathrm h}\), and
        \[
            r_t(\id_{v_r}^{\mathrm h})=s(v_r)=q.
        \]
      
        \item \textbf{Compute the right boundary coaction.} The right-coaction calculation is identical with \(r_s\), giving \(t(v_r)=s\). \qedhere
    \end{enumerate}
\end{proof}

\begin{proposition}[Right-mate boundary calculus in spans]
\label{prop:cgqg-right-mate-boundary-presentation}
    Let \(P,Q\) be sets.
    \begin{enumerate}
        \item Every right \(C_P\)-coaction on a set \(X\) in \(\Sp\) is represented by the graph coaction of a unique map
        \[
            r_X:X\longrightarrow P.
        \]

        \item Since \(C_P\) is commutative and
        \[
            C_P\otimes C_Q\cong C_{P\times Q},
        \]
        every right \(C_P^\circ\otimes C_Q\)-coaction on a set \(X\) is represented by the graph coaction of a unique map
        \[
            d_X:X\longrightarrow P\times Q.
        \]

        \item Writing
        \[
            d_X=(\ell_X,r_X),
        \]
        the Chikhladze right-mate convention identifies this comodule with the \(C_P\)--\(C_Q\)-bicomodule having boundary presentation
        \begin{equation}
            \begin{tikzcd}[column sep=huge]
                P
                &
                X
                    \arrow[l,"\ell_X"']
                    \arrow[r,"r_X"]
                &
                Q.
            \end{tikzcd}
            \label{eq:cgqg-right-mate-boundary-diagram}
        \end{equation}
        Conversely, this boundary diagram reconstructs the original right \(C_P^\circ\otimes C_Q\)-coaction on \(X\).
    \end{enumerate}
\end{proposition}
\begin{proof}
    \leavevmode
    \begin{enumerate}[beginpenalty=10000]
        \item \textbf{Normalize a right \(C_P\)-coaction.} Represent a right \(C_P\)-coaction on \(X\) by a span
        \[
            X
            \xleftarrow{\ell}
            R
            \xrightarrow{(u,b)}
            X\times P.
        \]
        Composing with the counit of \(C_P\) gives the span
        \[
            X
            \xleftarrow{\ell}
            R
            \xrightarrow{u}
            X.
        \]
        The comodule counit equation says that this morphism of \(\Sp\) is the identity span on \(X\). Hence there is an apex isomorphism
        \[
            \varphi:R\xrightarrow{\cong}X
        \]
        commuting with both exterior legs, so that
        \[
            \ell=\varphi,
            \qquad
            u=\varphi.
        \]
        Transporting the coaction along \(\varphi\) therefore puts it in graph form
        \[
            X
            \xleftarrow{1_X}
            X
            \xrightarrow{(1_X,r_X)}
            X\times P,
        \]
        where
        \[
            r_X:=b\circ\varphi^{-1}.
        \]
        Coassociativity is then the equality of the two graph spans of
        \[
            x\longmapsto
            \bigl(x,r_X(x),r_X(x)\bigr),
        \]
        so every such map \(r_X:X\to P\) indeed defines a right \(C_P\)-coaction.

        To prove uniqueness, suppose that the graph coactions of
        \[
            r_X,r_X':X\longrightarrow P
        \]
        represent the same morphism
        \[
            X\longrightarrow X\times P
        \]
        in \(\Sp\). An isomorphism between their apexes must commute with the two left legs \(1_X\), and is therefore \(1_X\). Commutation with the right legs then gives
        \[
            (1_X,r_X)=(1_X,r_X'),
        \]
        hence
        \[
            r_X=r_X'.
        \]

        \item \textbf{Apply the graph description to the product base.} The canonical Frobenius monoid \(C_P\) is commutative, so
        \[
            C_P^\circ=C_P.
        \]
        Moreover, taking Cartesian products of the four defining Frobenius spans gives
        \[
            C_P\otimes C_Q
            \cong
            C_{P\times Q}.
        \]
        Hence
        \[
            C_P^\circ\otimes C_Q
            \cong
            C_{P\times Q}.
        \]
        Applying Step~1 with base \(P\times Q\) shows that a right \(C_P^\circ\otimes C_Q\)-coaction on the fixed underlying set \(X\) is represented by the graph of a unique map
        \[
            d_X:X\longrightarrow P\times Q.
        \]

        \item \textbf{Expose the two boundary coordinates.} Write
        \[
            d_X=(\ell_X,r_X).
        \]
        The biduality of the \(P\)-coordinate is represented in spans by the diagonal evaluation and coevaluation
        \[
            P\times P
            \xleftarrow{\Delta_P}
            P
            \longrightarrow
            1,
            \qquad
            1
            \longleftarrow
            P
            \xrightarrow{\Delta_P}
            P\times P.
        \]
        Their snake composites are identity spans: the defining pullbacks impose equality of the repeated \(P\)-labels.

        Taking the mate of the \(P\)-coordinate of the graph coaction therefore leaves the carrier \(X\) unchanged and turns the two coordinates of \(d_X\) into the graph coactions
        \[
            x\longmapsto(\ell_X(x),x),
            \qquad
            x\longmapsto(x,r_X(x)).
        \]
        These are precisely the left \(C_P\)- and right \(C_Q\)-coactions of the boundary diagram
        \[
            P\xleftarrow{\ell_X}X\xrightarrow{r_X}Q.
        \]

        \item \textbf{Invert the mate construction.} Applying the opposite snake identity moves the exposed \(P\)-coordinate back to the right and recovers the graph
        \[
            d_X=(\ell_X,r_X):
            X\longrightarrow P\times Q.
        \]
        Thus, on the fixed underlying set \(X\), the right-comodule and boundary-bicomodule presentations determine one another uniquely. This is the explicit span form of the bidual convention used in Chikhladze's multitensor calculus~\cite[\S3]{Chikhladze2010QuantumModules}. \qedhere
    \end{enumerate}
\end{proof}

\begin{proposition}[Boundary presentations of the arrow object]
\label{prop:cgqg-arrow-right-mate-boundary}
The right \(C_A^\circ\otimes C_A\)-comodule underlying the arrow
object has boundary map
\[
 d_A^{R}(X)=(\ell_A(X),r_A(X))=(q,s)
\]
for a box with vertices \(p,q,r,s\) as in
\eqref{diag:cgqg-game-resource-box}. Equivalently, its right-mate
boundary presentation is
\begin{equation}
 \begin{tikzcd}[column sep=large]
 \Conf(A)&\WH_{\mathrm{sp}}(A)
    \arrow[l,"\ell_A"']\arrow[r,"r_A"]&\Conf(A).
 \end{tikzcd}
 \label{eq:cgqg-arrow-right-mate-boundary}
\end{equation}
The associated endocomonad on \(C_{\Conf(A)\times\Conf(A)}\) retains
both boundary maps
\begin{equation}
 d_A^{L}(X)=(p,r),\qquad d_A^{R}(X)=(q,s).
 \label{eq:cgqg-arrow-endocomonad-boundaries}
\end{equation}
Its comultiplication records horizontal box factorizations, and its
boundary counit is the span
\begin{equation}
 \Boxes_A
 \xleftarrow{v\mapsto\id_v^{\mathrm h}}\Sym^+_1(A)
 \xrightarrow{v\mapsto(s(v),t(v))}
 \Conf(A)\times\Conf(A).
 \label{eq:cgqg-arrow-endocomonad-counit}
\end{equation}
\end{proposition}
\begin{proof}
    \leavevmode
    \begin{enumerate}[beginpenalty=10000]
        \item \textbf{Identify the right-comodule boundary.} Propositions~\ref{prop:cgqg-ps-coactions-boundary} and~\ref{prop:cgqg-right-mate-boundary-presentation} identify the right coaction with \(d_A^{R}(X)=(q,s)\).

        \item \textbf{Compute the boundary counit.} A horizontal factorization with both factors in the support of~\eqref{eq:cgqg-arrow-endocomonad-counit} consists uniquely of two copies of the same horizontal identity. Thus this span preserves the comonoid structure.

        \item \textbf{Recover both endocomonad boundaries.} Applying the boundary counit to the first horizontal factor of \(\Delta(X)\) extracts the left vertical edge; applying it to the second extracts the right vertical edge. The associated endocomonad therefore has boundary diagram
        \[
        \begin{tikzcd}[column sep=8em]
        \Conf(A)^2 & \Boxes_A
          \arrow[l,"{X\mapsto(\boxnw{X},\boxsw{X})}"']
          \arrow[r,"{X\mapsto(\boxne{X},\boxse{X})}"] & \Conf(A)^2.
        \end{tikzcd}
        \]

        \item \textbf{Verify the comonad laws.} The relative comultiplication records horizontal factorizations. Their shared edge supplies the matching boundary, and the displayed counit selects horizontal identity factors. Associativity and the identity laws of the horizontal box groupoid give the comonad equations. \qedhere
    \end{enumerate}
\end{proof}

Put
\begin{equation}
    \mathsf{Comp}_A
    :=
    \left\{
        (X,Y)\in
        \WH_{\mathrm{sp}}(A)\times\WH_{\mathrm{sp}}(A)
        \;\middle|\;
        r_A(X)=\ell_A(Y)
    \right\}.
    \label{eq:cgqg-ps-cotensor-identification}
\end{equation}
Equivalently, \(\mathsf{Comp}_A\) is the boundary pullback
\begin{equation*}
    \begin{tikzcd}[row sep=large,column sep=huge]
        \mathsf{Comp}_A
            \arrow[r]
            \arrow[d]
            \arrow[dr,phantom,"\lrcorner",very near start]
        &
        \WH_{\mathrm{sp}}(A)
            \arrow[d,"\ell_A"]
        \\
        \WH_{\mathrm{sp}}(A)
            \arrow[r,"r_A"']
        &
        \Conf(A).
    \end{tikzcd}
\end{equation*}
Let
\[
    j_A:
    \mathsf{Comp}_A
    \longrightarrow
    \WH_{\mathrm{sp}}(A)\times\WH_{\mathrm{sp}}(A)
\]
be the graph of the inclusion, and let \(p_A\) be its converse span. Put
\[
    a_A:=j_Ap_A.
\]
Thus \(a_A\) is the partial identity on \(\WH_{\mathrm{sp}}(A)\times\WH_{\mathrm{sp}}(A)\) supported on the boundary equation
\[
    r_A(X)=\ell_A(Y),
\]
and
\[
    p_Aj_A=1_{\mathsf{Comp}_A},
    \qquad
    j_Ap_A=a_A.
\]

\begin{lemma}[Local exactness for the game cotensor]
\label{lem:cgqg-local-pastro-street-exactness}
    The object \(\mathsf{Comp}_A\) is an explicit representative of the cotensor
    \begin{equation}
        \WH_{\mathrm{sp}}(A)
        \boxtimes_{C_A}
        \WH_{\mathrm{sp}}(A)
        \cong
        \mathsf{Comp}_A.
        \label{eq:cgqg-game-cotensor-boundary-pullback}
    \end{equation}
    Moreover:
    \begin{enumerate}
        \item the graph inclusion
        \[
            j_A:
            \mathsf{Comp}_A
            \longrightarrow
            \WH_{\mathrm{sp}}(A)\times\WH_{\mathrm{sp}}(A)
        \]
        is a split, hence absolute, equalizer of the two boundary-coaction maps;

        \item the Pastro--Street cotensor idempotent is \(a_A=j_Ap_A\);

        \item the induced boundaries on \(\mathsf{Comp}_A\) are
        \[
            \ell_A^{(2)}(X,Y)=\ell_A(X),
            \qquad
            r_A^{(2)}(X,Y)=r_A(Y);
        \]

        \item every tensoring of this equalizer occurring in the forward Pastro--Street construction is again an equalizer;

        \item for every \(n\geq2\), the iterated cotensor is represented explicitly by
        \begin{equation}
            \mathsf{Comp}_A^{(n)}
            :=
            \left\{
                (X_1,\ldots,X_n)
                \in
                \WH_{\mathrm{sp}}(A)^n
                \;\middle|\;
                r_A(X_i)=\ell_A(X_{i+1})
                \text{ for }1\leq i<n
            \right\}.
            \label{eq:cgqg-local-ps-iterated-comp}
        \end{equation}
        with exterior boundaries
        \[
            \ell_A^{(n)}=\ell_A\pi_1,
            \qquad
            r_A^{(n)}=r_A\pi_n.
        \]
        If
        \[
            j_A^{(n)}:
            \mathsf{Comp}_A^{(n)}
            \longrightarrow
            \WH_{\mathrm{sp}}(A)^n
        \]
        is the graph inclusion and \(p_A^{(n)}\) its converse, then
        \[
            p_A^{(n)}j_A^{(n)}
            =
            1_{\mathsf{Comp}_A^{(n)}},
            \qquad
            j_A^{(n)}p_A^{(n)}
            =
            a_A^{(n)},
        \]
        where \(a_A^{(n)}\) is the partial identity supported on all adjacent boundary equations.
    \end{enumerate}
\end{lemma}
\begin{proof}
    \leavevmode
    \begin{enumerate}[beginpenalty=10000]
        \item \textbf{Write the two cotensor maps explicitly.} By Proposition~\ref{prop:cgqg-ps-coactions-boundary}, the right boundary coaction of the first box records \(r_A(X)\), while the left boundary coaction of the second records \(\ell_A(Y)\). Hence the defining cotensor equalizer compares the graph maps
        \[
            f_A,g_A:
            \WH_{\mathrm{sp}}(A)\times\WH_{\mathrm{sp}}(A)
            \longrightarrow
            \WH_{\mathrm{sp}}(A)\times
            \Conf(A)\times
            \WH_{\mathrm{sp}}(A)
        \]
        given by
        \[
            f_A(X,Y)
            =
            \bigl(X,r_A(X),Y\bigr),
            \qquad
            g_A(X,Y)
            =
            \bigl(X,\ell_A(Y),Y\bigr).
        \]
        They agree precisely on
        \[
            r_A(X)=\ell_A(Y),
        \]
        so their equalizing locus is \(\mathsf{Comp}_A\).

        The Pastro--Street idempotent on the ambient tensor product is
        \begin{equation}
            e_A^{\mathrm{PS}}
            :=
            (1\otimes\epsilon_A^C\otimes1)
            (1\otimes\mu_A^C\otimes1)
            (\gamma_r\otimes\gamma_l).
            \label{eq:cgqg-game-ps-relative-tensor-idempotent}
        \end{equation}
        The coactions record \((X,r_A(X),\ell_A(Y),Y)\); the canonical Frobenius multiplication imposes equality of the two middle labels and the counit forgets their common value. Hence
        \begin{equation}
            e_A^{\mathrm{PS}}=j_Ap_A=a_A.
            \label{eq:cgqg-ps-idempotent-equals-boundary-idempotent}
        \end{equation}

        \item \textbf{Construct the split equalizer.} Let
        \[
            d_A:
            \WH_{\mathrm{sp}}(A)\times
            \Conf(A)\times
            \WH_{\mathrm{sp}}(A)
            \longrightarrow
            \WH_{\mathrm{sp}}(A)\times\WH_{\mathrm{sp}}(A)
        \]
        be the span whose apex is \(\WH_{\mathrm{sp}}(A)\times\WH_{\mathrm{sp}}(A)\), with left leg
        \[
            (X,Y)
            \longmapsto
            \bigl(X,r_A(X),Y\bigr)
        \]
        and right leg
        \[
            (X,Y)\longmapsto(X,Y).
        \]
        Then
        \[
            d_Af_A
            =
            1_{\WH_{\mathrm{sp}}(A)\times\WH_{\mathrm{sp}}(A)},
            \qquad
            d_Ag_A
            =
            a_A,
        \]
        while
        \[
            f_Aa_A=g_Aa_A.
        \]

        \item \textbf{Verify the equalizer universal property.} Let
        \[
            z:
            Z
            \longrightarrow
            \WH_{\mathrm{sp}}(A)\times\WH_{\mathrm{sp}}(A)
        \]
        satisfy
        \[
            f_Az=g_Az.
        \]
        Then
        \[
            a_Az
            =
            d_Ag_Az
            =
            d_Af_Az
            =
            z.
        \]
        Since \(a_A=j_Ap_A\), this gives
        \[
            z=j_Ap_Az.
        \]
        The factorization through \(j_A\) is unique because \(p_Aj_A=1_{\mathsf{Comp}_A}\). Thus \(j_A\) is the required equalizer.

        \item \textbf{Conclude absoluteness.} The equations
        \[
            d_Af_A=1,
            \qquad
            d_Ag_A=a_A,
            \qquad
            f_Aa_A=g_Aa_A
        \]
        exhibit the equalizer as split. Split equalizers are absolute, hence preserved by every functor. In particular, every left or right tensoring of \(j_A\) occurring in the Pastro--Street construction is again an equalizer.

        \item \textbf{Identify the induced boundaries.} Restricting the tensor-product coactions along \(j_A\) leaves the left boundary of the first factor and the right boundary of the second, so
        \[
            \ell_A^{(2)}(X,Y)=\ell_A(X),
            \qquad
            r_A^{(2)}(X,Y)=r_A(Y).
        \]

        \item \textbf{Construct the threefold cotensor.} Applying the same boundary-pullback calculation once more gives
        \[
            \mathsf{Comp}_A^{(3)}
            =
            \left\{
                (X_1,X_2,X_3)
                \;\middle|\;
                r_A(X_1)=\ell_A(X_2),\;
                r_A(X_2)=\ell_A(X_3)
            \right\}.
        \]
        This represents either parenthesization of the threefold cotensor, and the associativity comparison is induced by the identity on this simultaneous boundary pullback.

        \item \textbf{Iterate the construction.} Repeating the same argument gives
        \[
            \mathsf{Comp}_A^{(n)}
            =
            \left\{
                (X_1,\ldots,X_n)
                \;\middle|\;
                r_A(X_i)=\ell_A(X_{i+1})
                \text{ for }1\leq i<n
            \right\}
        \]
        for every \(n\geq2\), with exterior boundaries \(\ell_A\pi_1\) and \(r_A\pi_n\). Its graph inclusion \(j_A^{(n)}\) has converse \(p_A^{(n)}\), so
        \[
            p_A^{(n)}j_A^{(n)}
            =
            1_{\mathsf{Comp}_A^{(n)}},
        \]
        while
        \[
            a_A^{(n)}
            :=
            j_A^{(n)}p_A^{(n)}
        \]
        is the partial identity imposing all adjacent boundary equations. Hence every finite cotensor and every splitting needed below is represented explicitly in \(\Sp\). \qedhere
    \end{enumerate}
\end{proof}

\begin{proposition}[Transport of the Pastro--Street retracts]
\label{prop:cgqg-ps-karoubi-transport}
    Let
    \[
        J:\Sp\longrightarrow\Kar(\Sp),
        \qquad
        X\longmapsto(X,1_X).
    \]
    Then the Pastro--Street target object
    \[
        \bigl(\WH_{\mathrm{sp}}(A),t_A^{\mathrm{PS}}\bigr)
    \]
    is isomorphic to \(J(\Conf(A))\) through \(i_s,r_s\); its Frobenius structure transports to \(J(C_A)\); and its cotensor Karoubi object is the split image of
    \[
        a_A=j_Ap_A,
    \]
    hence isomorphic to \(J(\mathsf{Comp}_A)\).

    The splitting of the composition idempotent is
    \begin{equation*}
        \begin{tikzcd}[column sep=huge]
            \mathsf{Comp}_A
                \arrow[r,bend left=18,"j_A"]
            &
            \WH_{\mathrm{sp}}(A)\times\WH_{\mathrm{sp}}(A)
                \arrow[l,bend left=18,"p_A"]
                \arrow[loop right,"a_A"]
        \end{tikzcd}
    \end{equation*}
    with
    \[
        p_Aj_A=1_{\mathsf{Comp}_A},
        \qquad
        j_Ap_A=a_A.
    \]
\end{proposition}
\begin{proof}
    \leavevmode
    \begin{enumerate}[beginpenalty=10000]
        \item \textbf{Split the target retract.} Proposition~\ref{prop:cgqg-incidence-counital-maps} gives
        \[
            r_si_s=1_{\Conf(A)},
            \qquad
            i_sr_s=t_A^{\mathrm{PS}}.
        \]
        Hence \(i_s,r_s\) exhibit \(\bigl(\WH_{\mathrm{sp}}(A),t_A^{\mathrm{PS}}\bigr)\) as isomorphic to \(J(\Conf(A))\) in \(\Kar(\Sp)\).

        \item \textbf{Transport the Frobenius structure.} Proposition~\ref{prop:cgqg-configuration-frobenius-retract} identifies the Frobenius structure transported through this retract with the canonical structure \(C_A\). Hence the target Frobenius object is identified with \(J(C_A)\).

        \item \textbf{Identify the cotensor idempotent.} The two boundary coactions record
        \[
            r_A(X)
            \qquad\text{and}\qquad
            \ell_A(Y).
        \]
        Therefore the Pastro--Street cotensor idempotent is the partial identity requiring
        \[
            r_A(X)=\ell_A(Y).
        \]
        By Lemma~\ref{lem:cgqg-local-pastro-street-exactness}, this is precisely
        \[
            a_A=j_Ap_A,
        \]
        whose split image is \(\mathsf{Comp}_A\). \qedhere
    \end{enumerate}
\end{proof}

Restrict weak multiplication along the split cotensor inclusion:
\begin{align}
    \mu_A
    &:=
    m\,j_A:
    \mathsf{Comp}_A
    \longrightarrow
    \WH_{\mathrm{sp}}(A),
    \label{eq:cgqg-ps-local-composition}\\
    u_A
    &:=
    i_s:
    C_A
    \longrightarrow
    \WH_{\mathrm{sp}}(A).
    \label{eq:cgqg-ps-local-unit}
\end{align}
The composition span has domain \(\mathsf{Comp}_A\) and support the pairs whose full intermediate horizontal edges agree. On this support its value is vertical box composition.

\begin{proposition}[Quantum-category structure]
\label{prop:cgqg-local-pastro-street}
    The arrow comonoid \(\WH_{\mathrm{sp}}(A)\), its boundary coactions \(\gamma_l,\gamma_r\), the composition \(\mu_A\), and the unit \(u_A\) form the Pastro--Street realization of a Day--Street quantum category in \(\Sp\).
\end{proposition}
\begin{proof}
    \leavevmode
    \begin{enumerate}[beginpenalty=10000]
        \item \textbf{Recall the forward Pastro--Street construction.} The forward direction of \cite[Thm.~6.1 and \S6.1]{PastroStreet2009} constructs the quantum category in the Karoubi completion from the target separable-Frobenius retract and the required relative tensor products. For the present weak Hopf monoid, their defining equalizers and tensor preservation are supplied explicitly by Lemma~\ref{lem:cgqg-local-pastro-street-exactness}.

        \item \textbf{Identify the target base.} Proposition~\ref{prop:cgqg-ps-karoubi-transport} identifies the target object with \(J(\Conf(A))\), while Proposition~\ref{prop:cgqg-configuration-frobenius-retract} identifies its transported separable Frobenius structure with \(J(C_A)\).

        \item \textbf{Identify the relative composition object.} Proposition~\ref{prop:cgqg-ps-coactions-boundary} identifies the relevant coactions with the boundary maps \(r_A\) and \(\ell_A\). Lemma~\ref{lem:cgqg-local-pastro-street-exactness} therefore identifies the cotensor with
        \[
            \mathsf{Comp}_A
            =
            \left\{
                (X,Y)
                \mid
                r_A(X)=\ell_A(Y)
            \right\},
        \]
        and identifies the corresponding Karoubi idempotent as
        \[
            a_A=j_Ap_A.
        \]

        \item \textbf{Supply the required exactness locally.} Lemma~\ref{lem:cgqg-local-pastro-street-exactness} shows that \(j_A\) is a split absolute equalizer and constructs every finite iterated cotensor as the simultaneous boundary pullback
        \[
            \mathsf{Comp}_A^{(n)}
            =
            \left\{
                (X_1,\ldots,X_n)
                \mid
                r_A(X_i)=\ell_A(X_{i+1})
                \text{ for }1\leq i<n
            \right\}.
        \]
        Hence all tensorings of the equalizers used in the forward construction are preserved.

        \item \textbf{Identify composition and unit.} On these split representatives the Pastro--Street composition is the restriction of weak multiplication,
        \[
            \mu_A
            =
            mj_A:
            \mathsf{Comp}_A
            \longrightarrow
            \WH_{\mathrm{sp}}(A),
        \]
        and the unit is
        \[
            u_A
            =
            i_s:
            C_A
            \longrightarrow
            \WH_{\mathrm{sp}}(A).
        \]
        Hence all the quantum-category structure maps have representatives in \(\Sp\).

        \item \textbf{Reflect the quantum-category equations to spans.} The forward Pastro--Street construction gives the required equations in \(\Kar(\Sp)\). Since
        \[
            J:\Sp\longrightarrow\Kar(\Sp)
        \]
        is fully faithful, they reflect to the displayed spans in \(\Sp\). \qedhere
    \end{enumerate}
\end{proof}

\begin{proposition}[Nullary mate of the unit]
\label{prop:cgqg-explicit-nullary-unit-mate}
The nullary multitensor unit corresponding to \(u_A=i_s\) is
\begin{equation}
 \eta_A^T=
 \left[
 \Conf(A)\times\Conf(A)
 \xleftarrow{h\mapsto(s(h),t(h))}\Sym^-_1(A)
 \xrightarrow{h\mapsto\id_h^{\mathrm v}}\Boxes_A
 \right].
 \label{eq:cgqg-explicit-nullary-unit-span}
\end{equation}
It is a map of comonads from \(T_0(C_A)\) to the arrow endocomonad.
\end{proposition}
\begin{proof}
    \leavevmode
    \begin{enumerate}[beginpenalty=10000]
        \item \textbf{Match the boundaries.} At \(h:p\to q\), the vertical identity box has source \((p,p)\) and target \((q,q)\) by~\eqref{eq:cgqg-arrow-endocomonad-boundaries}. These are the boundaries of \((p,q)\in T_0(C_A)\) from Proposition~\ref{prop:cgqg-background-span-multitensors}.

        \item \textbf{Take the nullary mate.} The direct unit has apex \(\Sym^-_1(A)\) and left leg \(t(h)\). The nullary Kan correspondence records both the incoming coevaluation label \(s(h)\) and the outgoing label \(t(h)\), giving
        \[
        \begin{tikzcd}[column sep=large,row sep=large]
        &\Sym^-_1(A)\arrow[dl,"{h\mapsto(s(h),t(h))}"']
          \arrow[dr,"{h\mapsto\id_h^{\mathrm v}}"]&\\
        \Conf(A)^2&&\Boxes_A.
        \end{tikzcd}
        \]

        \item \textbf{Verify comultiplication.} Horizontal factorizations of \(\id_h^{\mathrm v}\) are precisely the vertical identity boxes of factorizations \(h=h_2\circ h_1\). The matrix comultiplication of \(T_0(C_A)\) chooses their intermediate configuration, followed by the two witnesses \(h_1,h_2\). This is a bijection of the two composite apices preserving both exterior legs.

        \item \textbf{Verify the counit.} The counit selects the identity negative arrows. There is one such witness at each pair \((p,p)\), giving the counit of \(T_0(C_A)\). \qedhere
    \end{enumerate}
\end{proof}

\subsection{The quantum groupoid of a thin game}
\label{subsec:cgqg-game-quantum-groupoid}

We have constructed the quantum-category data. It remains to construct the inversion data of Definition~\ref{def:cgqg-background-quantum-groupoid} and verify equations~\textup{(G1)--(G3)}.

\begin{construction}[Transported inversion data]
\label{con:cgqg-transported-inversion-data}
    We transport the Pastro--Street inversion data of \(\WH_{\mathrm{sp}}(A)\) to the explicit configuration base and split composition object constructed above.

    \leavevmode
    \begin{enumerate}
        \item \textbf{Transport the object and arrow inversion maps.} Pastro--Street's weak-Hopf construction gives
        \[
            \nu
            =
            S:
            \WH_{\mathrm{sp}}(A)^\circ
            \longrightarrow
            \WH_{\mathrm{sp}}(A),
            \qquad
            \upsilon
            =
            t_A^{\mathrm{PS}}S^2t_A^{\mathrm{PS}}.
        \]
        Since
        \[
            S^2=1_{\WH_{\mathrm{sp}}(A)},
        \]
        transport through the target splitting \(i_s,r_s\) gives
        \begin{equation}
            \upsilon_A
            =
            1_{C_A},
            \qquad
            \nu_A(X)
            =
            X^{-1}.
            \label{eq:cgqg-vacant-inversion-simplification}
        \end{equation}

        \item \textbf{Identify the Pastro--Street composition comodules.} The remaining inversion datum is a left \(\WH_{\mathrm{sp}}(A)^{\otimes3}\)-comodule isomorphism
        \begin{equation}
            \theta_{\mathrm{PS}}:
            P_{A,l}^{\mathrm{PS}}
            \xrightarrow{\cong}
            P_{A,r}^{\mathrm{PS}},
            \label{eq:cgqg-ps-theta-typed}
        \end{equation}
        where \(P_{A,l}^{\mathrm{PS}}\) and \(P_{A,r}^{\mathrm{PS}}\) are the two Pastro--Street comodule structures on the common composition object; see~\cite[\S5.3, Thm.~6.1 and \S6.1]{PastroStreet2009}.

        By Proposition~\ref{prop:cgqg-ps-karoubi-transport}, their common underlying Karoubi object is
        \[
            \bigl(
                \WH_{\mathrm{sp}}(A)\times\WH_{\mathrm{sp}}(A),
                a_A
            \bigr),
            \qquad
            a_A=j_Ap_A,
        \]
        with split image \(J(\mathsf{Comp}_A)\). The splitting maps
        \[
            j_A:
            J(\mathsf{Comp}_A)
            \longrightarrow
            \bigl(
                \WH_{\mathrm{sp}}(A)\times\WH_{\mathrm{sp}}(A),
                a_A
            \bigr),
        \]
        and
        \[
            p_A:
            \bigl(
                \WH_{\mathrm{sp}}(A)\times\WH_{\mathrm{sp}}(A),
                a_A
            \bigr)
            \longrightarrow
            J(\mathsf{Comp}_A)
        \]
        satisfy
        \[
            p_Aj_A
            =
            1_{\mathsf{Comp}_A},
            \qquad
            j_Ap_A
            =
            a_A.
        \]
        Hence \(j_A\) and \(p_A\) are mutually inverse morphisms in
        \(\Kar(\Sp)\).

        \item \textbf{Transport the two comodule structures.} Transport the comodule structure of \(P_{A,l}^{\mathrm{PS}}\) along \(j_A,p_A\), and denote the resulting \(\WH_{\mathrm{sp}}(A)^{\otimes3}\)-comodule on \(\mathsf{Comp}_A\) by
        \[
            \mathsf{Comp}_{A,l}.
        \]
        Likewise, transport the comodule structure of \(P_{A,r}^{\mathrm{PS}}\) and denote the result by
        \[
            \mathsf{Comp}_{A,r}.
        \]

        With these typings, the common underlying span \(j_A\) gives comodule isomorphisms
        \begin{align}
            \kappa_l
            &:=
            j_A:
            \mathsf{Comp}_{A,l}
            \xrightarrow{\cong}
            P_{A,l}^{\mathrm{PS}},
            \label{eq:cgqg-kappa-left-explicit}\\
            \kappa_r
            &:=
            j_A:
            \mathsf{Comp}_{A,r}
            \xrightarrow{\cong}
            P_{A,r}^{\mathrm{PS}},
            \label{eq:cgqg-kappa-right-explicit}
        \end{align}
        whose inverses are both represented by \(p_A\), with the corresponding comodule typings.

        \item \textbf{Transport the inversion isomorphism.} Define
        \begin{equation}
            \theta_A
            :=
            p_A\,
            \theta_{\mathrm{PS}}\,
            j_A:
            \mathsf{Comp}_{A,l}
            \xrightarrow{\cong}
            \mathsf{Comp}_{A,r}.
            \label{eq:cgqg-transported-theta}
        \end{equation}
        In the symmetric category \(\Sp\), the defining
        Pastro--Street composite becomes
        \begin{equation}
          \theta_A
          =p_A(1\otimes S)(1\otimes m)(c\otimes1)
             (1\otimes\Delta)j_A.
          \label{eq:cgqg-explicit-theta-composite}
        \end{equation}
        Here \(c\) exchanges the first box and the left factor of the
        second box; multiplication pastes the remaining two boxes
        vertically, and \(S\) inverts their paste. This is the symmetric
        form of the inversion composite in
        \cite[\S6, p.~182]{PastroStreet2009}.
    \end{enumerate}
\end{construction}

\begin{proposition}[Box formula for the transported inversion]
\label{prop:cgqg-explicit-theta-boxes}
For \((X,Y)\in\mathsf{Comp}_A\), let \(k\) be the bottom edge of
\(X\), and write \(Y=h\trcorner v\). There is a unique horizontal
factorization \(Y=Y_R\circ_{\mathrm h}Y_L\) whose right factor has
top edge \(k\). With \(Z=Y_R\circ_{\mathrm v}X\), the map
\(\theta_A\) is the graph of the bijection
\begin{equation}
 \theta_A(X,Y)=(Y_L,Z^{-1}).
 \label{eq:cgqg-explicit-theta-box-formula}
\end{equation}
The two factors are explicitly
\[
 Y_R=k\trcorner v,
 \qquad Y_L=(k^{-1}\circ h)\trcorner\leftedge{k}{v}.
\]
\end{proposition}
\begin{proof}
    \leavevmode
    \begin{enumerate}[beginpenalty=10000]
        \item \textbf{Construct the required factorization.}
        Boundary matching gives \(t(k)=t(h)\), so
        \(h=k\circ(k^{-1}\circ h)\). The horizontal factorization
        clause of Lemma~\ref{lem:cgqg-vacant-cell-calculus} gives
        exactly the displayed \(Y_L,Y_R\). Since the top edge of
        \(Y_R\) is the bottom edge of \(X\), their vertical paste
        \(Z\) is defined.

        \item \textbf{Read the structural composite.}
        A witness of~\eqref{eq:cgqg-explicit-theta-composite} first
        factorizes \(Y\) horizontally and then requires its right
        factor to paste below \(X\). This forces that factor's top
        edge to be \(k\), so the witness just constructed is unique.
        The output is \((Y_L,Z^{-1})\). Its boundaries match because
        \(r_A(Y_L)=\boxsw{Y_R}=\boxsw{Z}=\ell_A(Z^{-1})\).

        \item \textbf{Construct an inverse.}
        Given \((U,V)\in\mathsf{Comp}_A\), put \(Z=V^{-1}\).
        The right edge of \(U\) and the bottom edge of \(Z\) form
        a compatible southwest corner: their common endpoint is
        \(\boxse{U}=\boxne{V}=\boxsw{Z}\). Vacancy gives the unique box
        \(Y_R\) with that left edge and that bottom edge.
        Vertical groupoid composition determines the unique \(X\)
        with \(Z=Y_R\circ_{\mathrm v}X\); put
        \(Y=Y_R\circ_{\mathrm h}U\). The pair \((X,Y)\) is in
        \(\mathsf{Comp}_A\), since both relevant boundary labels are
        \(\boxne{Y_R}\).

        \item \textbf{Verify the inverse equations.}
        Starting from \((U,V)\), the right factor selected in Step~1
        is the same \(Y_R\), by its top and right edges. Thus the
        displayed formula returns \((U,V)\). Starting from
        \((X,Y)\), the southwest corner in Step~3 reconstructs the
        original \(Y_R\); cancellation then reconstructs \(X\),
        and horizontal pasting reconstructs \(Y\). Hence the two
        constructions are inverse. The equality with the structural
        composite identifies this bijection with the transported
        Pastro--Street inversion. \qedhere
    \end{enumerate}
\end{proof}

\begin{proposition}[Transported inversion datum in spans]
\label{prop:cgqg-transported-inversion-in-spans}
    The transported comodules \(\mathsf{Comp}_{A,l}\) and \(\mathsf{Comp}_{A,r}\) are represented by comodule structures in \(\Sp\), and
    \[
        \theta_A
        =
        p_A\theta_{\mathrm{PS}}j_A:
        \mathsf{Comp}_{A,l}
        \xrightarrow{\cong}
        \mathsf{Comp}_{A,r}
    \]
    is represented by an isomorphism of those \(\Sp\)-comodules.

    The fixed split composition object
    \[
        \mathsf{Comp}_A
        \xrightarrow{\,j_A\,}
        \WH_{\mathrm{sp}}(A)\times\WH_{\mathrm{sp}}(A)
        \xrightarrow{\,p_A\,}
        \mathsf{Comp}_A
    \]
    and its splitting maps uniquely determine the transport of the Pastro--Street inversion map \(\theta_{\mathrm{PS}}\), namely
    \[
        \theta_A
        =
        p_A\theta_{\mathrm{PS}}j_A.
    \]

    Put
    \begin{equation}
        \varsigma_A
        :=
        (\bar s_A\otimes1_{C_A}\otimes\bar t_A)
        (\gamma_r\otimes1_{\WH_{\mathrm{sp}}(A)})
        j_A.
        \label{eq:cgqg-transported-varsigma}
    \end{equation}
    Equivalently,
    \begin{equation}
        \varsigma_A
        =
        (\bar s_A\otimes1_{C_A}\otimes\bar t_A)
        (1_{\WH_{\mathrm{sp}}(A)}\otimes\gamma_l)
        j_A,
        \label{eq:cgqg-transported-varsigma-second-form}
    \end{equation}
    since \(j_A\) equalizes the two boundary coactions.

    Together with
    \[
        \upsilon_A=1_{C_A},
        \qquad
        \nu_A=S_A,
    \]
    these data satisfy \textup{(G1)--(G3)}. In particular,
    \begin{equation}
        (1_{C_A}\otimes1_{C_A}\otimes\upsilon_A)
        c_{C_A,C_A\otimes C_A}
        \varsigma_A
        =
        \varsigma_A\theta_A.
        \label{eq:cgqg-transported-G3}
    \end{equation}
\end{proposition}
\begin{proof}
    \leavevmode
    \begin{enumerate}[beginpenalty=10000]
        \item \textbf{Identify the splitting in the Karoubi completion.} The relations
        \[
            p_Aj_A=1_{\mathsf{Comp}_A},
            \qquad
            j_Ap_A=a_A
        \]
        say precisely that \(j_A\) and \(p_A\) are inverse morphisms between \(J(\mathsf{Comp}_A)\) and the Pastro--Street composition idempotent in \(\Kar(\Sp)\).

        \item \textbf{Transport the two comodule structures.} Transport of a comodule structure along an isomorphism is uniquely determined by the fixed inverse pair \(j_A,p_A\). Hence the two Pastro--Street comodule structures transport uniquely to
        \[
            \mathsf{Comp}_{A,l},
            \qquad
            \mathsf{Comp}_{A,r},
        \]
        and \(j_A\), with its two respective typings, becomes the comodule isomorphisms
        \[
            \kappa_l:
            \mathsf{Comp}_{A,l}
            \xrightarrow{\cong}
            P_{A,l}^{\mathrm{PS}},
            \qquad
            \kappa_r:
            \mathsf{Comp}_{A,r}
            \xrightarrow{\cong}
            P_{A,r}^{\mathrm{PS}}.
        \]

        \item \textbf{Transport the Pastro--Street inversion map.} The Pastro--Street morphism
        \[
            \theta_{\mathrm{PS}}:
            P_{A,l}^{\mathrm{PS}}
            \xrightarrow{\cong}
            P_{A,r}^{\mathrm{PS}}
        \]
        is a comodule isomorphism. Therefore its conjugate
        \[
            \theta_A
            =
            p_A\theta_{\mathrm{PS}}j_A
        \]
        is a comodule morphism
        \[
            \mathsf{Comp}_{A,l}
            \longrightarrow
            \mathsf{Comp}_{A,r}.
        \]
        Moreover,
        \begin{align}
            j_A\theta_A
            &=
            j_Ap_A\theta_{\mathrm{PS}}j_A
            \notag\\
            &=
            a_A\theta_{\mathrm{PS}}j_A
            \notag\\
            &=
            \theta_{\mathrm{PS}}j_A,
            \label{eq:cgqg-theta-transport-intertwining}
        \end{align}
        because \(\theta_{\mathrm{PS}}\) is a morphism between the corresponding Karoubi objects and hence is fixed by the relevant composition idempotents.

        \item \textbf{Construct the inverse of the transported map.} Define
        \[
            \theta_A^{-1}
            :=
            p_A\theta_{\mathrm{PS}}^{-1}j_A.
        \]
        Using \(j_Ap_A=a_A\), together with the fact that \(\theta_{\mathrm{PS}}\) is a morphism of the corresponding Karoubi objects, one obtains
        \begin{align*}
            \theta_A^{-1}\theta_A
            &=
            \bigl(
                p_A\theta_{\mathrm{PS}}^{-1}j_A
            \bigr)
            \bigl(
                p_A\theta_{\mathrm{PS}}j_A
            \bigr)
            \\
            &=
            p_A\theta_{\mathrm{PS}}^{-1}
            a_A
            \theta_{\mathrm{PS}}j_A
            \\
            &=
            p_Aj_A
            =
            1_{\mathsf{Comp}_A}.
        \end{align*}
        The same calculation in the opposite order gives
        \[
            \theta_A\theta_A^{-1}
            =
            1_{\mathsf{Comp}_A}.
        \]
        Hence \(\theta_A\) is a comodule isomorphism.

        \item \textbf{Reflect the transported data to spans.} The source and target of every transported coaction and of \(\theta_A\) lie in the image of the fully faithful embedding
        \[
            J:\Sp\longrightarrow\Kar(\Sp).
        \]
        Therefore each transported morphism is represented uniquely by a morphism in \(\Sp\).

        \item \textbf{Identify the transported \(\varsigma\)-map.} Let
        \[
            \varsigma_{\mathrm{PS}}
        \]
        denote the map appearing in the Pastro--Street \textup{(G3)} equation before splitting the target and composition idempotents. By definition, \(\varsigma_{\mathrm{PS}}\) is obtained from the source and target maps and either of the two boundary coactions on the composition object.

        Transport the target object through
        \[
            C_A
            \xrightarrow{i_s}
            \bigl(
                \WH_{\mathrm{sp}}(A),
                t_A^{\mathrm{PS}}
            \bigr)
            \xrightarrow{r_s}
            C_A
        \]
        and the composition object through \(j_A,p_A\). Since the transported base-valued source and target maps are
        \[
            \bar s_A=r_t,
            \qquad
            \bar t_A=r_s,
        \]
        and the transported boundary coactions are \(\gamma_l,\gamma_r\), the conjugate of \(\varsigma_{\mathrm{PS}}\) is precisely
        \[
            \varsigma_A
            =
            (\bar s_A\otimes1_{C_A}\otimes\bar t_A)
            (\gamma_r\otimes1)
            j_A.
        \]
        Because \(j_A\) is the cotensor equalizer,
        \[
            (\gamma_r\otimes1)j_A
            =
            (1\otimes\gamma_l)j_A,
        \]
        which also gives
        \[
            \varsigma_A
            =
            (\bar s_A\otimes1_{C_A}\otimes\bar t_A)
            (1\otimes\gamma_l)
            j_A.
        \]

        Equivalently, in \(\Kar(\Sp)\),
        \begin{equation}
            \varsigma_{\mathrm{PS}}j_A
            =
            (i_s\otimes i_s\otimes i_s)\varsigma_A.
            \label{eq:cgqg-varsigma-transport-intertwining}
        \end{equation}
        Indeed, the left-hand side lands in the tensor cube of the target retract, while
        \[
            i_sr_s=t_A^{\mathrm{PS}}.
        \]

        \item \textbf{Verify \textup{(G1)--(G2)}.} By Proposition~\ref{prop:cgqg-incidence-counital-maps},
        \[
            \bar s_A\nu_A
            =
            r_tS_A
            =
            r_s
            =
            \bar t_A,
        \]
        giving \textup{(G1)}. Likewise,
        \[
            \bar t_A\nu_A
            =
            r_sS_A
            =
            r_t
            =
            \upsilon_A\bar s_A,
        \]
        since
        \[
            \upsilon_A=1_{C_A}.
        \]
        This gives \textup{(G2)}.

        \item \textbf{Verify \textup{(G3)} explicitly.} Put
        \[
            I_A:=i_s\otimes i_s\otimes i_s.
        \]
        Since
        \[
            r_si_s=1_{C_A},
        \]
        the morphism \(I_A\) is split monic.

        By the transport of \(\upsilon_{\mathrm{PS}}\) to \(\upsilon_A\), naturality of the braiding, and~\eqref{eq:cgqg-varsigma-transport-intertwining},
        \begin{align*}
            &
            I_A
            (1_{C_A}\otimes1_{C_A}\otimes\upsilon_A)
            c_{C_A,C_A\otimes C_A}
            \varsigma_A
            \\
            &\qquad=
            (1\otimes1\otimes\upsilon_{\mathrm{PS}})
            c\,
            \varsigma_{\mathrm{PS}}j_A.
        \end{align*}
        The Pastro--Street \textup{(G3)} equation gives
        \[
            (1\otimes1\otimes\upsilon_{\mathrm{PS}})
            c\,
            \varsigma_{\mathrm{PS}}j_A
            =
            \varsigma_{\mathrm{PS}}
            \theta_{\mathrm{PS}}j_A.
        \]
        Using~\eqref{eq:cgqg-theta-transport-intertwining} and then~\eqref{eq:cgqg-varsigma-transport-intertwining},
        \begin{align*}
            \varsigma_{\mathrm{PS}}
            \theta_{\mathrm{PS}}j_A
            &=
            \varsigma_{\mathrm{PS}}j_A\theta_A
            \\
            &=
            I_A\varsigma_A\theta_A.
        \end{align*}
        Hence
        \[
            I_A
            (1_{C_A}\otimes1_{C_A}\otimes\upsilon_A)
            c_{C_A,C_A\otimes C_A}
            \varsigma_A
            =
            I_A\varsigma_A\theta_A.
        \]
        Since \(I_A\) is split monic, cancellation gives
        \[
            (1_{C_A}\otimes1_{C_A}\otimes\upsilon_A)
            c_{C_A,C_A\otimes C_A}
            \varsigma_A
            =
            \varsigma_A\theta_A.
        \]
        This is precisely \textup{(G3)}.

        \item \textbf{Conclude the construction in \(\Sp\).} Thus the transported comodule structures
        \[
            \mathsf{Comp}_{A,l},
            \qquad
            \mathsf{Comp}_{A,r},
        \]
        the isomorphism
        \[
            \theta_A:
            \mathsf{Comp}_{A,l}
            \xrightarrow{\cong}
            \mathsf{Comp}_{A,r},
        \]
        and the map \(\varsigma_A\) are represented in \(\Sp\). Together with
        \[
            \upsilon_A=1_{C_A},
            \qquad
            \nu_A=S_A,
        \]
        they satisfy \textup{(G1)--(G3)}. \qedhere
    \end{enumerate}
\end{proof}

\begin{definition}[Quantum groupoid associated to a thin game]
\label{def:cgqg-game-quantum-groupoid}
    Let \(\QG_{\mathrm{sp}}(A)\) be the Day--Street quantum groupoid with object-of-objects \(C_A\), arrow object \(\WH_{\mathrm{sp}}(A)\), boundary coactions \(\gamma_l,\gamma_r\), composition object \(\mathsf{Comp}_A\), composition \(\mu_A\), unit \(u_A\), and inversion data
    \[
        \upsilon_A: C_A^{\circ\circ}\to C_A,
        \qquad
        \nu_A:\WH_{\mathrm{sp}}(A)^\circ\to\WH_{\mathrm{sp}}(A),
        \qquad
        \theta_A:\mathsf{Comp}_{A,l}\xrightarrow{\cong}\mathsf{Comp}_{A,r}.
    \]
\end{definition}

\begin{theorem}[Canonical quantum groupoid]
\label{thm:cgqg-game-qg}
    Every thin concurrent game canonically determines \(\QG_{\mathrm{sp}}(A)\). Its object-of-objects is the configuration set with its canonical separable commutative Frobenius structure, and its arrow inversion is total box inversion.
\end{theorem}
\begin{proof}
    \leavevmode
    \begin{enumerate}[beginpenalty=10000]
        \item \textbf{Pass from symmetry to boxes.} Polarized decomposition gives the vacant double groupoid.
  
        \item \textbf{Construct the weak Hopf monoid.} Theorem~\ref{thm:cgqg-span-weak-hopf} gives the incidence weak Hopf monoid.
  
        \item \textbf{Identify the Pastro--Street base and composition object.} Propositions~\ref{prop:cgqg-configuration-frobenius-retract} and~\ref{prop:cgqg-ps-karoubi-transport} identify the target base and cotensor object, and Proposition~\ref{prop:cgqg-local-pastro-street} supplies the quantum category.
  
        \item \textbf{Transport inversion.} Equation~\eqref{eq:cgqg-vacant-inversion-simplification} gives the object and arrow inversion maps, and Proposition~\ref{prop:cgqg-transported-inversion-in-spans} supplies \(\theta_A\) and verifies \textup{(G1)--(G3)} in \(\Sp\). \qedhere
    \end{enumerate}
\end{proof}

\subsection{Tensor and negation}

Here tensor product of quantum groupoids is taken componentwise, and \(\cong\) denotes an isomorphism of the complete quantum-groupoid data.

\begin{theorem}[Monoidal compatibility]
\label{thm:cgqg-qg-tensor}
    For thin games \(A,B\),
    \begin{align}
        \WH_{\mathrm{sp}}(A\otimes B)
        &\cong
        \WH_{\mathrm{sp}}(A)\otimes\WH_{\mathrm{sp}}(B),
        \label{eq:cgqg-span-wh-tensor}\\
        \QG_{\mathrm{sp}}(A\otimes B)
        &\cong
        \QG_{\mathrm{sp}}(A)\otimes\QG_{\mathrm{sp}}(B).
        \label{eq:cgqg-qg-tensor}
    \end{align}
\end{theorem}
\begin{proof}
    \leavevmode
    \begin{enumerate}[beginpenalty=10000]
        \item \textbf{Factor boxes componentwise.} By Proposition~\ref{prop:cgqg-resource-tensor},
        \[
            T(A\otimes B)
            \cong
            T(A)\times T(B).
        \]
        Hence a polarized box of \(A\otimes B\) is canonically a pair consisting of one polarized box of \(A\) and one polarized box of \(B\).

        \item \textbf{Factor the incidence witness objects.} Vertical composability, horizontal factorization, vertical and horizontal identity boxes, and total box inversion are all computed componentwise under this identification. Therefore the witness objects defining
        \[
            m_{A\otimes B},
            \qquad
            \eta_{A\otimes B},
            \qquad
            \Delta_{A\otimes B},
            \qquad
            \varepsilon_{A\otimes B},
            \qquad
            S_{A\otimes B}
        \]
        are Cartesian products of the corresponding witness objects for \(A\) and \(B\).

        \item \textbf{Identify the weak Hopf structure.} Since the tensor product in \(\Sp\) is Cartesian product, the preceding componentwise identifications give
        \[
            \WH_{\mathrm{sp}}(A\otimes B)
            \cong
            \WH_{\mathrm{sp}}(A)\otimes\WH_{\mathrm{sp}}(B)
        \]
        as weak Hopf monoids. This proves~\eqref{eq:cgqg-span-wh-tensor}.

        \item \textbf{Factor the configuration base.} By~\eqref{eq:cgqg-conf-tensor},
        \[
            \Conf(A\otimes B)
            \cong
            \Conf(A)\times\Conf(B),
        \]
        and the canonical Frobenius structure on this product is the tensor product of the canonical Frobenius structures on \(\Conf(A)\) and \(\Conf(B)\). Hence
        \[
            C_{A\otimes B}
            \cong
            C_A\otimes C_B.
        \]

        \item \textbf{Factor the boundary and composition data.} The boundary maps, boundary coactions, split composition object, composition, and unit are all obtained componentwise. In particular,
        \[
            \mathsf{Comp}_{A\otimes B}
            \cong
            \mathsf{Comp}_A\times\mathsf{Comp}_B,
        \]
        since the boundary-matching equation for \(A\otimes B\) is precisely the pair of boundary-matching equations for \(A\) and \(B\).

        \item \textbf{Factor the inversion data.} The target splitting, arrow inversion, transported composition comodules, and transported Pastro--Street map \(\theta\) are likewise products of the corresponding structures for \(A\) and \(B\).

        \item \textbf{Conclude monoidal compatibility.} Thus every datum of the game quantum groupoid is computed componentwise under tensor product, giving
        \[
            \QG_{\mathrm{sp}}(A\otimes B)
            \cong
            \QG_{\mathrm{sp}}(A)\otimes
            \QG_{\mathrm{sp}}(B).
        \]
        This proves \eqref{eq:cgqg-qg-tensor}. \qedhere
    \end{enumerate}
\end{proof}

\begin{proposition}[Compact self-duality of spans]
\label{prop:cgqg-span-compact-closed}
    Every set \(X\) is canonically self-dual in \(\Sp\), with evaluation and coevaluation
    \begin{equation*}
        \begin{tikzcd}
            & X
                \arrow[dl,"\Delta_X"']
                \arrow[dr]
            & \\
            X\times X
            &&
            1
        \end{tikzcd}
        \qquad
        \begin{tikzcd}
            & X
                \arrow[dl]
                \arrow[dr,"\Delta_X"]
            & \\
            1
            &&
            X\times X.
        \end{tikzcd}
    \end{equation*}
    Hence \(\Sp\) is compact closed.
\end{proposition}
\begin{proof}
    \leavevmode
    \begin{enumerate}[beginpenalty=10000]
        \item \textbf{Compute the first snake composite.} Compose the coevaluation and evaluation spans in the first snake identity. The defining pullbacks impose equality of the repeated \(X\)-coordinates, so the resulting apex is canonically \(X\).

        \item \textbf{Identify the first exterior legs.} Under this canonical apex identification, both exterior legs are \(1_X\). Hence the first snake composite is the identity span on \(X\).

        \item \textbf{Compute the second snake composite.} The reflected pullback calculation again has apex canonically \(X\), with both exterior legs equal to \(1_X\). Hence the second snake composite is also the identity span.

        \item \textbf{Conclude compact self-duality.} The displayed evaluation and coevaluation therefore exhibit \(X\) as its own left and right dual. Since this construction is available for every set, \(\Sp\) is compact closed. \qedhere
    \end{enumerate}
\end{proof}

If
\[
    K=(K,m_K,\eta_K,\Delta_K,\varepsilon_K,S_K)
\]
is a weak Hopf monoid in \(\Sp\), write \(K^\vee\) for the compactly transposed weak Hopf structure
\begin{align*}
    m_{K^\vee}&:=\Delta_K^\vee,
    &
    \eta_{K^\vee}&:=\varepsilon_K^\vee,
    \\
    \Delta_{K^\vee}&:=m_K^\vee,
    &
    \varepsilon_{K^\vee}&:=\eta_K^\vee,
    \\
    S_{K^\vee}&:=S_K^\vee.
\end{align*}
If \(Q\) is the Pastro--Street quantum groupoid obtained from a weak
Hopf monoid \(K\) in \(\Sp\), write
\[
    Q^\vee
\]
for the Pastro--Street quantum groupoid obtained from the compactly
transposed weak Hopf monoid \(K^\vee\) in \(\Kar(\Sp)\).

Thus the target object, the two boundary coactions, the composition
idempotent, and the inversion data of \(Q^\vee\) are those determined by
the weak Hopf structure of \(K^\vee\) itself. When these target and
composition idempotents are equipped with explicit splittings in
\(\Sp\), we use those splittings to present \(Q^\vee\) in \(\Sp\).

\begin{lemma}[Weak Hopf structure under compact transposition]
\label{lem:cgqg-compact-transpose-weak-hopf}
    Let
    \[
        K=(K,m,\eta,\Delta,\varepsilon,S)
    \]
    be a weak Hopf monoid in \(\Sp\). Under the canonical compact self-duality of \(\Sp\), the transposed structure
    \[
        K^\vee
        :=
        \bigl(
            K,
            \Delta^\vee,
            \varepsilon^\vee,
            m^\vee,
            \eta^\vee,
            S^\vee
        \bigr)
    \]
    is again a weak Hopf monoid.

    If \(S^2=1_K\), then
    \[
        (S^\vee)^2=1_K.
    \]
    Moreover,
    \[
        (K^\vee)^\vee\cong K
    \]
    canonically as weak Hopf monoids.
\end{lemma}
\begin{proof}
    \leavevmode
    \begin{enumerate}[beginpenalty=10000]
        \item \textbf{Fix the tensor-order convention for transposition.}
        Under the canonical compact self-duality of
        Proposition~\ref{prop:cgqg-span-compact-closed}, every object
        \(X\) is identified with its dual. For a span morphism
        \[
            f:X\longrightarrow Y,
        \]
        write
        \[
            f^\vee:Y\longrightarrow X
        \]
        for its compact transpose.

        Compact transposition is contravariant:
        \[
            (g\circ f)^\vee
            =
            f^\vee\circ g^\vee,
            \qquad
            1_X^\vee=1_X.
        \]
        The canonical compact dual of a tensor product reverses tensor
        order. Throughout this lemma we compose that canonical
        identification with the symmetry of \(\Sp\), and therefore use
        the coherent identification
        \[
            (X\otimes Y)^\vee
            \cong
            X^\vee\otimes Y^\vee.
        \]
        With this convention,
        \[
            (f\otimes g)^\vee
            \cong
            f^\vee\otimes g^\vee,
        \]
        and the symmetry is fixed by transposition.

        \item \textbf{Transpose the monoid and comonoid structures.}
        Since
        \[
            \Delta:
            K\longrightarrow K\otimes K,
        \]
        its transpose has type
        \[
            \Delta^\vee:
            K\otimes K\longrightarrow K
        \]
        under the preceding tensor-order identification. Likewise,
        \[
            \varepsilon^\vee:I\longrightarrow K.
        \]
        Thus
        \[
            m_{K^\vee}:=\Delta^\vee,
            \qquad
            \eta_{K^\vee}:=\varepsilon^\vee
        \]
        have the types of multiplication and unit.

        Similarly,
        \[
            m^\vee:
            K\longrightarrow K\otimes K,
            \qquad
            \eta^\vee:
            K\longrightarrow I,
        \]
        so
        \[
            \Delta_{K^\vee}:=m^\vee,
            \qquad
            \varepsilon_{K^\vee}:=\eta^\vee
        \]
        have the types of comultiplication and counit.

        \item \textbf{Transpose associativity, coassociativity, unit,
        and counit.}
        Transposing the coassociativity equation for \(\Delta\) gives
        the associativity equation for \(\Delta^\vee\). Transposing the
        counit equations for \(\varepsilon\) gives the two unit
        equations for \(\varepsilon^\vee\).

        Conversely, transposing associativity of \(m\) gives
        coassociativity of \(m^\vee\), and transposing the unit equations
        for \(\eta\) gives the counit equations for \(\eta^\vee\).
        Hence the displayed maps define a monoid and a comonoid on
        \(K^\vee\).

        \item \textbf{Transpose multiplicativity.}
        The multiplicativity equation
        \[
            \Delta m
            =
            (m\otimes m)
            (1\otimes c\otimes1)
            (\Delta\otimes\Delta)
        \]
        transposes to
        \[
            m^\vee\Delta^\vee
            =
            (\Delta^\vee\otimes\Delta^\vee)
            (1\otimes c\otimes1)
            (m^\vee\otimes m^\vee),
        \]
        after the coherent tensor-order identifications of Step~1.
        This is precisely the multiplicativity equation for the
        multiplication \(\Delta^\vee\) and comultiplication \(m^\vee\)
        of \(K^\vee\).

        \item \textbf{Transpose the weak unit and weak counit equations.}
        Compact transposition exchanges the weak-unit and weak-counit equations \eqref{eq:cgqg-game-weak-unit-equations}--\eqref{eq:cgqg-game-weak-counit-equations}. The symmetry occurring in those equations is preserved by transposition under the convention of Step~1. Hence
        \[
            \bigl(
                K,
                \Delta^\vee,
                \varepsilon^\vee,
                m^\vee,
                \eta^\vee
            \bigr)
        \]
        is a weak bimonoid.

        \item \textbf{Transpose the counital antipode equations.}
        Transposing the defining composites of the canonical counital maps gives
        \[
            \varepsilon_t^{K^\vee}=(\varepsilon_t^K)^\vee,
            \qquad
            \varepsilon_s^{K^\vee}=(\varepsilon_s^K)^\vee.
        \]
        Hence the two equations
        \[
            m(1\otimes S)\Delta=\varepsilon_t,
            \qquad
            m(S\otimes1)\Delta=\varepsilon_s
        \]
        transpose to the corresponding counital convolution equations
        for \(S^\vee\) in the transposed weak bimonoid. Thus the first
        two antipode equations hold for \(S^\vee\).

        \item \textbf{Transpose the third antipode equation.}
        Transposing
        \[
            m^{(3)}
            (S\otimes1\otimes S)
            \Delta^{(3)}
            =
            S
        \]
        gives
        \[
            m_{K^\vee}^{(3)}
            (S^\vee\otimes1\otimes S^\vee)
            \Delta_{K^\vee}^{(3)}
            =
            S^\vee.
        \]
        Hence \(S^\vee\) is an antipode and \(K^\vee\) is a weak Hopf
        monoid.

        \item \textbf{Transpose involutivity.}
        If
        \[
            S^2=1_K,
        \]
        then
        \[
            1_K
            =
            (S^2)^\vee
            =
            S^\vee S^\vee,
        \]
        so
        \[
            (S^\vee)^2=1_K.
        \]

        \item \textbf{Transpose twice.}
        The two snake identities identify the double compact transpose
        of every span with the original span. These identifications are
        compatible with tensor product and symmetry, so they identify
        every structure map of \((K^\vee)^\vee\) with the corresponding
        structure map of \(K\). Therefore
        \[
            (K^\vee)^\vee
            \cong
            K
        \]
        canonically as weak Hopf monoids. \qedhere
    \end{enumerate}
\end{proof}

\begin{theorem}[Negation--duality]
\label{thm:cgqg-qg-negation}
    For every thin game \(A\),
    \begin{align}
        \WH_{\mathrm{sp}}(A^\perp)
        &\cong
        \WH_{\mathrm{sp}}(A)^\vee,
        \label{eq:cgqg-span-negation-dual}\\
        \QG_{\mathrm{sp}}(A^\perp)
        &\cong
        \QG_{\mathrm{sp}}(A)^\vee.
        \label{eq:cgqg-span-qg-negation}
    \end{align}
\end{theorem}
\begin{proof}
    \leavevmode
    \begin{enumerate}[beginpenalty=10000]
        \item \textbf{Calculate the compact transpose of a span.}
        Let
        \[
            X\xleftarrow{\ell}S\xrightarrow{r}Y
        \]
        represent a morphism in \(\Sp\). Composing with the canonical
        diagonal coevaluation and evaluation spans leaves the apex
        canonically \(S\), with the two exterior legs exchanged.
        Hence its compact transpose is represented by
        \[
            Y\xleftarrow{r}S\xrightarrow{\ell}X.
        \]

        \item \textbf{Transpose the polarized double groupoid.}
        By Proposition~\ref{prop:cgqg-resource-negation},
        \[
            T(A^\perp)
            \cong
            T(A)^{\mathsf t}.
        \]
        Under this identification, the horizontal and vertical edge
        groupoids are exchanged. Consequently vertical box composition
        for \(A^\perp\) is horizontal box composition for \(A\), while
        horizontal box factorization for \(A^\perp\) is vertical box
        factorization for \(A\).

        \item \textbf{Transpose the incidence structure maps.}
        Since compact transposition of a span exchanges its two legs,
        the preceding identification gives
        \[
            m_{A^\perp}
            \cong
            \Delta_A^\vee,
            \qquad
            \eta_{A^\perp}
            \cong
            \varepsilon_A^\vee,
        \]
        and
        \[
            \Delta_{A^\perp}
            \cong
            m_A^\vee,
            \qquad
            \varepsilon_{A^\perp}
            \cong
            \eta_A^\vee.
        \]
        Total box inversion commutes with transposition:
        \[
            (X^{\mathsf t})^{-1}
            =
            (X^{-1})^{\mathsf t},
        \]
        so
        \[
            S_{A^\perp}
            \cong
            S_A^\vee.
        \]

        By
        Lemma~\ref{lem:cgqg-compact-transpose-weak-hopf}, these
        identifications assemble to an isomorphism of weak Hopf monoids
        \[
            \WH_{\mathrm{sp}}(A^\perp)
            \cong
            \WH_{\mathrm{sp}}(A)^\vee.
        \]
        This proves~\eqref{eq:cgqg-span-negation-dual}.

        \item \textbf{Apply the Pastro--Street construction to the
        transposed weak Hopf monoid.}
        By definition,
        \[
            \QG_{\mathrm{sp}}(A)^\vee
        \]
        is the Pastro--Street quantum groupoid obtained from
        \[
            \WH_{\mathrm{sp}}(A)^\vee.
        \]
        The weak-Hopf isomorphism of Step~3 therefore identifies this
        Pastro--Street quantum groupoid in \(\Kar(\Sp)\) with the one
        obtained from
        \[
            \WH_{\mathrm{sp}}(A^\perp).
        \]

        \item \textbf{Identify the target base.}
        The Pastro--Street target idempotent is determined functorially
        by the weak-bimonoid structure. Hence the weak-Hopf isomorphism
        of Step~3 carries the target idempotent of
        \(\WH_{\mathrm{sp}}(A^\perp)\) to the target idempotent of
        \(\WH_{\mathrm{sp}}(A)^\vee\).

        For \(A^\perp\), the explicit splitting constructed above is
        \[
            \Conf(A^\perp)
            \xrightarrow{i_s}
            \WH_{\mathrm{sp}}(A^\perp)
            \xrightarrow{r_s}
            \Conf(A^\perp).
        \]
        Since
        \[
            \Conf(A^\perp)=\Conf(A),
        \]
        Proposition~\ref{prop:cgqg-configuration-frobenius-retract}
        identifies the corresponding object-of-objects with the same
        canonical separable commutative Frobenius monoid
        \[
            C_A.
        \]

        \item \textbf{Identify the composition object from the
        transposed coactions.}
        The Pastro--Street composition object for
        \(\WH_{\mathrm{sp}}(A)^\vee\) is, by definition, the split
        equalizer of the two boundary coactions determined by that
        transposed weak Hopf structure.

        Under the weak-Hopf isomorphism of Step~3 these are precisely the
        boundary coactions of \(\WH_{\mathrm{sp}}(A^\perp)\). Therefore
        Lemma~\ref{lem:cgqg-local-pastro-street-exactness}, applied to
        \(A^\perp\), supplies the explicit split representative
        \[
            \mathsf{Comp}_{A^\perp}
            =
            \left\{
                (X,Y)
                \in
                \WH_{\mathrm{sp}}(A^\perp)^2
                \;\middle|\;
                r_{A^\perp}(X)
                =
                \ell_{A^\perp}(Y)
            \right\}.
        \]
        Its inclusion and converse split the composition idempotent
        calculated from the transposed weak Hopf structure.

        \item \textbf{Identify composition and unit.}
        On this split composition object, the Pastro--Street composition
        is the restriction of the multiplication of
        \(\WH_{\mathrm{sp}}(A^\perp)\), namely
        \[
            \mu_{A^\perp}
            =
            m_{A^\perp}j_{A^\perp},
        \]
        and the unit is the target-base inclusion
        \[
            u_{A^\perp}
            =
            i_s.
        \]
        Under the weak-Hopf isomorphism of Step~3 these are exactly the
        composition and unit obtained by applying the Pastro--Street
        construction to \(\WH_{\mathrm{sp}}(A)^\vee\).

        \item \textbf{Transport the inversion data.}
        The Pastro--Street inversion maps are constructed from the weak
        Hopf antipode, the target idempotent, and the two comodule
        structures on the Pastro--Street composition object. The
        weak-Hopf isomorphism of Step~3 therefore transports the
        inversion data of \(\WH_{\mathrm{sp}}(A)^\vee\) to those of
        \(\WH_{\mathrm{sp}}(A^\perp)\).

        In the explicit split presentation for \(A^\perp\), the arrow
        inversion is total box inversion and the composition-comodule
        isomorphism is the transported map
        \[
            \theta_{A^\perp}:
            \mathsf{Comp}_{A^\perp,l}
            \xrightarrow{\cong}
            \mathsf{Comp}_{A^\perp,r}.
        \]

        \item \textbf{Conclude quantum-groupoid duality.}
        Thus the Pastro--Street quantum groupoid associated to the
        compactly transposed weak Hopf monoid
        \(\WH_{\mathrm{sp}}(A)^\vee\) is presented in \(\Sp\) by the
        explicit game quantum groupoid of \(A^\perp\). Hence
        \[
            \QG_{\mathrm{sp}}(A^\perp)
            \cong
            \QG_{\mathrm{sp}}(A)^\vee.
        \]
        This proves~\eqref{eq:cgqg-span-qg-negation}. \qedhere
    \end{enumerate}
\end{proof}

\section{Thin strategies as quantum modules}
\label{sec:cgqg-strategy-quantum-modules}

We construct the module associated to a distributor using
Definition~\ref{def:cgqg-background-quantum-module}. The configuration
Frobenius bases express boundary coactions as graphs, and the module
equations become identifications of spans.

\subsection{Configuration-boundary span calculus}

\begin{theorem}[Configuration comodules and boundary spans]
\label{thm:cgqg-split-cotensor}
There is a symmetric monoidal bicategory \(\mathbf{BndSp}\) whose objects
are sets and whose \(1\)-cells \(P\rightsquigarrow Q\) are diagrams
\[
 \begin{tikzcd}[column sep=large,row sep=large]
 &X\arrow[dl,"\ell_X"']\arrow[dr,"r_X"]&\\
 P&&Q.
 \end{tikzcd}
\]
The identity \(1\)-cell on \(P\) is
\[
 \begin{tikzcd}[column sep=large]
 P&P\arrow[l,"1_P"']\arrow[r,"1_P"]&P.
 \end{tikzcd}
\]
Its \(2\)-cells are isomorphism classes of spans
\(X\xleftarrow{u}S\xrightarrow{v}Y\) satisfying
\[
 \ell_Xu=\ell_Yv,\qquad r_Xu=r_Yv.
\]
Horizontal composition of \(1\)-cells is boundary pullback. Vertical
composition of \(2\)-cells is pullback over the intermediate carrier,
horizontal composition is induced by pullback over the common boundary,
and identity \(2\)-cells are identity spans. The assignment \(P\mapsto C_P\),
with graph coactions on \(1\)-cells, gives a symmetric monoidal
biequivalence with the configuration-base fragment of
\(\mathbf{Comod}(\Sp)\). For \(X:P\rightsquigarrow Q\) and
\(Y:Q\rightsquigarrow R\), the cotensor inclusion is the graph
\[
 j:X\times_QY\hookrightarrow X\times Y
\]
and is a split absolute equalizer.
\end{theorem}
\begin{proof}
    \leavevmode
    \begin{enumerate}[beginpenalty=10000]
        \item \textbf{Identify coactions and \(2\)-cells.}
        Proposition~\ref{prop:cgqg-right-mate-boundary-presentation}
        identifies the coactions with the graphs of \(\ell_X,r_X\).
        Intertwining them is precisely preservation of the two boundary
        functions by the legs of a \(2\)-span. The stated pullback formulas
        for vertical and horizontal composition agree with composition of
        comodule maps; associativity, units, and interchange are the canonical
        comparison isomorphisms for iterated pullbacks.

        \item \textbf{Construct the cotensor.}
        Put \(K=X\times_QY\), with its defining pullback
        \[
        \begin{tikzcd}[column sep=large,row sep=large]
        K\arrow[r,"\pi_Y"]\arrow[d,"\pi_X"']
          \arrow[dr,phantom,"\lrcorner",very near start]
          &Y\arrow[d,"\ell_Y"]\\
        X\arrow[r,"r_X"']&Q.
        \end{tikzcd}
        \]
        The two cotensor maps are the graphs
        \[
        f(x,y)=(x,r_X(x),y),\qquad
        g(x,y)=(x,\ell_Y(y),y).
        \]
        Let \(p:X\times Y\to K\) be the converse of \(j\). Define
        \(h:X\times Q\times Y\to X\times Y\) to be the partial span
        supported on \(q=r_X(x)\), with value \((x,y)\). Then
        \[
        pj=1_K,\qquad hf=1_{X\times Y},\qquad hg=jp.
        \]
        If a span \(z:Z\to X\times Y\) satisfies \(fz=gz\), these equations give
        \(z=h fz=h gz=jpz\). Thus \(pz\) factors \(z\) through \(j\),
        uniquely because \(pj=1_K\).

        \item \textbf{Exhibit the splitting.}
        The preceding equations display the cotensor as the split equalizer
        \[
        \begin{tikzcd}[column sep=large,row sep=large]
        K\arrow[r,bend left=15,"j"]
          &X\times Y\arrow[l,bend left=15,"p"]
          \arrow[r,shift left=.8ex,"f"]
          \arrow[r,shift right=.8ex,"g"']
          &X\times Q\times Y\arrow[l,bend left=28,"h"'].
        \end{tikzcd}
        \]
        Split equalizers are absolute, so tensoring preserves this cotensor.

        \item \textbf{Construct the bicategorical and tensor comparisons.}
        Iterated composites are simultaneous boundary pullbacks.
        Rebracketing their tuples gives the associators, and removing a
        repeated boundary gives the unitors. Cartesian product gives
        \(C_P\otimes C_{P'}\cong C_{P\times P'}\) and the bijection
        \[
        \begin{aligned}
        (X\times_QY)\times(X'\times_{Q'}Y')
          &\longrightarrow
          (X\times X')\times_{Q\times Q'}(Y\times Y'),\\
        ((x,y),(x',y'))&\longmapsto((x,x'),(y,y')).
        \end{aligned}
        \]
        Each coherence equation identifies the same rebracketing or
        permutation of tuples, giving the symmetric monoidal biequivalence.
        \qedhere
    \end{enumerate}
\end{proof}

We calculate in \(\mathbf{BndSp}\). For sets \(P,Q\), write
\(E_{P,Q}=P\times Q\). The configuration biduality is represented by
\[
 e_P:\ P^2\xleftarrow{\Delta_P}P\longrightarrow1,
 \qquad
 n_P:\ 1\longleftarrow P\xrightarrow{\Delta_P}P^2.
\]
The snake composites identify the repeated labels and give identity spans.
Put
\[
 L_{P,Q}=E_{P,P}\times E_{P,Q}\times E_{Q,Q}.
\]
The canonical action has boundary span
\[
 \begin{tikzcd}[column sep=large,row sep=large]
 &P\times P\times Q\times Q
   \arrow[dl,"\iota_{\act}"']\arrow[dr,"\tau_{\act}"]&\\
 L_{P,Q}&&E_{P,Q},
 \end{tikzcd}
\]
where
\begin{align}
 \iota_{\act}(a',a,b,b')&=((a',a),(a,b),(b,b')),\\
 \tau_{\act}(a',a,b,b')&=(a',b').
 \label{eq:cgqg-bndsp-canonical-action}
\end{align}
This is \(\act_{P,Q}\) from~\eqref{eq:cgqg-background-canonical-module-action}.

\begin{proposition}[Configuration-base right-mate quantum modules]
\label{prop:cgqg-configuration-qmodule-right-mate}
For canonical configuration bases \(C_{\Conf(A)},C_{\Conf(B)}\), Chikhladze's
formal module data are equivalently represented in \(\mathbf{BndSp}\) by
an endocomonad \(M\) on \(E_{\Conf(A),\Conf(B)}=\Conf(A)\times\Conf(B)\) and an action
\[
 \kappa_M:\act_{\Conf(A),\Conf(B)}(A_1\otimes M\otimes B_1)
 \Longrightarrow M\act_{\Conf(A),\Conf(B)}
\]
satisfying \eqref{eq:cgqg-background-module-counit}--%
\eqref{eq:cgqg-background-module-unit}. Module maps correspond to the
right-mate comonad transformations intertwining these actions.
\end{proposition}
\begin{proof}
    \leavevmode
    \begin{enumerate}[beginpenalty=10000]
        \item \textbf{Transport the comonad data.}
        Proposition~\ref{prop:cgqg-right-mate-boundary-presentation} and
        Theorem~\ref{thm:cgqg-split-cotensor} identify the relevant
        configuration-base comodule hom-categories with boundary-span
        hom-categories. The mutually inverse mate operations preserve
        identities and pullback composition by the snake identities, hence
        transport comonads and their transformations.

        \item \textbf{Transport the action.}
        Under the same correspondence the bidual evaluations give the
        canonical boundary action \eqref{eq:cgqg-bndsp-canonical-action}.
        Thus the action cell, its four equations, and the intertwining
        equation for module maps are transported bijectively. \qedhere
    \end{enumerate}
\end{proof}

\subsection{The distributor comonad and the direct action}
\label{subsec:cgqg-distributor-carrier}

Let \(F:\Sym(A)^\op\times \Sym(B)\to\mathbf{Set}\) be a distributor. Set
\begin{equation}
 W_F=\coprod_{a\in \Conf(A),\,b\in \Conf(B)}F(a,b),
 \qquad d_F(w)=(a,b)\quad(w\in F(a,b)).
 \label{eq:cgqg-distributor-total-set}
\end{equation}
For \(\Omega_A:a'\to a\), \(\Omega_B:b\to b'\), write
\[
 \Omega_{B*}\Omega_A^*w=F(\Omega_A,\Omega_B)(w).
\]
The endpoint transports commute and are functorial.

\begin{proposition}[Diagonal distributor comonad]
\label{prop:cgqg-diagonal-distributor-comonad}
The endo-\(1\)-cell
\[
 \begin{tikzcd}[column sep=large,row sep=large]
 &W_F\arrow[dl,"d_F"']\arrow[dr,"d_F"]&\\
 E_{\Conf(A),\Conf(B)}&&E_{\Conf(A),\Conf(B)}
 \end{tikzcd}
\]
is a comonad \(M_F\), with comultiplication \(w\mapsto(w,w)\) and
counit \(w\mapsto d_F(w)\). Its carrier comonoid in \(\Sp\) is the
diagonal comonoid on \(W_F\).
\end{proposition}
\begin{proof}
    \leavevmode
    \begin{enumerate}[beginpenalty=10000]
        \item \textbf{Construct the comonad cells.}
        The diagonal lands in \(W_F\times_{\Conf(A)\times \Conf(B)}W_F\) and
        preserves the two boundary maps. The graph of \(d_F\) is a
        \(2\)-cell from \(M_F\) to the identity \(1\)-cell on \(E_{\Conf(A),\Conf(B)}\).

        \item \textbf{Verify the equations.}
        Both iterated comultiplications send \(w\) to \((w,w,w)\).
        Either counit removes one repeated copy and returns \(w\).
        Composing with the cotensor inclusion gives the carrier diagonal;
        composing the boundary counit with the terminal map gives
        \(W_F\to1\). \qedhere
    \end{enumerate}
\end{proof}

By Proposition~\ref{prop:cgqg-arrow-right-mate-boundary}, a box \(X\) of
\(A\) has endocomonad boundaries
\[
 d_A^{L}(X)=(\boxnw{X},\boxsw{X}),\qquad d_A^{R}(X)=(\boxne{X},\boxse{X}).
\]
Put \(\act=\act_{\Conf(A),\Conf(B)}\) and \(D_F=\WH_{\mathrm{sp}}(A)\otimes M_F\otimes \WH_{\mathrm{sp}}(B)\).
The carrier of \(\act D_F\) is
\begin{equation}
 K_F=\{(X,w,Y):w\in F(a,b),\ \boxse{X}=a,\ \boxne{Y}=b\},
 \label{eq:cgqg-qmodule-action-source}
\end{equation}
with boundary maps
\begin{equation}
 \begin{split}
 d^{L}_{K_F}(X,w,Y)&=((\boxnw{X},\boxsw{X}),(a,b),(\boxnw{Y},\boxsw{Y})),\\
 d^{R}_{K_F}(X,w,Y)&=(\boxne{X},\boxse{Y}).
 \end{split}
 \label{eq:cgqg-qmodule-action-source-boundaries}
\end{equation}
The carrier of \(M_F\act\) is
\[
 J_F=\Conf(A)\times \Conf(B)\times W_F.
\]
For \(w'\in F(a',b')\), its boundaries are
\begin{equation}
 \begin{split}
 d^{L}_{J_F}(a,b,w')&=((a',a),(a,b),(b,b')),\\
 d^{R}_{J_F}(a,b,w')&=(a',b').
 \end{split}
 \label{eq:cgqg-qmodule-action-target-boundaries}
\end{equation}
Both descriptions follow by composing with
\eqref{eq:cgqg-bndsp-canonical-action}; the labels \(a,b\) are the
contraction witnesses in \(J_F\).

Define the action apex
\begin{equation}
 \operatorname{Act}_F=
 \left\{(v,w,u)\ \middle|\
 \begin{array}{l}
 v:a'\to a\text{ in }\Sym^+(A),\\
 w\in F(a,b),\\
 u:b\to b'\text{ in }\Sym^+(B)
 \end{array}\right\}.
 \label{eq:cgqg-distributor-action-witness-apex}
\end{equation}
The direct action is the span
\begin{equation}
 \begin{tikzcd}[column sep=large,row sep=large]
 &\operatorname{Act}_F
   \arrow[dl,"\iota_F"']\arrow[dr,"\tau_F"]&\\
 K_F&&J_F,
 \end{tikzcd}
 \qquad \kappa_F=[\iota_F,\tau_F],
 \label{eq:cgqg-distributor-action-span}
\end{equation}
with legs
\begin{equation}
 \iota_F(v,w,u)=(\id_v^{\mathrm h},w,\id_u^{\mathrm h}),\qquad
 \tau_F(v,w,u)=(a,b,u_*v^*w).
 \label{eq:cgqg-distributor-action-legs}
\end{equation}
Its two source boundaries are \(((a',a),(a,b),(b,b'))\), and its two
target boundaries are \((a',b')\). Thus it defines
\(\kappa_F:\act D_F\Rightarrow M_F\act\). The endpoint transport appears as
\[
 \begin{tikzcd}[column sep=huge,row sep=large]
 F(a,b)\arrow[r,"v^*"]\arrow[d,"u_*"']
   &F(a',b)\arrow[d,"u_*"]\\
 F(a,b')\arrow[r,"v^*"']&F(a',b').
 \end{tikzcd}
\]

\begin{proposition}[The direct action and its multitensor mate]
\label{prop:cgqg-direct-action-mate}
The multitensor \(T_3(\WH_{\mathrm{sp}}(A),M_F,\WH_{\mathrm{sp}}(B))\) has carrier
\[
 \{(X,w,Y):\boxsw{X}=\boxse{X}=a(w),\ \boxnw{Y}=\boxne{Y}=b(w)\}.
\]
The Kan-extension mate of \(\kappa_F\) has apex \(\operatorname{Act}_F\)
and legs
\[
 (v,w,u)\longmapsto(\id_v^{\mathrm h},w,\id_u^{\mathrm h}),\qquad
 (v,w,u)\longmapsto u_*v^*w.
\]
\end{proposition}
\begin{proof}
    \leavevmode
    \begin{enumerate}[beginpenalty=10000]
        \item \textbf{Identify the matching locus.}
        Equality of the source boundaries of \(K_F\) and \(J_F\)
        imposes \(\boxsw{X}=a\) and \(\boxnw{Y}=b\), in addition to
        \(\boxse{X}=a\) and \(\boxne{Y}=b\). These are exactly the two-boundary
        equations of Proposition~\ref{prop:cgqg-background-span-multitensors}.

        \item \textbf{Construct the correspondence of \(2\)-cells.}
        On this locus, a boundary-preserving span to \(M_F\act\)
        determines a span to \(M_F\) by forgetting the uniquely determined
        contraction labels \(a,b\). Inserting those labels is inverse
        to this operation. Applying this correspondence to
        \eqref{eq:cgqg-distributor-action-legs} gives the displayed mate. \qedhere
    \end{enumerate}
\end{proof}

\subsection{The quantum-module equations}

\begin{theorem}[Distributor-induced quantum module]
\label{thm:cgqg-distributor-qmodule}
Every distributor \(F:\Sym(A)^\op\times \Sym(B)\to\mathbf{Set}\) canonically
determines a Chikhladze quantum module
\[
 Q_F^{\mathrm{sp}}:
 \QG_{\mathrm{sp}}(A)\xnrightarrow{}\QG_{\mathrm{sp}}(B),
\]
whose right-mate presentation has comonad \(M_F\) and direct action
\(\kappa_F\).
\end{theorem}
\begin{proof}
    \leavevmode
    \begin{enumerate}[beginpenalty=10000]
        \item \textbf{Verify the counit equation.}
        The comonad counit equation is
        \[
        \begin{tikzcd}[column sep=large,row sep=large]
        \act D_F\arrow[r,"\kappa_F"]\arrow[d,"\act\epsilon_{D_F}"']
          &M_F\act\arrow[dl,"\epsilon_{M_F}\act"]\\
        \act.&
        \end{tikzcd}
        \]
        On the direct route the two arrow counits select horizontal
        identity boxes, and the distributor counit is defined on every
        \(w\). Its apex is therefore \(\operatorname{Act}_F\).
        The route through \(\kappa_F\) has the same apex. Both spans
        take \((v,w,u)\) to the contraction witness \((a',a,b,b')\)
        and have the same input \((\id_v^{\mathrm h},w,\id_u^{\mathrm h})\).

        \item \textbf{Verify the comultiplication equation.}
        The diagram to be checked is
        \[
        \begin{tikzcd}[column sep=large,row sep=large]
        \act D_F\arrow[r,"\act\delta_{D_F}"]\arrow[d,"\kappa_F"']
          &\act D_F^2\arrow[r,"\kappa_F D_F"]
          &M_F\act D_F\arrow[d,"M_F\kappa_F"]\\
        M_F\act\arrow[rr,"\delta_{M_F}\act"']&&M_F^2\act.
        \end{tikzcd}
        \]
        Equivalently,
        \begin{equation}
        (\delta_{M_F}\act)\kappa_F
          =(M_F\kappa_F)(\kappa_FD_F)(\act\delta_{D_F}).
        \label{eq:cgqg-direct-action-comultiplication}
        \end{equation}
        A witness on the upper route chooses horizontal factorizations
        of \(X,Y\) and duplicates \(w\). The two actions require all
        four box factors to be horizontal identities. Horizontally
        composable horizontal identities lie on the same vertical edge,
        giving the unique surviving factorizations
        \[
        X=\id_v^{\mathrm h}=\id_v^{\mathrm h}\circ_{\mathrm h}\id_v^{\mathrm h},\qquad
        Y=\id_u^{\mathrm h}=\id_u^{\mathrm h}\circ_{\mathrm h}\id_u^{\mathrm h}.
        \]
        Both routes therefore have apex \(\operatorname{Act}_F\).
        Their output consists of two copies of \(u_*v^*w\) and the
        contraction labels \(a,b\), with identical exterior legs.

        \item \textbf{Verify associativity of the action.}
        Use \(\pi,\lambda\) from
        Definition~\ref{def:cgqg-background-quantum-module} and put
        \[
        \chi=\act\pi\cong \act\lambda,\qquad
        D_F^{(2)}=\WH_{\mathrm{sp}}(A)\otimes \WH_{\mathrm{sp}}(A)\otimes M_F\otimes \WH_{\mathrm{sp}}(B)\otimes \WH_{\mathrm{sp}}(B).
        \]
        Write
        \[
        L_F=\act(\phi_A\otimes1_{M_F}\otimes\phi_B),\qquad
        R_F=\act(1_{\WH_{\mathrm{sp}}(A)}\otimes\kappa_F\otimes1_{\WH_{\mathrm{sp}}(B)}).
        \]
        The action equation is the commutativity of
        \[
        \begin{tikzcd}[column sep=large,row sep=large]
        \chi D_F^{(2)}\arrow[r,"L_F"]\arrow[d,"R_F"']
          &\act D_F\pi\arrow[d,"\kappa_F\pi"]\\
        \act D_F\lambda\arrow[r,"\kappa_F\lambda"']&M_F\chi.
        \end{tikzcd}
        \]
        On the multiplication-first route, the action selects
        horizontal identity vertical pastes. By
        Lemma~\ref{lem:cgqg-vacant-cell-calculus}, every vertical
        factor of such a paste is a horizontal identity. The
        action-first route has the same support. The common apex is
        the set of tuples \((v_1,v_2,w,u_1,u_2)\) with
        \[
        \begin{gathered}
        v_1:a_0\to a_1,\quad v_2:a_1\to a_2,\quad
        w\in F(a_2,b_0),\\
        u_1:b_0\to b_1,\quad u_2:b_1\to b_2,
        \end{gathered}
        \]
        where the four arrows are positive. Both exterior input legs
        give
        \((\id_{v_1}^{\mathrm h},\id_{v_2}^{\mathrm h},w,
        \id_{u_1}^{\mathrm h},\id_{u_2}^{\mathrm h})\). The outputs are
        \[
        (u_2\circ u_1)_*(v_2\circ v_1)^*w
          =u_{2*}v_1^*u_{1*}v_2^*w
        \]
        by functoriality and commuting endpoint actions. Both output
        spans retain the same contraction labels \(a_1,a_2,b_0,b_1\).
        Thus the complete spans agree.

        \item \textbf{Verify the unit equation.}
        Let \(\jmath\) be the canonical unit insertion and put
        \(U_F=\act(\phi_A^0\otimes1_{M_F}\otimes\phi_B^0)\). Under
        \(\act\jmath\cong1\), the unit equation is
        \[
        \begin{tikzcd}[column sep=large,row sep=large]
        M_F\arrow[r,"U_F"]\arrow[dr,equal]
          &\act D_F\jmath\arrow[d,"\kappa_F\jmath"]\\
          &M_F.
        \end{tikzcd}
        \]
        The quantum-category units select vertical identity boxes.
        Intersecting their support with the horizontal identities
        selected by the action gives the double identities:
        \[
        \{\id_h^{\mathrm v}:h\in\Sym^-_1(A)\}
        \cap\{\id_v^{\mathrm h}:v\in\Sym^+_1(A)\}
          =\{\Theta_a:a\in \Conf(A)\},
        \]
        and similarly for \(B\). Over \(w\in F(a,b)\) there is
        exactly one surviving witness, \((\id_a,w,\id_b)\), and
        its output is \(w\). Both exterior legs are identities.
        The four equations establish the required right-mate module action;
        Proposition~\ref{prop:cgqg-configuration-qmodule-right-mate}
        transports it to the formal Chikhladze quantum module. \qedhere
    \end{enumerate}
\end{proof}

\subsection{Naturality and tensor product}

\begin{proposition}[Functoriality and tensor compatibility]
\label{prop:cgqg-distributor-qmodule-functorial}
\label{prop:cgqg-distributor-qmodule-tensor}
For fixed \(A,B\), a natural transformation \(\eta:F\Rightarrow F'\)
induces a quantum-module map
\(Q_\eta^{\mathrm{sp}}:Q_F^{\mathrm{sp}}\to Q_{F'}^{\mathrm{sp}}\),
functorially. External products give canonical isomorphisms
\begin{equation}
 Q_{F\extprod F'}^{\mathrm{sp}}
   \cong Q_F^{\mathrm{sp}}\otimes Q_{F'}^{\mathrm{sp}},
 \label{eq:cgqg-distributor-qmodule-tensor}
\end{equation}
over the object tensor comparisons of Theorem~\ref{thm:cgqg-qg-tensor}.
They are natural and satisfy associativity, unit, and symmetry coherence.
\end{proposition}
\begin{proof}
    \leavevmode
    \begin{enumerate}[beginpenalty=10000]
        \item \textbf{Construct the comonad transformation.}
        Define \(\widehat\eta(w)=\eta_{a,b}(w)\) for \(w\in F(a,b)\).
        Its graph preserves \(d_F\), the diagonal, and the counit.
        Hence it gives a comonad transformation \(M_F\Rightarrow M_{F'}\).

        \item \textbf{Intertwine the actions.}
        Naturality gives
        \(\widehat\eta(u_*v^*w)=u_*v^*\widehat\eta(w)\).
        Sending \((v,w,u)\) to \((v,\widehat\eta(w),u)\) identifies
        the apices of the two composites in
        \[
        \begin{tikzcd}[column sep=large,row sep=large]
        \act D_F\arrow[r,"\kappa_F"]
          \arrow[d,"{\act(1\otimes\widehat\eta\otimes1)}"']
          &M_F\act\arrow[d,"\widehat\eta \act"]\\
        \act D_{F'}\arrow[r,"\kappa_{F'}"']&M_{F'}\act.
        \end{tikzcd}
        \]
        Both exterior legs agree by the displayed naturality equation.
        Graphs preserve identity maps and composition, proving functoriality.

        \item \textbf{Construct the tensor comparison.}
        The external product has carrier
        \(W_{F\extprod F'}\cong W_F\times W_{F'}\).
        The configuration and box product identifications give
        \(\operatorname{Act}_{F\extprod F'}
          \cong\operatorname{Act}_F\times\operatorname{Act}_{F'}\).
        Under these bijections, the two action legs, the comonad
        diagonal, and its counit are componentwise products.
        This gives~\eqref{eq:cgqg-distributor-qmodule-tensor}.

        \item \textbf{Verify tensor coherence.}
        The associativity and unit comparisons rebracket tuples or
        remove the singleton coordinate of the tensor unit. The symmetry
        comparison exchanges the two factors. Both sides of each
        coherence equation have the same action apex and the same
        permutation of its coordinates. Naturality follows by applying
        the component maps to these coordinates. \qedhere
    \end{enumerate}
\end{proof}

\begin{proposition}[Positive restriction and quantum-module maps]
\label{prop:cgqg-positive-restriction}
For a distributor \(F:\Sym(A)^{\op}\times \Sym(B)\to\mathbf{Set}\), put
\[
 F^+=F\big|_{\Sym^+(A)^{\op}\times\Sym^+(B)}.
\]
The quantum module \(Q_F^{\mathrm{sp}}\) depends only on \(F^+\).
For fixed games \(A,B\), there is a canonical bijection
\begin{equation}
 \operatorname{Hom}_{\mathrm{qmod}}
   (Q_F^{\mathrm{sp}},Q_{F'}^{\mathrm{sp}})
 \cong\operatorname{Nat}(F^+,F'^+),
 \label{eq:cgqg-positive-restriction-module-maps}
\end{equation}
where the left-hand side consists of the fixed-endpoint quantum-module
maps of Definition~\ref{def:cgqg-background-quantum-module}.
\end{proposition}
\begin{proof}
    \leavevmode
    \begin{enumerate}[beginpenalty=10000]
        \item \textbf{Identify the data used in the construction.}
        The positive symmetry subgroupoids are wide, so restriction
        retains every fibre \(F(a,b)\) and the endpoint map \(d_F\).
        It therefore retains the diagonal comonad. The action
        in~\eqref{eq:cgqg-distributor-action-legs} uses only
        \(u_*v^*\) for positive \(u,v\), and is determined by \(F^+\).

        \item \textbf{Recover a function from a module map.}
        Represent its underlying comonad transformation by a boundary
        span \(W_F\xleftarrow{a}S\xrightarrow{b}W_{F'}\).
        Preservation of the boundary counit says that
        \[
         [W_F\xleftarrow{a}S\xrightarrow{d_{F'}\circ b}\Conf(A)\times \Conf(B)]
         =[W_F\xleftarrow{1}W_F\xrightarrow{d_F}\Conf(A)\times \Conf(B)].
        \]
        An apex isomorphism therefore identifies \(a\) with a
        bijection. The module map is the graph of the unique function
        \(f=b\circ a^{-1}\), and \(d_{F'}\circ f=d_F\). Such a graph
        automatically preserves the diagonal comultiplication.

        \item \textbf{Identify action compatibility.}
        The two composite action spans have the same possible inputs
        \((\id_v^{\mathrm h},w,\id_u^{\mathrm h})\) and one witness for each such input.
        They agree exactly when
        \[
          f(u_*v^*w)=u_*v^*f(w)
        \]
        for all positive endpoint arrows. Together with endpoint
        preservation, these equations are precisely naturality of
        \(F^+\Rightarrow F'^+\).

        \item \textbf{Construct the inverse correspondence.}
        A natural transformation of the restrictions gives an
        endpoint-preserving function with these equations. Its graph
        preserves the two comonad maps and intertwines the actions,
        hence is a quantum-module map. The two constructions are
        inverse and preserve identities and composition. \qedhere
    \end{enumerate}
\end{proof}

\begin{corollary}[Classical distributor specialization]
\label{cor:cgqg-qmodule-classical-variance}
If the negative symmetry groupoids of \(A,B\) are discrete,
\(Q_F^{\mathrm{sp}}\) is the ordinary span realization of \(F\).
\end{corollary}
\begin{proof}
    \leavevmode
    \begin{enumerate}[beginpenalty=10000]
        \item \textbf{Identify the arrow structures.}
        Polarized decomposition gives
        \[
        \Sym^+(A)=\Sym(A),\qquad \Sym^+(B)=\Sym(B).
        \]
        Every box is a horizontal identity,
        identified with its positive edge. Vertical multiplication is
        groupoid composition and horizontal comultiplication is diagonal.

        \item \textbf{Identify the action.}
        The action span is the ordinary distributor action
        \((v,w,u)\mapsto u_*v^*w\), over the endpoint-indexed carrier \(W_F\). \qedhere
    \end{enumerate}
\end{proof}

\subsection{Thin strategies and copycat}
\label{subsec:cgqg-thin-strategy-qmodules}

\begin{theorem}[Game-generated quantum module]
\label{thm:cgqg-game-generated-qmodule}
Every thin strategy \(\sigma:A\to B\) canonically determines
\[
 Q_\sigma^{\mathrm{sp}}=
 Q_{\lVert\sigma\rVert}^{\mathrm{sp}}:
 \QG_{\mathrm{sp}}(A)\xnrightarrow{}\QG_{\mathrm{sp}}(B).
\]
Globular positive morphisms induce quantum-module maps, and tensor
products of strategies induce tensor products of these modules up to
the canonical isomorphisms.
\end{theorem}
\begin{proof}
    \leavevmode
    \begin{enumerate}[beginpenalty=10000]
        \item \textbf{Construct the module.}
        Apply Theorem~\ref{thm:cgqg-distributor-qmodule} to the
        positive-witness distributor of
        Theorem~\ref{thm:cgqg-background-positive-witness-double-functor}.

        \item \textbf{Transport morphisms and tensor products.}
        The collapse sends globular positive morphisms to natural
        transformations and supplies its symmetric monoidal comparisons.
        Proposition~\ref{prop:cgqg-distributor-qmodule-functorial}
        carries these to the stated module maps and tensor isomorphisms. \qedhere
    \end{enumerate}
\end{proof}

\begin{proposition}[The quantum module of copycat]
\label{prop:cgqg-copycat-quantum-module}
For every thin game \(A\), composition of the polarized frames gives
a natural isomorphism
\begin{equation}
 \zeta_A:\lVert\cc_A\rVert\xrightarrow{\cong}\Sym(A)(-,-),
 \qquad
 (\theta^-,\cc_Ax,\theta^+)\longmapsto\theta^+\circ\theta^-.
 \label{eq:cgqg-copycat-witness-comparison}
\end{equation}
Its quantum image is the canonical isomorphism
\[
 Q_{\zeta_A}^{\mathrm{sp}}:
 Q_{\cc_A}^{\mathrm{sp}}\xrightarrow{\cong}Q_{\Sym(A)(-,-)}^{\mathrm{sp}}.
\]
\end{proposition}
\begin{proof}
    \leavevmode
    \begin{enumerate}[beginpenalty=10000]
        \item \textbf{Identify the copycat witnesses.}
        The \(+\)-covered configurations of copycat are \(\cc_Ax\),
        with display \(x\otimes x\). A witness from \(a\) to \(b\)
        is therefore the factorization triangle
        \[
        \begin{tikzcd}[column sep=huge,row sep=large]
        a\arrow[rr,"\zeta_A(w)"]\arrow[dr,"\theta^-"']
          &&b\\
          &x\arrow[ur,"\theta^+"']&.
        \end{tikzcd}
        \]

        \item \textbf{Construct the inverse.}
        For \(\omega:a\to b\), Theorem~\ref{thm:cgqg-polarized-decomposition}
        gives the unique factorization \(\omega=\omega^+\circ\omega^-\),
        through \(x\). Send \(\omega\) to
        \((\omega^-,\cc_Ax,\omega^+)\). Uniqueness makes this the
        inverse of~\eqref{eq:cgqg-copycat-witness-comparison}.

        \item \textbf{Verify naturality.}
        For endpoint symmetries \(\Omega_A:a'\to a\), \(\Omega_B:b\to b'\),
        write the transported witness as
        \((\vartheta^-,\cc_Ay,\vartheta^+)\).
        The two components of the copycat symmetry are the same
        map \(\varphi:x\to y\). Endpoint transport gives
        \[
        \vartheta^-=\varphi\circ\theta^-\circ\Omega_A,\qquad
        \vartheta^+\circ\varphi=\Omega_B\circ\theta^+.
        \]
        Multiplying these equations proves commutativity of
        \[
        \begin{tikzcd}[column sep=large,row sep=large]
        \Wit_{\cc_A}^+(a,b)\arrow[r,"\zeta_A"]
          \arrow[d,"\Omega_{B*}\Omega_A^*"']
          &\Sym(A)(a,b)\arrow[d,"{\omega\mapsto\Omega_B\circ\omega\circ\Omega_A}"]\\
        \Wit_{\cc_A}^+(a',b')\arrow[r,"\zeta_A"']&\Sym(A)(a',b').
        \end{tikzcd}
        \]
        This is the copycat comparison of
        \cite[Proposition~5.12]{ClairambaultOlimpieriPaquet2025}.

        \item \textbf{Identify the quantum-module isomorphism.}
        The graph of \(\zeta_A\) preserves the endpoint map, diagonal,
        and counit. For positive \(v:a'\to a\), \(u:b\to b'\),
        naturality gives
        \[
        \zeta_A(u_*v^*w)=u\circ\zeta_A(w)\circ v.
        \]
        The two action spans therefore intertwine under this graph.
        Its inverse gives the inverse module map. Tensor decomposition
        identifies \(\zeta_{A\otimes B}\) with the product of
        \(\zeta_A,\zeta_B\), since composition of symmetry arrows is
        componentwise. \qedhere
    \end{enumerate}
\end{proof}

\section{Quantum realization as a symmetric monoidal double functor}
\label{sec:cgqg-composition-scope}

We assemble the game quantum groupoids and distributor-induced modules into
a double category. A horizontal arrow consists of a distributor together
with its quantum module, and horizontal composition is induced by the
distributor coend. Copycat gives the generated horizontal identity, and companion
distributors supply the vertical arrows and squares.

\subsection{The generated horizontal bicategory}
\label{subsec:cgqg-generated-qmodule-bicategory}

For a thin game \(A\), define
\begin{equation}
 \mathbb Q_A=(A,\Sym(A),\QG_{\mathrm{sp}}(A)).
 \label{eq:cgqg-generated-quantum-object}
\end{equation}
For \(F:\Sym(A)^\op\times \Sym(B)\to\mathbf{Set}\), define
\begin{equation}
 \mathbb M_F=(F,Q_F^{\mathrm{sp}}):
 \mathbb Q_A\xnrightarrow{}\mathbb Q_B.
 \label{eq:cgqg-generated-horizontal-arrow-pair}
\end{equation}
A natural transformation \(\eta:F\Rightarrow F'\) gives the globular cell
\(\mathbb M_\eta=(\eta,Q_\eta^{\mathrm{sp}})\).
These quantum components are supplied by
Theorem~\ref{thm:cgqg-distributor-qmodule} and
Proposition~\ref{prop:cgqg-distributor-qmodule-functorial}.

For composable distributors \(F:\Sym(A)\xnrightarrow{}\Sym(B)\) and
\(G:\Sym(B)\xnrightarrow{}\Sym(C)\), put
\begin{equation}
 \mathbb M_G\star\mathbb M_F=\mathbb M_{G\circ F},
 \qquad
 (G\circ F)(a,c)=\int^{b\in \Sym(B)}F(a,b)\times G(b,c).
 \label{eq:cgqg-generated-horizontal-composition}
\end{equation}
For fixed \(a,c\), this coend has the coequalizer presentation in sets
\[
 \begin{tikzcd}[column sep=large,row sep=large]
 \displaystyle\coprod_{r:b\to b'}F(a,b)\times G(b',c)
   \arrow[r,shift left=.7ex,"d_0"]
   \arrow[r,shift right=.7ex,"d_1"']
   &\displaystyle\coprod_b F(a,b)\times G(b,c)
       \arrow[r,"q"]&(G\circ F)(a,c),
 \end{tikzcd}
\]
where \(d_0(r,x,y)=(r_*x,y)\), \(d_1(r,x,y)=(x,r^*y)\). Thus
\begin{equation}
 [r_*x,y]=[x,r^*y],
 \qquad r:b\to b',\quad x\in F(a,b),\quad y\in G(b',c).
 \label{eq:cgqg-full-coend-relation}
\end{equation}
The exterior transports
\[
 h^*[x,y]=[h^*x,y],\qquad k_*[x,y]=[x,k_*y]
\]
commute with these relations and make the coend a distributor.
The generated horizontal identity is
\begin{equation}
 \mathbb I_A^{\mathrm{gen}}=\mathbb M_{\Sym(A)(-,-)}.
 \label{eq:cgqg-generated-horizontal-unit}
\end{equation}
Proposition~\ref{prop:cgqg-copycat-quantum-module} identifies its quantum
component with \(Q_{\cc_A}^{\mathrm{sp}}\).

\begin{proposition}[Generated horizontal bicategory]
\label{prop:cgqg-generated-qmodule-bicategory}
The objects \(\mathbb Q_A\), arrows \(\mathbb M_F\), and cells
\(\mathbb M_\eta\) form a bicategory with composition \(\star\) and
identities \(\mathbb I_A^{\mathrm{gen}}\).
\end{proposition}
\begin{proof}
    \leavevmode
    \begin{enumerate}[beginpenalty=10000]
        \item \textbf{Construct composites of arrows and cells.}
        The full coend and the hom-distributor are distributors, so
        Theorem~\ref{thm:cgqg-distributor-qmodule} supplies their quantum
        modules. Horizontal composition of natural transformations is
        \[
        [x,y]\longmapsto[\eta(x),\zeta(y)].
        \]
        Naturality identifies the images of the two representatives in
        \eqref{eq:cgqg-full-coend-relation}. Applying
        \(Q_{(-)}^{\mathrm{sp}}\) gives the quantum component.

        \item \textbf{Construct the associator.}
        For \(F,G,H\), the two iterated coends are quotients of triples
        by the same two middle-symmetry relations. The bijection
        \[
        [x,[y,z]]\longmapsto[[x,y],z]
        \]
        therefore induces
        \[
        \mathfrak a_{H,G,F}:
        (\mathbb M_H\star\mathbb M_G)\star\mathbb M_F
        \xrightarrow{\cong}
        \mathbb M_H\star(\mathbb M_G\star\mathbb M_F).
        \]
        The inverse rebrackets in the opposite direction. Both are
        natural, and their quantum images are inverse module maps.

        \item \textbf{Construct the unitors.}
        The co-Yoneda maps and their inverses are
        \[
        \begin{aligned}
        [h,w]&\longmapsto h^*w,
        &w&\longmapsto[\id_a,w],\\
        [w,k]&\longmapsto k_*w,
        &w&\longmapsto[w,\id_b],
        \end{aligned}
        \qquad w\in F(a,b).
        \]
        The two inverse equations are instances of
        \eqref{eq:cgqg-full-coend-relation}. Their quantum images give
        \[
        \begin{tikzcd}[column sep=large,row sep=large]
        \mathbb I_B^{\mathrm{gen}}\star\mathbb M_F
          \arrow[dr,"\lambda_F"']
          &&\mathbb M_F\star\mathbb I_A^{\mathrm{gen}}
              \arrow[dl,"\rho_F"]\\
          &\mathbb M_F.&
        \end{tikzcd}
        \]

        \item \textbf{Verify coherence and interchange.}
        Every pentagon path sends a quadruple representative to the
        same rebracketed quadruple. Both triangle paths remove the
        identity factor by the same endpoint transport. The interchange
        equation applies the component natural transformations to the
        same pair \((x,y)\). These are equalities of distributor maps;
        functoriality of \(Q_{(-)}^{\mathrm{sp}}\) gives the corresponding
        equalities of quantum-module maps. \qedhere
    \end{enumerate}
\end{proof}

Let \(\mathbb R_A^{\mathrm{reg}}\) denote the regular quantum module
of \(\QG_{\mathrm{sp}}(A)\), with carrier \(\WH_{\mathrm{sp}}(A)\) and action induced
by multiplication. Its comparison with \(\mathbb I_A^{\mathrm{gen}}\)
is a comparison with the quantum component \(Q_{\Sym(A)(-,-)}^{\mathrm{sp}}\).

\begin{proposition}[Copycat and the regular quantum module]
\label{prop:cgqg-copycat-unit-iff}
The quantum component \(Q_{\Sym(A)(-,-)}^{\mathrm{sp}}\) of the generated
horizontal identity and the regular quantum module
\(\mathbb R_A^{\mathrm{reg}}\) are isomorphic if and only if
\[
 \Sym^-(A)=\Disc(\Conf(A)).
\]
Through \(Q_{\zeta_A}^{\mathrm{sp}}\), this also characterizes when the
quantum module of copycat is isomorphic to the regular quantum module.
\end{proposition}
\begin{proof}
    \leavevmode
    \begin{enumerate}[beginpenalty=10000]
        \item \textbf{Identify the generated carrier.}
        Its carrier is \(\coprod_{a,b}\Sym(A)(a,b)\), with diagonal
        comultiplication and terminal counit. An invertible morphism
        of \(\Sp\) is represented by a bijection: composing with its
        inverse gives identity spans, forcing a unique witness over
        each exterior element.

                \item \textbf{Derive discreteness from an identity comparison.}
        A quantum-module isomorphism induces an isomorphism of the
        underlying endocomonads. By the argument of
        Proposition~\ref{prop:cgqg-positive-restriction}, its underlying
        span is therefore the graph of a boundary-preserving bijection
        \[
            f:
            \coprod_{a,b}\Sym(A)(a,b)
            \xrightarrow{\cong}
            \WH_{\mathrm{sp}}(A).
        \]
        The boundary counit of the distributor comonad is defined on
        every element of its carrier, whereas the boundary counit of the
        arrow endocomonad is
        \[
            \Boxes_A
            \xleftarrow{v\mapsto\id_v^{\mathrm h}}
            \Sym^+_1(A)
            \xrightarrow{v\mapsto(s(v),t(v))}
            \Conf(A)\times\Conf(A)
        \]
        by Proposition~\ref{prop:cgqg-arrow-right-mate-boundary}.
        Preservation of the boundary counit therefore implies that
        every value of \(f\) is a horizontal identity box. Since \(f\)
        is surjective, every box in \(\Boxes_A\) is a horizontal
        identity.

        In particular, for every \(h\in\Sym^-_1(A)\), the vertical
        identity \(\id_h^{\mathrm v}\) is also a horizontal identity.
        A box which is both a vertical identity and a horizontal
        identity is a double identity, so \(h\) is an identity arrow.
        Hence
        \[
            \Sym^-(A)=\Disc(\Conf(A)).
        \]

        \item \textbf{Construct the comparison under discreteness.}
        Polarized decomposition gives \(\Sym(A)=\Sym^+(A)\), and every
        box is \(\id_v^{\mathrm h}\) for a unique \(v\in \Sym(A)\). The bijection
        \[
        \coprod_{a,b}\Sym(A)(a,b)\longrightarrow \WH_{\mathrm{sp}}(A),\qquad
        v\longmapsto\id_v^{\mathrm h}
        \]
        preserves both endocomonad boundaries. Horizontal
        comultiplication is the diagonal and its counit is terminal.
        The regular and distributor actions both compose the positive
        edges, so this bijection is a quantum-module isomorphism.
        Composing with \(Q_{\zeta_A}^{\mathrm{sp}}\) gives the stated
        comparison with copycat. \qedhere
    \end{enumerate}
\end{proof}

\begin{proposition}[Full symmetry transport in composition]
\label{prop:cgqg-full-transport-example}
There are thin games \(I,B\), distributors \(F_0,F_1:\Sym(I)\xnrightarrow{}\Sym(B)\),
and \(N:\Sym(B)\xnrightarrow{}\Sym(I)\), such that
\[
 Q_{F_0}^{\mathrm{sp}}=Q_{F_1}^{\mathrm{sp}},
 \qquad
 |(N\circ F_0)(*,*)|=2,\qquad |(N\circ F_1)(*,*)|=1.
\]
\end{proposition}
\begin{proof}
    \leavevmode
    \begin{enumerate}[beginpenalty=10000]
        \item \textbf{Construct the game and distributors.}
        Let \(I\) be the empty game and let \(B\) have two concurrent
        negative events, with all selected bijections. Put
        \(b=\{e_1,e_2\}\). Its automorphism group is \(C_2\), and
        \(\Sym^+(B)\) is discrete. Define \(F_0,F_1\) to have fibre
        \(\{0,1\}\) at \(b\) and empty fibres elsewhere, with respectively
        the trivial and swapping \(C_2\)-actions. Define \(N\) to have
        singleton fibre at \(b\), trivial action, and empty fibres
        elsewhere. Symmetries preserve configuration cardinality, so
        these prescriptions define distributors.

        \item \textbf{Identify their quantum modules.}
        The two carriers and endpoint maps agree. Their diagonals
        and counits agree as well. The action in
        \eqref{eq:cgqg-distributor-action-legs} uses positive endpoint
        transport, which is the identity transport in this example.
        Hence the two comonads and action spans coincide.

        \item \textbf{Compute the composites.}
        Formula~\eqref{eq:cgqg-full-coend-relation} gives the orbit sets
        \[
        (N\circ F_i)(*,*)=F_i(*,b)/C_2.
        \]
        The trivial action has orbits \(\{0\},\{1\}\); the swapping
        action has the single orbit \(\{0,1\}\). This gives the two
        displayed cardinalities. Proposition~\ref{prop:cgqg-positive-restriction}
        identifies the common restriction retained by their quantum
        modules. The retained distributor in a generated arrow specifies
        the additional symmetry transport used by this full coend. \qedhere
    \end{enumerate}
\end{proof}

\subsection{Representable vertical arrows and squares}
\label{subsec:cgqg-representable-quantum-maps}

A functor \(u:\Sym(A)\to \Sym(C)\) has companion and conjoint distributors
\begin{equation}
 u_*(a,c)=\Sym(C)(u(a),c),\qquad
 u^*(c,a)=\Sym(C)(c,u(a)).
 \label{eq:cgqg-companion-distributors}
\end{equation}
Applying Theorem~\ref{thm:cgqg-distributor-qmodule} gives the generated arrows
\[
 \begin{tikzcd}[column sep=huge]
 \mathbb Q_A
   \arrow[r,bend left=18,Rightarrow,"\mathbb M_{u_*}"]
   &\mathbb Q_C
       \arrow[l,bend left=18,Rightarrow,"\mathbb M_{u^*}"].
 \end{tikzcd}
\]

\begin{proposition}[Representable adjunction and composition]
\label{prop:cgqg-representable-quantum-adjunction}
\label{prop:cgqg-representable-vertical-composition}
There is a canonical adjunction
\(\mathbb M_{u_*}\dashv\mathbb M_{u^*}\) in the generated bicategory.
For composable \(u,v\), there are coherent canonical isomorphisms
\[
 \mathbb M_{v_*}\star\mathbb M_{u_*}\cong\mathbb M_{(v\circ u)_*},
 \qquad
 \mathbb M_{u^*}\star\mathbb M_{v^*}\cong\mathbb M_{(v\circ u)^*}.
\]
\end{proposition}
\begin{proof}
    \leavevmode
    \begin{enumerate}[beginpenalty=10000]
        \item \textbf{Construct the unit and counit.}
        At distributor level the unit and counit are
        \[
        \begin{aligned}
        \eta_u:\Sym(A)(-,-)&\longrightarrow u^*\circ u_*,
        &h:a\to a'&\longmapsto[\id_{u(a)},u(h)],\\
        \epsilon_u:u_*\circ u^*&\longrightarrow \Sym(C)(-,-),
        &[\alpha,\beta]&\longmapsto\beta\circ\alpha,
        \end{aligned}
        \]
        where \(\alpha:c\to u(a)\), \(\beta:u(a)\to c'\).
        Coend relations and functoriality of \(u\) give their
        well-definedness and naturality.

        \item \textbf{Verify the adjunction identities.}
        The two triangle composites are
        \[
        u_*\xrightarrow{u_*\eta_u}u_*\circ u^*\circ u_*
          \xrightarrow{\epsilon_u u_*}u_*,
        \]
        \[
        u^*\xrightarrow{\eta_u u^*}u^*\circ u_*\circ u^*
          \xrightarrow{u^*\epsilon_u}u^*.
        \]
        For \(\alpha:u(a)\to c\), the first inserts identity arrows
        and then composes to \(\alpha\). For \(\beta:c\to u(a)\),
        the second similarly returns \(\beta\). Applying the
        module-map functor proves the generated adjunction.

        \item \textbf{Construct composition comparisons.}
        The companion and conjoint comparisons are respectively
        \[
        \begin{aligned}
        [\alpha:u(a)\to c,\ \beta:v(c)\to d]
          &\longmapsto\beta\circ v(\alpha),\\
        [\alpha:d\to v(c),\ \beta:c\to u(a)]
          &\longmapsto v(\beta)\circ\alpha.
        \end{aligned}
        \]
        Their inverses insert \(c=u(a)\) and the corresponding
        identity arrow. Coend relations give both inverse equations.

        \item \textbf{Verify coherence.}
        For three functors \(u,v,w\), both companion comparison paths
        send \([\alpha,\beta,\gamma]\) to
        \(\gamma\circ w(\beta)\circ(w\circ v)(\alpha)\).
        Both conjoint paths give the corresponding reversed-endpoint
        composite. Identity functors insert identity arrows.
        These equations carry to the quantum components by functoriality. \qedhere
    \end{enumerate}
\end{proof}

A distributor square consists of functors \(u:\Sym(A)\to \Sym(C)\),
\(v:\Sym(B)\to \Sym(D)\) and a natural family
\begin{equation}
 \alpha_{a,b}:F(a,b)\longrightarrow F'(u(a),v(b)).
 \label{eq:cgqg-distributor-square}
\end{equation}
It is displayed by
\[
 \begin{tikzcd}[column sep=huge,row sep=large]
 \Sym(A)\arrow[r,Rightarrow,"F"]\arrow[d,"u"']
   &\Sym(B)\arrow[d,"v"]\\
 \Sym(C)\arrow[r,Rightarrow,"F'"']&\Sym(D)
 \arrow[phantom,from=1-1,to=2-2,"\alpha" description].
 \end{tikzcd}
\]
Its companion mate is
\begin{equation}
 \begin{split}
 \bar\alpha:v_*\circ F&\Longrightarrow F'\circ u_*,\\
 [w,\beta:v(b)\to d]&\longmapsto
       [\id_{u(a)},\beta_*\alpha(w)].
 \end{split}
 \label{eq:cgqg-distributor-square-mate-formula}
\end{equation}
The associated \textbf{representable quantum square} is the pair
consisting of \(\alpha\) and
\begin{equation}
 Q_{\bar\alpha}^{\mathrm{sp}}:
 Q_{v_*\circ F}^{\mathrm{sp}}
       \longrightarrow Q_{F'\circ u_*}^{\mathrm{sp}}.
 \label{eq:cgqg-representable-quantum-square}
\end{equation}
The companion construction turns the square into a globular cell:
\[
 \begin{tikzcd}[column sep=huge,row sep=large]
 \mathbb M_{v_*}\star\mathbb M_F
   \arrow[r,"\mathbb M_{\bar\alpha}"]
   &\mathbb M_{F'}\star\mathbb M_{u_*}.
 \end{tikzcd}
\]
Both sides have endpoints \(\mathbb Q_A,\mathbb Q_D\).
For identity vertical boundaries, the co-Yoneda unitors identify
\(Q_{\bar\alpha}^{\mathrm{sp}}\) with \(Q_\alpha^{\mathrm{sp}}\).

\begin{lemma}[Companion mates and pasting]
\label{lem:cgqg-companion-mate-calculus}
Formula~\eqref{eq:cgqg-distributor-square-mate-formula} is well defined
and natural. Identity squares and both kinds of square pasting correspond
to identity and pasting of companion mates.
\end{lemma}
\begin{proof}
    \leavevmode
    \begin{enumerate}[beginpenalty=10000]
        \item \textbf{Verify the coend relation and naturality.}
        For \(r:b\to b'\), the representatives
        \([r_*w,\beta]\), \([w,\beta\circ v(r)]\) have the same image because
        \(\alpha(r_*w)=v(r)_*\alpha(w)\).
        For \(h:a'\to a\), naturality gives
        \(\alpha(h^*w)=u(h)^*\alpha(w)\); moving \(u(h)\)
        across the target coend proves naturality at the first endpoint.
        At the last endpoint it is postcomposition of \(\beta\).

        \item \textbf{Compute vertical pasting.}
        Let \(\alpha':F'\Rightarrow F''(u'-,v'-)\) be a second square.
        The composite is
        \[
        w\longmapsto\alpha'_{u(a),v(b)}(\alpha_{a,b}(w)).
        \]
        Its mate sends \([w,\rho:(v'\circ v)(b)\to d]\) to
        \[
        [\id_{u'(u(a))},\,\rho_*\alpha'_{u(a),v(b)}(\alpha_{a,b}(w))].
        \]
        Pasting the two mates inserts two companion factors.
        The comparisons of
        Proposition~\ref{prop:cgqg-representable-vertical-composition}
        compose those factors and give exactly the displayed value.

        \item \textbf{Compute horizontal pasting.}
        For horizontally adjacent squares, use the diagram
        \[
        \begin{tikzcd}[column sep=large,row sep=large]
        \Sym(A)\arrow[r,Rightarrow,"F"]\arrow[d,"u"']
          &\Sym(B)\arrow[r,Rightarrow,"G"]\arrow[d,"v"]
          &\Sym(C)\arrow[d,"t"]\\
        \Sym(A')\arrow[r,Rightarrow,"F'"']
          &\Sym(B')\arrow[r,Rightarrow,"G'"']&\Sym(C')
        \arrow[phantom,from=1-1,to=2-2,"\alpha" description]
        \arrow[phantom,from=1-2,to=2-3,"\beta" description].
        \end{tikzcd}
        \]
        Their paste sends \([x,y]\) to \([\alpha(x),\beta(y)]\).
        Its well-definedness is the calculation
        \[
        \begin{aligned}
        [\alpha(r_*x),\beta(y)]
          &=[v(r)_*\alpha(x),\beta(y)]\\
          &=[\alpha(x),v(r)^*\beta(y)]\\
          &=[\alpha(x),\beta(r^*y)].
        \end{aligned}
        \]
        For \(\rho:t(c)\to d\), the mate of this paste sends
        \[
        [[x,y],\rho]\longmapsto
          [\id_{u(a)},[\alpha(x),\rho_*\beta(y)]].
        \]
        Pasting the two companion mates gives the same formula after
        rebracketing the coends and eliminating identity companion factors.

        \item \textbf{Identify identities and quantum components.}
        For an identity square, the mate is the co-Yoneda identity
        comparison. The preceding representative calculations therefore
        establish identities and both pasting laws for distributor mates.
        Functoriality of \(Q_{(-)}^{\mathrm{sp}}\) establishes the same
        equations for their quantum-module maps. \qedhere
    \end{enumerate}
\end{proof}

\begin{theorem}[Generated quantum-module double category]
\label{thm:cgqg-generated-qmodule-double-category}
There is a double category \(\mathbf{QMod}^{\mathrm{rep}}_{\mathrm{sp}}\)
whose objects are \(\mathbb Q_A\), horizontal arrows are \(\mathbb M_F\),
vertical arrows are functors \(u:\Sym(A)\to \Sym(C)\) equipped with their canonical
representable quantum adjunctions, and squares are the representable
quantum squares. Horizontal composition is \(\star\), horizontal identities
are \(\mathbb I_A^{\mathrm{gen}}\), and vertical composition is composition
of functors.
\end{theorem}
\begin{proof}
    \leavevmode
    \begin{enumerate}[beginpenalty=10000]
        \item \textbf{Construct the two compositions.}
        Proposition~\ref{prop:cgqg-generated-qmodule-bicategory}
        supplies the horizontal bicategory.
        Proposition~\ref{prop:cgqg-representable-vertical-composition}
        identifies the representable arrows of composite functors.
        Define square pasting by pasting the distributor squares and
        taking the quantum image of the resulting companion mate.

        \item \textbf{Identify the pasted quantum maps.}
        Lemma~\ref{lem:cgqg-companion-mate-calculus} identifies this
        definition with pasting the quantum components, including the
        horizontal associators, unitors, and companion comparisons.
        It also identifies the identity squares.

        \item \textbf{Verify the double-category equations.}
        Let \(\mathbf{Dist}_{\mathrm{TCG}}\) be the distributor double
        category labelled by games \(A\), with groupoid \(\Sym(A)\) at \(A\).
        Forgetting the quantum components defines
        \[
        U:\mathbf{QMod}^{\mathrm{rep}}_{\mathrm{sp}}
          \longrightarrow\mathbf{Dist}_{\mathrm{TCG}}.
        \]
        It is faithful on squares because each distributor square
        determines its quantum component. The identities, associativity,
        and interchange equations hold in \(\mathbf{Dist}_{\mathrm{TCG}}\).
        The construction of pasting in Steps~1--2 preserves these
        equations under \(U\), and faithfulness reflects them to the
        generated squares. \qedhere
    \end{enumerate}
\end{proof}

\subsection{Symmetric monoidal structure}
\label{subsec:cgqg-qmodule-double-tensor}

Tensor products are defined by
\[
 \mathbb Q_A\otimes\mathbb Q_B=\mathbb Q_{A\otimes B},
 \qquad
 \mathbb M_F\otimes\mathbb M_{F'}=\mathbb M_{F\extprod F'}.
\]
The unit is \(\mathbb Q_I\). Theorem~\ref{thm:cgqg-qg-tensor} and
Proposition~\ref{prop:cgqg-distributor-qmodule-tensor} identify the quantum
components with the corresponding tensor products.

\begin{theorem}[Symmetric monoidal target and realization]
\label{thm:cgqg-qmodule-symmetric-monoidal-double-category}
The generated target is a symmetric monoidal double category. The
canonical realization
\[
 \mathcal R:\mathbf{Dist}_{\mathrm{TCG}}
      \longrightarrow\mathbf{QMod}^{\mathrm{rep}}_{\mathrm{sp}}
\]
is a symmetric monoidal double functor satisfying \(U\circ\mathcal R=1\).
The canonical prescribed components also identify \(\mathcal R\circ U\)
with the identity, so \(U\) and \(\mathcal R\) give a canonical
symmetric monoidal equivalence of the generated target with
\(\mathbf{Dist}_{\mathrm{TCG}}\).
\end{theorem}
\begin{proof}
    \leavevmode
    \begin{enumerate}[beginpenalty=10000]
        \item \textbf{Tensor vertical arrows and squares.}
        Use products of functors. Hom-sets in product groupoids give
        \[
        (u\times v)_*\cong u_*\extprod v_*,
        \qquad
        (u\times v)^*\cong u^*\extprod v^*.
        \]
        Tensor distributor squares by external product of their natural
        transformations, and use their quantum companion mates.
        The formula for a companion mate is componentwise under these
        product identifications.

        \item \textbf{Compare tensor and horizontal composition.}
        The map
        \[
        \xi:\ (G\extprod G')\circ(F\extprod F')
          \longrightarrow(G\circ F)\extprod(G'\circ F')
        \]
        sends \([(x,x'),(y,y')]\) to \(([x,y],[x',y'])\).
        The two coends impose independent middle-symmetry relations
        in the two coordinates. Thus the inverse pairs representatives
        coordinatewise, and both inverse equations follow from the
        respective coend relations. Its quantum image is the tensor
        comparison for composition, by
        Proposition~\ref{prop:cgqg-distributor-qmodule-tensor}.

        \item \textbf{Verify the tensor equations.}
        Associators, unitors, and symmetries rebracket product coordinates,
        remove singleton coordinates, or permute factors. These operations
        commute with the two independent coend relations and the
        componentwise action spans. Therefore the tensor comparisons
        satisfy their coherence equations on distributor squares, and the
        faithful \(U\) reflects those equations to the quantum squares.

        \item \textbf{Construct the realization functor.}
        Define \(\mathcal R\) by adjoining the quantum groupoid and module
        components and taking companion mates on squares. The definitions
        of \(\star\), identities, and pasting make it a double functor.
        The comparisons of Steps~1--3 make it symmetric monoidal.
        Forgetting those components recovers the labelled distributor
        data, giving \(U\circ\mathcal R=1\). Conversely, every object,
        horizontal arrow, vertical arrow, and square of the generated
        target already carries precisely the components prescribed by
        \(\mathcal R\). Adjoining those components after forgetting
        them recovers the original data and all the chosen comparisons.
        Hence \(\mathcal R\circ U\) is canonically the identity, proving
        the stated equivalence. \qedhere
    \end{enumerate}
\end{proof}

\subsection{The global quantum realization}
\label{subsec:cgqg-global-quantum-realization}

For a positive morphism \(f:\sigma\to\tau\) with vertical boundaries
\(h:A\to C\), \(k:B\to D\), write
\[
 \alpha_f:\lVert\sigma\rVert\Longrightarrow
 \lVert\tau\rVert(\Sym(h)-,\Sym(k)-)
\]
for its distributor square. For composable strategies
\(\sigma:A\to B\), \(\tau:B\to C\), interaction and hiding identify a
\(+\)-covered composite configuration with a causally compatible matching
pair \((x^\sigma,x^\tau)\). Write its common middle configuration as \(b\).
The collapse compositor is
\begin{equation}
 \begin{split}
 &\gamma_{\sigma,\tau}
   (\theta_A^-,x^\tau\odot x^\sigma,\theta_C^+)\\
 &\qquad=
 [(\theta_A^-,x^\sigma,\id_b),
   (\id_b,x^\tau,\theta_C^+)].
 \end{split}
 \label{eq:cgqg-full-compositor-witness-formula}
\end{equation}
This is the full-coend construction of
\cite[Proposition~5.13]{ClairambaultOlimpieriPaquet2025}.

For exterior symmetries \(h:a'\to a\), \(k:c\to c'\), let
\([x,y]=\gamma_{\sigma,\tau}(w)\). The symmetry lifting of the composite
restricts to symmetries of its two interacting configurations, with common
middle symmetry \(r:b\to b'\). Uniqueness of endpoint transport identifies
the two component witnesses of the transported composite with
\(r_*h^*x\) and \((r^{-1})^*k_*y\). Consequently,
\begin{equation}
 \begin{split}
 \gamma_{\sigma,\tau}(k_*h^*w)
   &=[r_*h^*x,(r^{-1})^*k_*y]\\
   &=[h^*x,k_*y]
    =k_*h^*\gamma_{\sigma,\tau}(w).
 \end{split}
 \label{eq:cgqg-full-compositor-naturality}
\end{equation}
The middle equality is~\eqref{eq:cgqg-full-coend-relation}. Thus the
naturality square is
\[
 \begin{tikzcd}[column sep=large,row sep=large]
 \lVert\tau\odot\sigma\rVert(a,c)
   \arrow[r,"\gamma_{\sigma,\tau}"]\arrow[d,"k_*h^*"']
   &(\lVert\tau\rVert\circ\lVert\sigma\rVert)(a,c)
       \arrow[d,"k_*h^*"]\\
 \lVert\tau\odot\sigma\rVert(a',c')
   \arrow[r,"\gamma_{\sigma,\tau}"']
   &(\lVert\tau\rVert\circ\lVert\sigma\rVert)(a',c').
 \end{tikzcd}
\]

\begin{proposition}[A two-move interaction]
\label{prop:cgqg-two-move-oplax-example}
There are thin strategies \(\sigma:I\to B\), \(\tau:B\to I\) for which
the compositor has the form
\[
 \gamma_{\sigma,\tau}:\{*\}\longrightarrow
       \{[\varnothing,\varnothing],[b,b]\},
 \qquad *\longmapsto[\varnothing,\varnothing].
\]
\end{proposition}
\begin{proof}
    \leavevmode
    \begin{enumerate}[beginpenalty=10000]
        \item \textbf{Construct the game and strategies.}
        Let \(B\) have two concurrent events \(b^-,b^+\), of the indicated
        polarities, with discrete symmetry. Put \(b=\{b^-,b^+\}\).
        The strategy \(\sigma:I\to B\) has the same events and adds the
        causal order \(b^-<b^+\). The strategy \(\tau:B\to I\) has
        display in \(B^\perp\) and adds the causal order \(b^+<b^-\).
        Their display maps are the identity on the two named events:
        \[
        \begin{tikzcd}[column sep=large,row sep=large]
        \sigma:\quad b^-\arrow[r]&b^+
          &\qquad \tau:\quad b^+\arrow[r]&b^-.
        \end{tikzcd}
        \]
        In each strategy the new dependency runs from a negative event
        to a positive event in its displayed game. Courtesy therefore
        holds. The sole negative event is available initially and has a
        unique occurrence, giving receptivity. With discrete symmetry,
        symmetry-receptivity and thinness follow from the identity maps.
        Thus both are thin strategies.

        \item \textbf{Compute the distributor composite.}
        Each strategy has exactly two \(+\)-covered configurations:
        the empty configuration and the full configuration.
        All endpoint symmetries are identities. The coend over the
        discrete groupoid \(\Sym(B)\) therefore consists of the two matching
        pairs \([\varnothing,\varnothing]\) and \([b,b]\).

        \item \textbf{Compute interaction and hiding.}
        The full matching pair induces the causal cycle
        \[
        \begin{tikzcd}[column sep=huge]
        b^-\arrow[r,bend left=20,"\sigma"]
          &b^+\arrow[l,bend left=20,"\tau"].
        \end{tikzcd}
        \]
        The empty pair is causally compatible. Hence
        \(\Conf^+(\tau\odot\sigma)=\{\varnothing\}\), and
        \eqref{eq:cgqg-full-compositor-witness-formula} gives the displayed
        map. Its quantum image is the graph of the same proper inclusion. \qedhere
    \end{enumerate}
\end{proof}

\begin{theorem}[Symmetric monoidal oplax quantum realization]
\label{thm:cgqg-master-quantum-double-functor}
The assignments
\begin{align}
 A&\longmapsto\mathbb Q_A,\\
 h:A\to C&\longmapsto\Sym(h),\\
 \sigma:A\to B&\longmapsto
       \mathbb M_{\lVert\sigma\rVert}
       =(\lVert\sigma\rVert,Q_\sigma^{\mathrm{sp}}),\\
 f&\longmapsto(\alpha_f,Q_{\overline{\alpha_f}}^{\mathrm{sp}})
\end{align}
extend canonically to a symmetric monoidal oplax double functor
\begin{equation}
 \mathfrak Q:\mathbf{TCG}\longrightarrow
       \mathbf{QMod}^{\mathrm{rep}}_{\mathrm{sp}}.
 \label{eq:cgqg-master-double-functor}
\end{equation}
The vertical arrow \(\Sym(h)\) carries the representable adjunction
\[
 \mathbb M_{\Sym(h)_*}\dashv\mathbb M_{\Sym(h)^*}.
\]
The horizontal identity comparison is the quantum image of
\begin{equation}
 \zeta_A:\lVert\cc_A\rVert\xrightarrow{\cong}\Sym(A)(-,-).
 \label{eq:cgqg-master-horizontal-unit-comparison}
\end{equation}
For composable \(\sigma,\tau\), the compositor has quantum component
\begin{equation}
 \Gamma_{\sigma,\tau}^{\mathrm{sp}}
   =Q_{\gamma_{\sigma,\tau}}^{\mathrm{sp}}:
 Q_{\lVert\tau\odot\sigma\rVert}^{\mathrm{sp}}
       \longrightarrow
 Q_{\lVert\tau\rVert\circ\lVert\sigma\rVert}^{\mathrm{sp}}.
 \label{eq:cgqg-master-horizontal-compositor}
\end{equation}
The identity and tensor comparisons are invertible.
\end{theorem}
\begin{proof}
    \leavevmode
    \begin{enumerate}[beginpenalty=10000]
        \item \textbf{Compose the two realizations.}
        Keep the game labels in
        Theorem~\ref{thm:cgqg-background-positive-witness-double-functor}.
        The resulting factorization is
        \[
        \begin{tikzcd}[column sep=large,row sep=large]
        \mathbf{TCG}\arrow[rr,"\mathfrak Q"]
          \arrow[dr,"\lVert-\rVert"']
          &&\mathbf{QMod}^{\mathrm{rep}}_{\mathrm{sp}}\\
          &\mathbf{Dist}_{\mathrm{TCG}}
              \arrow[ur,"\mathcal R"']&.
        \end{tikzcd}
        \]
        Theorem~\ref{thm:cgqg-qmodule-symmetric-monoidal-double-category}
        constructs \(\mathcal R\). Their composite has the displayed
        object, vertical, horizontal, and square assignments.

        \item \textbf{Construct the identity and composition comparisons.}
        Proposition~\ref{prop:cgqg-copycat-quantum-module} supplies the
        identity comparison and its inverse. The graph of
        \eqref{eq:cgqg-full-compositor-witness-formula} gives
        \(Q_{\gamma_{\sigma,\tau}}^{\mathrm{sp}}\); its action compatibility
        follows from~\eqref{eq:cgqg-full-compositor-naturality} and
        Proposition~\ref{prop:cgqg-distributor-qmodule-functorial}.
        These are the quantum images of the collapse comparisons.

        \item \textbf{Verify horizontal coherence.}
        For three interacting strategy witnesses, both compositor
        paths produce their class in the threefold full coend, with
        the parenthesizations identified by \(\mathfrak a\).
        The copycat comparison composes the polarized frames, and
        the co-Yoneda unitors then apply the corresponding endpoint
        transport. These are the associativity and unit equations of
        the positive-witness collapse. Applying the module-map functor
        preserves the equations and gives horizontal coherence.

        \item \textbf{Verify square and tensor compatibility.}
        Lemma~\ref{lem:cgqg-companion-mate-calculus} identifies the
        companion mates of pasted non-globular squares with the pasted
        quantum components. The collapse is symmetric monoidal, and
        \(\mathcal R\) has the tensor comparisons constructed in
        Theorem~\ref{thm:cgqg-qmodule-symmetric-monoidal-double-category}.
        Their composite therefore satisfies the symmetric monoidal
        double-functor equations. The identity and tensor comparisons
        are images of natural isomorphisms and are invertible. \qedhere
    \end{enumerate}
\end{proof}

The realization satisfies \(U\circ\mathfrak Q=\lVert-\rVert\), with game labels
retained. It assigns to each game its canonical quantum groupoid, to each
strategy its distributor-induced quantum module, and to each positive
morphism its representable quantum square.

\section{Conclusion}
\label{sec:cgqg-conclusion}

The polarized symmetry of a thin concurrent game determines a quantum
groupoid. Unique negative--positive factorization produces the vacant
box structure; its two composition directions yield the weak Hopf
operations; and the configuration set supplies the explicit Frobenius
base. Tensor and negation transport to corresponding operations on this
quantum structure. The construction therefore gives an algebraic
interpretation of symmetry and polarity within an established model
of concurrency.

Thin strategies extend this interpretation through their associated
quantum modules. Their positive-restriction characterization identifies the
endpoint transport retained by the module construction. Together
with their generating distributors, these modules support a symmetric
monoidal oplax double-categorical realization. The necessity of the
retained distributor and the exact criterion for agreement with the
regular quantum module are established by the comparison results above. The identity and composition comparisons explain how the
realization relates copycat, symmetry transport, and causal
compatibility.

A natural next goal is a higher categorical extension of this
construction. The role of unique polarized factorization suggests
studying groupoids of factorizations and their comparison symmetries,
with contractible uniqueness taking the place of strict uniqueness.
This raises the question of how the box construction, its quantum
algebraic structure, and its extension to strategies can be developed
while retaining these comparisons and their coherence data. Such an
extension would clarify the role of thinness in the emergence of
quantum structure from concurrent games.

These developments provide concrete objects for further study at the
interface of semantics and quantum algebra. Alongside the higher
categorical generalization, two questions arise directly: which quantum
groupoids can be obtained from thin concurrent games, and which
properties of strategies can be expressed through their associated
quantum modules?

\section{Related work}
\label{sec:cgqg-related-work}

The semantic input is the theory of thin concurrent games of Castellan,
Clairambault, and Winskel~\cite{CastellanClairambaultWinskel2019}. Their
construction uses event structures and symmetry to model concurrent
computation and its interaction with higher-order behaviour.
Clairambault, Olimpieri, and
Paquet~\cite{ClairambaultOlimpieriPaquet2025} develop the polarized
factorization and positive-witness distributors used here, and establish
a symmetric monoidal oplax collapse to distributors. Our strategy
realization equips that collapse with the quantum groupoids and quantum
modules constructed from its game-semantic data.

The algebraic mechanism is related to the construction of weak Hopf
algebras from matched pairs of groupoids and cocycle data by
Andruskiewitsch and Natale~\cite{AndruskiewitschNatale2005}, formulated
through double groupoids. We identify the required vacant double
groupoid in polarized configuration symmetry and realize its box
operations in spans. Pastro and Street~\cite{PastroStreet2009} establish
the passage from weak Hopf monoids in braided monoidal categories to
quantum groupoids, developing the categorical setting introduced by Day
and Street~\cite{DayStreet2004}. We make this passage explicit for games
by constructing the configuration Frobenius base, the boundary and
composition objects, and the inversion data.

Chikhladze~\cite{Chikhladze2010QuantumModules} develops quantum modules
as quantum analogues of distributors. We use the direct action
formulation to construct modules from ordinary distributors between the
symmetry groupoids of games. Their diagonal carriers and positive
endpoint transports determine explicit action spans. The resulting
equations establish the module structure, and natural transformations
and external products give its functorial and monoidal behaviour. The positive-restriction characterization specifies exactly which
endpoint actions these modules retain, and the composition example
shows why the generating distributor is needed for the full coend.
Full distributor composition defines horizontal composition in the
generated target, which is canonically equivalent to the game-labelled
distributor double category equipped with its prescribed quantum data.

A neighbouring connection between logical interaction and quantum
topology is developed in Melli\`es's ribbon tensorial
logic~\cite{Mellies2018RibbonTensorialLogic}. Motivated by the graphical
representation of proofs, that work equips tensorial logic with
ribbon-group exchange and proves a faithful interpretation of proofs as
ribbon tangles. The topological structure guides the design of the
calculus. Here the input is the existing model of thin concurrent games,
and the construction extracts a quantum groupoid from each game's
polarized configuration symmetries. The two developments relate logical
interaction to quantum-group-theoretic structures through different
constructions: graphical representations of proofs and algebraic
structures associated with games.

\begin{thebibliography}{99}
\addcontentsline{toc}{section}{References}

\bibitem{Mellies2018RibbonTensorialLogic}
Paul-Andr\'e Melli\`es.
Ribbon tensorial logic.
In \emph{Proceedings of the 33rd Annual ACM/IEEE Symposium on
Logic in Computer Science (LICS 2018)}, ACM, 2018.
\href{https://doi.org/10.1145/3209108.3209129}
     {doi:10.1145/3209108.3209129}.

\bibitem{AndruskiewitschNatale2005}
Nicol\'as Andruskiewitsch and Sonia Natale.
Double categories and quantum groupoids.
\emph{Publicaciones Matem\'aticas del Uruguay}, 10 (2005), 11--51.
\href{https://arxiv.org/abs/math/0308228}{arXiv:math/0308228}.

\bibitem{BohmNillSzlachanyi1999}
Gabriella B\"ohm, Florian Nill, and Korn\'el Szlach\'anyi.
Weak Hopf algebras. I. Integral theory and \(C^*\)-structure.
\emph{Journal of Algebra}, 221(2) (1999), 385--438.
\href{https://arxiv.org/abs/math/9805116}{arXiv:math/9805116}.

\bibitem{CastellanClairambaultWinskel2019}
Simon Castellan, Pierre Clairambault, and Glynn Winskel.
Thin games with symmetry and concurrent Hyland--Ong games.
\emph{Logical Methods in Computer Science}, 15(1) (2019), article 18.
\href{https://doi.org/10.23638/LMCS-15(1:18)2019}{doi:10.23638/LMCS-15(1:18)2019}.

\bibitem{Chikhladze2010QuantumModules}
Dimitri Chikhladze.
Quantum modules.
Preprint, 2010.
\href{https://arxiv.org/abs/1008.1399}{arXiv:1008.1399}.

\bibitem{Chikhladze2011QuantumCategories}
Dimitri Chikhladze.
A category of quantum categories.
\emph{Theory and Applications of Categories}, 25(1) (2011), 1--37.
\url{https://tac.mta.ca/tac/volumes/25/1/25-01abs.html}.

\bibitem{ClairambaultOlimpieriPaquet2025}
Pierre Clairambault, Federico Olimpieri, and Hugo Paquet.
From thin concurrent games to generalized species of structures.
\emph{Logical Methods in Computer Science}, 21(4) (2025), article 12.
\href{https://doi.org/10.46298/lmcs-21(4:12)2025}{doi:10.46298/lmcs-21(4:12)2025}.

\bibitem{ContrerasKellerMehta2022}
Ivan Contreras, Molly Keller, and Rajan Amit Mehta.
Frobenius objects in the category of spans.
\emph{Reviews in Mathematical Physics} (2022), article 2250036.
\href{https://doi.org/10.1142/S0129055X22500362}{doi:10.1142/S0129055X22500362}.

\bibitem{DayStreet2004}
Brian Day and Ross Street.
Quantum categories, star autonomy, and quantum groupoids.
In \emph{Galois Theory, Hopf Algebras, and Semiabelian Categories},
Fields Institute Communications, vol.~43, 2004, pp.~187--226.

\bibitem{Hayashi1993}
Takahiro Hayashi.
Quantum group symmetry of partition functions of IRF models and its
application to Jones' index theory.
\emph{Communications in Mathematical Physics}, 157 (1993), 331--345.

\bibitem{Hayashi1998}
Takahiro Hayashi.
Face algebras. I. A generalization of quantum group theory.
\emph{Journal of the Mathematical Society of Japan}, 50(2) (1998), 293--315.
\href{https://doi.org/10.2969/jmsj/05020293}{doi:10.2969/jmsj/05020293}.

\bibitem{MackSchomerus1992}
Gerhard Mack and Volker Schomerus.
Quasi Hopf quantum symmetry in quantum theory.
\emph{Nuclear Physics B}, 370 (1992), 185--230.

\bibitem{McCrudden2002}
Paddy McCrudden.
Opmonoidal monads.
\emph{Theory and Applications of Categories}, 10(19) (2002), 469--485.
\url{https://tac.mta.ca/tac/volumes/10/19/10-19abs.html}.

\bibitem{NikshychVainerman2000}
Dmitri Nikshych and Leonid Vainerman.
A Galois correspondence for \(\mathrm{II}_1\) factors and quantum groupoids.
\emph{Journal of Functional Analysis}, 178(1) (2000), 113--142.
\href{https://arxiv.org/abs/math/0001020}{arXiv:math/0001020}.

\bibitem{NikshychVainerman2002}
Dmitri Nikshych and Leonid Vainerman.
Finite quantum groupoids and their applications.
In \emph{New Directions in Hopf Algebras},
Mathematical Sciences Research Institute Publications, vol.~43,
Cambridge University Press, 2002, pp.~211--262.
\href{https://arxiv.org/abs/math/0006057}{arXiv:math/0006057}.

\bibitem{PastroStreet2009}
Craig Pastro and Ross Street.
Weak Hopf monoids in braided monoidal categories.
\emph{Algebra \& Number Theory}, 3(2) (2009), 149--207.
\href{https://doi.org/10.2140/ant.2009.3.149}{doi:10.2140/ant.2009.3.149}.
\href{https://msp.org/ant/2009/3-2/ant-v3-n2-p02-p.pdf}{Published version}.
All numbered references to this work use the published version.

\end{thebibliography}
\end{document}
